\documentclass[11pt,a4paper]{article}
\usepackage[T1]{fontenc}
\usepackage[utf8]{inputenc}
\usepackage{lmodern,microtype}
\usepackage[margin=.85in]{geometry}
\usepackage{amsmath,amssymb,amsthm,mathtools,aliascnt}
\usepackage{booktabs,tabularx,array,enumitem,needspace,float}
\usepackage[section]{placeins}
\usepackage{xcolor,tikz,mdframed,graphicx}
\usetikzlibrary{arrows.meta,positioning,calc,decorations.pathreplacing}
\usepackage[hidelinks]{hyperref}
\usepackage[nameinlink,noabbrev]{cleveref}
\usepackage{fancyhdr}
\hypersetup{pdftitle={Sync HotStuff RR, Part II: Sequencer Rotation},pdfsubject={Regional rotation with Espresso ordering and explicit L1 multisignature control}}
\setlength{\headheight}{14pt}
\pagestyle{fancy}\fancyhf{}
\lhead{\small Sync HotStuff RR, Part II}\rhead{\small Sequencer Rotation}
\cfoot{\thepage}
\setlength{\parindent}{0pt}\setlength{\parskip}{4pt}
\setlength{\emergencystretch}{2em}
\renewcommand{\topfraction}{.92}
\renewcommand{\textfraction}{.06}
\renewcommand{\floatpagefraction}{.85}
\makeatletter
\setlength{\@fptop}{0pt}\setlength{\@fpbot}{0pt plus 1fil}
\makeatother
\setlist{nosep,leftmargin=*}
\newcolumntype{Y}{>{\raggedright\arraybackslash}X}
\definecolor{rrRed}{HTML}{A51C30}
\definecolor{rrBlue}{HTML}{2563EB}
\definecolor{rrTeal}{HTML}{0F766E}
\newcommand{\RR}[1]{{\color{rrRed}#1}}
\newcommand{\QC}{\mathcal C}
\newcommand{\FC}{\mathcal F}
\newcommand{\Hash}{\mathsf H}
\newcommand{\rank}{\operatorname{rank}}
\newcommand{\TB}{\mathit{TB}}
\newcommand{\tagg}[1]{\ensuremath{\mathsf{#1}}}
\newcommand{\defterm}[1]{\textcolor{rrBlue}{\textit{#1}}}
\newtheorem{definition}{Definition}[section]
\newaliascnt{assumption}{definition}
\newtheorem{assumption}[assumption]{Assumption}
\aliascntresetthe{assumption}
\crefname{assumption}{assumption}{assumptions}
\Crefname{assumption}{Assumption}{Assumptions}
\newaliascnt{proposition}{definition}
\newtheorem{proposition}[proposition]{Proposition}
\aliascntresetthe{proposition}
\crefname{proposition}{proposition}{propositions}
\Crefname{proposition}{Proposition}{Propositions}
\newcounter{protocol}
% Protocol boxes float to the top of a page, so long boxes leave no gaps.
\newenvironment{protocol}[3][htbp]{\begin{figure}[#1]\refstepcounter{protocol}\label{#2}
\begin{mdframed}[linewidth=.6pt,skipabove=0pt,skipbelow=0pt,
innertopmargin=7pt,innerbottommargin=7pt]\small
\textbf{Protocol \theprotocol. #3}\par\smallskip}{\end{mdframed}\end{figure}}
\crefname{protocol}{Protocol}{Protocols}
\Crefname{protocol}{Protocol}{Protocols}
\title{Sync HotStuff RR, Part II: Sequencer Rotation}
\author{}\date{19 October 2026}
\begin{document}
\maketitle
\begin{abstract}
Part~I runs one synchronous BFT engine inside one region. This document
lets several regions take turns producing one chain. Each regional turn
is a \emph{session} with a fixed DATA quota. When a session commits its
last DATA, the next eligible region takes over at once, without waiting
for global confirmation. A Publish Layer sends complete regional results
to Espresso. Every verifier reads the same finalized L1 controls and
Espresso history, and so accepts the same continuous prefix of DATA,
called the global prefix~$G$.

A multisignature committee can force a handoff. It does so with an
explicit L1 instruction that takes effect at a future Espresso time.
A Force closes the current production branch and restarts from $G$.
It never removes DATA that is already in $G$. Missing output alone never
causes a Force. On each node, a Session Manager prepares, enters and
retires the local programs of different sessions. The document states
which guarantees hold unconditionally and which depend on regional
service, publication service and committee behavior.
\end{abstract}

\section{Overview}
\label{sec:overview}
Part~I~\cite{part1} gives one region a synchronous BFT engine.
This document connects several regional engines into one chain.

\paragraph{A running example.}
We use one example throughout. The regions are A, B, C and D, in this
cyclic order.
\begin{itemize}
\item A's turn produces DATA 101--110. A commits all of them inside its
region and hands off to B.
\item B continues with DATA 111--115.
\item Meanwhile, verifiers have globally accepted only DATA up to 108.
\end{itemize}
Now the multisignature committee orders a Force. When the order takes
effect:
\begin{itemize}
\item everything up to 108 stays in the chain;
\item the unconfirmed DATA 109--115 of the old branch can no longer be
accepted;
\item if B is selected again, B starts a \emph{new} program that produces
a new DATA 109 after 108;
\item B's old program, which continued after 110, shuts down separately.
The two programs never share votes, locks or timers.
\end{itemize}
The rest of the document explains when a handoff happens, who chooses
the next region, why the new program continues after 108, and how old
and new programs run side by side.

\subsection{How the protocol works}
\paragraph{Two kinds of handoff.}
A \defterm{regular handoff} happens when a session finishes its quota.
A commits its last DATA, collects a regional commitment receipt, and
gives the result to B. B checks the data, its own selection and the
execution state, and then continues, possibly before A's last
DATA is globally accepted.

A \defterm{Force} is a handoff ordered by the multisignature committee.
The committee submits an L1 instruction that names a future Espresso
time. When verifiers reach the first Espresso block at or after that
time, they all apply the same instruction: they close the old branch
and continue from the current global prefix. A Force happens only
through such an instruction. Silence of a region, a monitoring alarm or
a local timeout does not cause one. The committee may use slow delivery
as a reason to act, but that judgment belongs to its own policy.

\paragraph{Three stages of a block.}
A DATA block passes through three separate stages.
\begin{enumerate}
\item \emph{Regional commitment.} The producing region commits the block
under its BFT rules and collects enough FINAL signatures.
\item \emph{Publication.} The Publish Layer sends the block and its
evidence to Espresso, which fixes its position in a finalized log.
\item \emph{Global acceptance.} Verifiers check the regional evidence,
the L1 controls and the rotation rules. If the block extends the
current global prefix, they append it.
\end{enumerate}
Only the third stage adds a block to the \defterm{global prefix} $G_u$
of observer $u$. We write ``$G=108$'' for ``the global prefix ends at
DATA 108''.

The \defterm{production branch} $R_u$ is what regions have produced so
far. It can run ahead of $G_u$: in the example, $R$ reaches 115 while
$G$ is still 108. The \defterm{export session} is the session that owns
the next DATA height after $G_u$. In the example it is still A's
session, because 109 is A's block, even though B is the region that is
currently producing. A Force closes the part of $R_u$ beyond $G_u$; it
never touches $G_u$.

A later stage, the rollup DA publication, has its own accepted prefix
$C_{\rm pub}$. It may lag behind $G_u$, but it may not skip any DATA
in $G_u$. The Publish Layer of this document, which delivers regional
results to Espresso, is a different service.

A regional receipt remains valid evidence of regional commitment even
if the branch containing the block is later closed. Global acceptance is
stronger: no Force ever removes DATA from $G$. These statements concern
input content, order and provenance. Correct execution and native state
continuity belong to the application and publication interfaces.

\begin{table}[htbp]\centering\small
\caption{Participants and responsibilities.}\label{tab:roles}
\begin{tabularx}{\linewidth}{@{}>{\raggedright\arraybackslash}p{.19\linewidth}Y@{}}\toprule
Role & Responsibility\\\midrule
L1 & Publishes immutable authorities and the ordered control log.\\
Regional replica & Runs the DATA/DRAIN engine of its sessions and
serves complete commitment evidence.\\
Multisignature service & Announces future status changes and Forces
on L1. It cannot choose a retained prefix or declare data global.\\
Publish Layer & Delivers complete regional results to Espresso. The
leader, replicas and external verifiers may act as senders.\\
Espresso & Supplies a finalized ordered log, authenticated timestamps,
complete namespace contents, and bounded finalization and visibility.\\
Observer / certifier & Verifies both sources, replays control and
exports, and derives global output and forced boundaries.\\
\bottomrule
\end{tabularx}
\end{table}

\subsection{The four view indices}
The full view is
\[
 (\beta,r,g,t).
\]
\begin{itemize}
\item $\beta$ is the \defterm{branch version}. Each effective Force
increases it by one. It is not an Espresso epoch or a time interval; it
only separates the old production branch from the new one.
\item $r$ is the \defterm{turn index}. It selects the immutable L1
authority (voting snapshot and DATA quota) used by a regional turn. It
is not a physical region identifier.
\item $g$ is the \defterm{generation}: how many times the session has
run regional recovery.
\item $t$ is the \defterm{term}: which leader turn within the
generation.
\end{itemize}
A \defterm{session} is the pair $(\beta,r)$: one region's production
turn on one branch version. Its voting snapshot and DATA quota stay
fixed through all of its recoveries. Every node and certifier derives
these numbers by replaying the same authenticated history, so no other
session identifier is needed.

\begin{table}[htbp]\centering\small
\caption{How the view changes.}\label{tab:views}
\begin{tabularx}{\linewidth}{@{}>{\raggedright\arraybackslash}p{.19\linewidth}>{\raggedright\arraybackslash}p{.25\linewidth}Y@{}}\toprule
Change & New view & Trigger and effect on timers\\\midrule
Leader rotation & $(\beta,r,g,t+1)$ & A terminal QC below $H_{\beta,r}$;
running commit timers continue.\\
Regional recovery & $(\beta,r,g+1,0)$ & A quit proof, the $\Delta_i$ wait
and finalized authority; the quit cancels generation-$g$ timers.\\
Regular handoff & $(\beta,r+1,0,0)$ & Ready on a REGULAR witness
with the terminal receipt; the old engine stops through Completion.\\
Force & $(\beta+1,r',0,0)$ & $r'=r_c$ if the export session has zero
accepted DATA, otherwise $r_c+1$; an unassigned boundary keeps its
target index. Old programs close separately.\\
\bottomrule\end{tabularx}
\end{table}

For example, $(7,20,0,0)\to(7,20,0,1)$ is a leader rotation inside A,
$(7,20,0,1)\to(7,20,1,0)$ is a recovery inside A, and
$(7,20,1,0)\to(7,21,0,0)$ is a regular handoff to B. None of these
changes $\beta$.

A Force can reuse a turn index. Suppose export session $(7,20)$ has
already added DATA to $G$. A Force then opens $(8,21,0,0)$, even if
a regular successor already used $(7,21)$. If $(7,20)$ had added no
DATA to $G$, the Force would open $(8,20,0,0)$ instead. In both cases
the new pair differs from every old pair, so votes, locks and evidence
stay separate. Where the new session starts is fixed by the DATA anchor
it keeps, not by its turn index. An L1 control that is finalized but
not yet effective does not stop production or global acceptance.

\subsection{What the protocol leaves to others}
\label{sec:scope}
The protocol fixes the following boundaries.
\begin{itemize}
\item Only a finalized, authorized MSIG instruction can cause a
cross-region Force. A Force needs no waiting period, no Ready and no
prior output.
\item Region eligibility changes only through explicit L1 STATUS
records. The protocol never disables or restores a region by itself.
\item Deciding \emph{when} to intervene is the multisignature
committee's policy. This document specifies the control interface and
the guarantees that follow from it, not a health policy.
\item The protocol does not enforce fair service among regions or
eventual inclusion of a client request. Request persistence and
resubmission belong to the execution layer.
\end{itemize}
Regional BFT timers, BLAME and Quit remain in place. Inside a region,
the only changes from Part~I are those listed in \cref{sec:protocols}.

\section{Model and assumptions}
\label{sec:model}
\subsection{L1 control state}
\label{sec:l1}
Nodes authenticate finalized L1 headers, authorities and a complete,
ordered control history. L1 timestamps only label L1 blocks. They are
neither regional commit times nor Espresso activation times.

\begin{table}[htbp]\centering\small
\caption{L1 authority and control records.}\label{tab:l1}
\begin{tabularx}{\linewidth}{@{}p{.27\linewidth}Y@{}}\toprule
Field & Meaning\\\midrule
$\tagg{AUTH}_{r,0}$ & Fixed regional snapshots, leader schedule and
DATA quota $\TB_r$ for index $r$.\\
$\tagg{AUTH}_{r,g}$, $g>0$ & Recovery leader schedule; it does not
change the initial snapshot or quota.\\
Control sequence $c$ & Unique increasing index binding the chain,
record kind and payload.\\
Activation time $\theta_a$ & Espresso timestamp threshold, announced
before its first crossing can become final.\\
STATUS & Explicit ON/OFF for named regions. It changes eligibility for
later selections, not fixed voting weights.\\
MSIG FORCE & Target branch version and activation time. It names no
retained DATA height and implies no OFF.\\\bottomrule
\end{tabularx}
\end{table}

\paragraph{Authority records.}
For each $(r,g)$, L1 supplies at most one immutable authority.
$\tagg{AUTH}_{r,0}$ fixes $\mathit{cfg}_r$ and $\TB_r$. The configuration
$\mathit{cfg}_r$ contains the members, weights, $F_i$, $\Delta_i$ and $K$
of each region. For $g>0$, only the leader schedule changes. Nodes wait
until the evidence they need arrives. Later controls and slashing never
reweight an old certificate or reset a session quota.

Authority is indexed by $r$ alone, not by branch. A new branch version
that reuses $r$ reads the same authority. Its pair $(\beta,r)$ still
names a different engine. An authority, once assigned, is never
rebound.

\paragraph{Ingesting controls.}
The L1 control history has sequence numbers and a binding history
commitment. A verifier keeps both pending and applied records. Finalization on L1
only announces a record; it becomes active at its activation time. Activation times need not
follow sequence order; records that become due together are ordered by
$(\theta_a,c)$. Malformed, unauthorized or duplicate records grant
nothing. To change an instruction, the committee submits a new record;
finalized history is never edited.

STATUS and MSIG are independent. An MSIG carries no status change. To
disable a region \emph{and} force a handoff, the committee schedules
OFF and MSIG explicitly, possibly for the same threshold. A STATUS
record takes effect even if a nearby MSIG turns out to be stale.
Neither record is inferred from the other.

\paragraph{Espresso positions and when controls apply.}
Let $J_q$ be Espresso block $q$, with authenticated timestamp $\theta_q$.
Within a block, namespace-transaction and record indices give a complete
order, also when timestamps are equal. For a threshold $d$, let
\[
 q(d)=\min\{q:\theta_q\ge d\}
\]
be the first block whose timestamp reaches $d$. For example, if blocks
500, 501, 502 have timestamps 119, 122, 125, then $q(120)=501$.

A control with threshold $d$ takes effect at the start of $J_{q(d)}$,
before any EXPORT in that block is examined. We call this start-of-block
step the \defterm{prelude}. The prelude first applies due STATUS
records and then the first matching MSIG. There is no other path to a
cross-region Force. We write ``$G$ at the start of $J_q$'' for the global
prefix just before the block's EXPORTs.

Write $\chi_E(q)$ for the eligibility snapshot after the prelude of
$J_q$ and before its EXPORTs. A region is eligible exactly when its
effective L1 flag is ON and its designated authority gives it positive
total weight. Replay derives this historical snapshot. A monitoring
report, missing output or local timeout does not change it.

OFF affects later selections only: a handoff that is already fixed
keeps its region and its snapshot weights. ON restores future
eligibility only; a closed branch version stays closed. To stop
the current branch, the committee uses MSIG. A Force without OFF leaves
the source region eligible. Eligibility is never restored automatically:
if every region is OFF, selection waits for an explicit change. No
control removes accepted DATA.

\paragraph{Genesis.}
The system starts from an initial checkpoint. It fixes the chain
identifier, the cyclic region order $\mathcal L_R$, the genesis DATA
anchor, the first replay position, the Espresso quorum trust anchor,
the control prefix and the initial branch version $\beta=0$. The initial
authority and state are either trusted or verified from their stated
sources; starting this interpreter does not make Espresso authenticate
earlier application history.

The genesis anchor carries synthetic commitment evidence and fixed,
authenticated L1 and Espresso references. A synthetic boundary with
source context $(i_b,\beta_b,r_b)=(\bot,\bot,-1)$ selects the first
eligible region in $\mathcal L_R$ at the first eligible authenticated
position, and waits if none exists. It opens session $(\beta,0)$ under
the branch version current at that moment. Normally this is $(0,0)$; if
a Force acted while genesis was still waiting, it is a later version.
At the selection height $q_s$, set $q^{\rm open}_{\beta,0}=q_s$ and
$\lambda^{\rm open}_{\beta,0}=\max\{\lambda(\text{genesis}),
\lambda(\tagg{AUTH}_{0,0})\}$, where the second term is the L1 height
at which the authority was published. A later checkpoint used as a
starting point must instead include its current boundary, its reference
bounds and the complete interpreter state.

\subsection{Data delivery and the Publish Layer}
\label{sec:delivery}
The \defterm{Publish Layer} is a delivery service run by existing
participants. It runs no consensus of its own and does not decide which
DATA is valid. A regional leader, another replica or an external
verifier can submit. The rotation rules rely only on the service
guarantee in \cref{ass:services}(i); \cref{sec:publish-implementation}
describes one implementation.

A \defterm{complete regional result} is a DATA block together with its
direct or covering FINAL receipt and the ancestry, boundaries and source
evidence needed to validate it. A proposal QC, a block hash, or a receipt
without its payload is not complete. Correct replicas offer complete
results to the Publish Layer as soon as they have them. The service may
batch results and builds EXPORT records as specified in
\cref{sec:exports}.

Regions may also send results directly to other regions. This lets a
regular successor become Ready before the terminal result is finalized
in Espresso. Without direct delivery, the successor reads the result
from Espresso after finalization, which takes longer. There is no
separate cross-region transport service.

Espresso carries finalized records, including DATA and commitment
evidence, to online readers. Readers still check the application and
rotation rules. L1 supplies its own authority and control proofs. An
execution-state snapshot may be transferred directly or rebuilt from
authenticated input history. Rebuilding state is part of Ready,
not of Espresso inclusion.

\subsection{Stake and faults}
\label{sec:stake}
Member $v$ of region $i$ has voting weight $\omega_{i,v}>0$. Let
$W_i=\sum_{v\in V_i}\omega_{i,v}$. The stake configuration fixes a bound
$F_i<W_i/2$, and Byzantine members of region $i$ hold total weight at
most $F_i$. Every regional certificate needs matching signatures from
distinct members with total weight greater than $F_i$. So every
certificate has a correct signer, and the correct members alone can form
one. This replaces the $f+1$ of $n=2f+1$ in Part~I.

Membership and weights stay fixed for a whole session, including its
recoveries. Slashing and status updates do not change them. A later
session checks old certificates against the old snapshot. The bound
$F_i$ counts every key of the snapshot that is compromised while its
certificates are still accepted. Compromise of old keys after that is
outside this model.

Correct replicas keep their protocol state and committed blocks.
Signatures cannot be forged, and hashes are binding.

\subsection{Assumptions}
\label{sec:assumptions}
\begin{assumption}[Finalized L1 records]
\label{ass:authority}
At every finalized L1 height, the authority records and the ordered
control history through that height are common and authenticated.
Finalized records never change. A verifier obtains a complete control
prefix, not an arbitrary subset of valid events.
\end{assumption}

\begin{assumption}[Authenticated synchronous communication and global clock]
\label{ass:network}
\leavevmode\par
\begin{itemize}
\item Messages use authenticated communication channels.
\item Within region $i$, the actual delay of a correct-to-correct
message is $\delta\le\Delta_i$; $\delta$ may differ across messages.
\item All nodes share a global clock.
\end{itemize}
\end{assumption}
The delay bound covers delivery and processing of complete validation
dependencies. Each replica records its own entry and voting events on
the global clock.

\begin{assumption}[Honest multisignature committee]
\label{ass:committee}
The committee acts as one honest, coordinated control service. It
submits only authorized, immutable, future-dated STATUS and MSIG FORCE
records, with activation times that satisfy the advance-announcement
condition below. It keeps no private protected prefix and cannot select
an older retained anchor. Honesty here imposes no response deadline and
no minimum interval between valid instructions.
\end{assumption}

\Needspace{8\baselineskip}
\begin{assumption}[Espresso history and time]
\label{ass:espresso}
Finalized Espresso history is consistent, persistent and publicly
verifiable, including the complete namespace content of each block, so a
verifier can also confirm that a block holds no record for this chain.
Each block carries a non-decreasing timestamp that runs at most
$\epsilon_E\ge0$ ahead of real time: if $J_q$ first becomes final at real
time $t_E(q)$, then $\theta_q\le t_E(q)+\epsilon_E$.
\end{assumption}
Long-term retrieval is a separate service, for example archival readers;
a finality certificate alone does not provide it. A returning node
catches up, or starts from an authorized checkpoint, before it confirms
new output; catching up has no fixed time bound.

\begin{assumption}[Bounded external services]
\label{ass:services}
There are fixed bounds $T_{\rm pub},T_{\rm L1},T_E>0$, valid under the
declared input-size, admission and load limits, such that:
\begin{enumerate}[label=(\roman*)]
\item \emph{Publish Layer.} If a correct participant has a complete
regional result available at real time $s$, an Espresso-admissible record
carrying that result, with its DATA and commitment evidence intact, is
submitted to Espresso by $s+T_{\rm pub}$. Submission does not rely on the
original leader alone: a working submission path and the required fee
resources exist.
\item \emph{L1.} A valid control operation actually submitted at real
time $s$ is finalized by $s+T_{\rm L1}$, and by then every correct online
participant can retrieve and verify it.
\item \emph{Espresso.} An admissible complete record submitted at real
time $s_E$ is in a finalized block by $s_E+T_E$, and by then every correct
participant that stays online and keeps up with the history can retrieve
the complete record and verify its inclusion and finality. Every finalized
block, with its complete relevant namespace, becomes retrievable and
verifiable to those participants within $T_E$ of finalization, also when
it holds no record for this chain.
\end{enumerate}
\end{assumption}
Participants are regional replicas, external observers and certifiers.
Each bound covers the whole path: inclusion, finalization, dissemination,
retrieval and validation. Each clock starts at an actual event: in (i),
when the complete result is available, not when a leader proposes the
block or one replica commits it; in (ii) and (iii), at actual submission,
which excludes any detection or decision time before it. ``Submitted''
means the complete record reaches the submission interface under its
admission rules; a local send attempt is not enough. A result already
recorded in Espresso satisfies (i). Forming the FINAL receipt is the
regional protocol's job.

These are delivery bounds only. Espresso delivers full payloads and
evidence, but does not promise that an EXPORT passes this chain's
acceptance rules: a stale base, a closed branch or a missing predecessor
can still stop it from advancing $G$, and retries needed for acceptance
are a separate task. The bounds never trigger a Force, and $T_E$ does not
replace the clock-lead bound of \cref{ass:espresso} or the
advance-announcement rule below.

For a newly available complete result, (i) and (iii) give
\[
 s_{\rm submit}\le s+T_{\rm pub},\qquad
 s_{\rm visible}\le s+T_{\rm pub}+T_E,
\]
where $s_{\rm visible}$ is the time by which the result and its Espresso
evidence are available to the readers in (iii). This bounds neither when
the regional result first forms, nor when a new region restores its
execution state.

\paragraph{Advance announcement.}
A control must be visible to everyone before it can take effect.
For a control submitted at real time $s$, the committee chooses
\[
 \theta_a>s+T_{\rm L1}+\epsilon_E+\eta,
 \qquad \eta>0.
\]
Then the control is available to correct participants before the first
block crossing $\theta_a$ can become final. This is a model condition,
not a conclusion from average L1 or Espresso block times. The L1 bound
already includes dissemination, so no second announcement delay is
added. $T_{\rm L1}$ starts at actual submission; detection, deliberation
and signature collection before submission are not covered by any
bound. Offline or unsynchronized participants catch up before they
confirm new output. Missing control history pauses verification.

\paragraph{Verification interfaces.}
We assume that regional receipt validity and the application-validity
checks below are sound under the regional model. A deployed certifier
must authenticate both source histories and bind the control version it
used. A signature or a TEE alone does not show that an old control
snapshot is complete. The publication interface additionally requires
sound execution and DA correspondence checks.

The bounded preparation assumption for a newly selected region is
stated with the Ready rules, in \cref{ass:ready-service}. The conditions
needed for progress are collected in \cref{sec:progress}.

\section{Regional engine: definitions}
\label{sec:regional}
This section defines the objects that one regional engine uses: blocks
and their positions, the session quota, certificates and receipts, and
the rules for admitting views. The engine is Part~I's engine with
sessions, a quota, DRAIN blocks and FINAL receipts added.
\Cref{tab:notation} lists the main symbols. Anchors and entry times
carry the same four indices as the view.
\begin{table}[htbp]\centering\small
\caption{Notation.}\label{tab:notation}
\begin{tabularx}{\linewidth}{@{}p{.32\linewidth}Y@{}}\toprule
Symbol & Meaning\\\midrule
$(\beta,r,g,t)$; $j$ & Branch version, turn index, generation, term; fresh position.\\
$(\beta,r)$ & Session: one region\textquotesingle s turn on one branch version.\\
$i_{\beta,r}$; $V_i$ & Session region; its regional members.\\
$\omega_{i,v},W_i,F_i$ & Member weight, total weight, and Byzantine bound $F_i<W_i/2$.\\
$\Delta_i$; $K$ & Regional synchronous delay bound; maximum DATA per term.\\
$\tagg{AUTH}_{r,g}$; $\mathit{cfg}_r$ & Immutable authority; fixed session snapshot.\\
$A_{\beta,r,g,t}$ & Anchor: the block just before the first fresh position of that view.\\
$h(B)$; $k_{\beta,r}(B)$ & DATA height; physical position in the session, counting DRAIN.\\
$\TB_r,H_{\beta,r},T_{\beta,r}$ & Session DATA quota, last DATA height, and terminal DATA.\\
$m_{\beta,r,g,t}$; $L(\beta,r,g,t)$ & Number of fresh positions in the term; designated leader.\\
$\QC_{\beta,r,g,t}(B)$; ANCHOR QC & Proposal certificate; re-certification of the opening.\\
$\FC_{\beta,r}(B)$ & Regional FINAL receipt, or its same-session coverage.\\
$Q_{\beta,r}^0,Q_v,Q^*$ & Inherited commitment evidence; entry lock; selected justification.\\
$e_v(\beta,r,g,t),\tau_v(P),a_v$ & Real-time entry, vote and quit events.\\
$\rho,\lambda(B)$ & Finalized L1 reference and its L1 height.\\
$\xi(B),\vartheta(B)$ & Signed Espresso selection reference and its timestamp.\\
$J_q,\theta_q$ & Espresso block and authenticated timestamp.\\
$\chi_E(q),\mathcal L_R$ & Eligibility snapshot after the prelude of $J_q$; cyclic region order.\\
$q_s,q_{\beta,r}^{\rm open}$ & Selection Espresso height; the session's lower bound for Espresso references.\\
$\lambda_{\beta,r}^{\rm open}$ & The session's lower bound for L1 references.\\
$G_u,R_u,C_{\rm pub}$ & Global DATA prefix; production branch; accepted DA publication prefix.\\
$T_{\rm L1},\epsilon_E$ & L1 announcement-service bound; maximum Espresso clock lead.\\
$T_{\rm pub},T_E$ & Complete-result submission bound; Espresso finalization and visibility bound.\\
$T_{\rm ready}$ & Preparation bound, measured from the public opening position.\\
$a$ & Accepted DATA count of the current export session; used for quota and for the Force target index.\\
$E$; Ready & Regional entry proof; the Ready condition of \cref{def:ready}.\\
\bottomrule\end{tabularx}
\end{table}

\subsection{Blocks, positions, and L1 references}
\label{sec:blocks}
The \defterm{anchor} $A_{\beta,r,g,t}$ is the block just before the first
fresh position of view $(\beta,r,g,t)$. In short: the anchor is the block
the view continues from. We always write all four indices.
\begin{itemize}
\item $A_{\beta,r,0,0}$ is the DATA block inherited from the preceding
session; the first session inherits genesis. It keeps the regional
commitment evidence of the session that produced it: a direct or
covering receipt (\cref{sec:coverage}). Genesis has synthetic evidence.
\item For $g>0$, NEWVIEW chooses the \defterm{recovery anchor}
$A_{\beta,r,g,0}$ (\cref{alg:recovery}). It is certified but may be
uncommitted.
\item For $t>0$, $A_{\beta,r,g,t}$ is the terminal DATA block of term
$(\beta,r,g,t-1)$, certified by that term's terminal QC.
\end{itemize}
The anchor of a view never changes. A generation starts at
$A_{\beta,r,g,0}$, and its later terms use their own anchors. ANCHOR
voting re-certifies only $A_{\beta,r,g,0}$, once per generation.

In the running example, A's session starts after anchor 100, and B's
regular session starts after anchor 110. After the Force, B's new session
starts after anchor 108. The inherited anchor keeps its original
producing session, height and hash: DATA 108 produced by $(7,20)$ stays
a $(7,20)$ block even when $(8,21)$ continues from it.

Each block has two positions.
\begin{itemize}
\item The \defterm{DATA height} $h(B)$ counts the DATA blocks on the chain
from genesis to $B$.
\item The \defterm{physical position} $k_{\beta,r}(B)$ counts DATA and DRAIN
blocks within session $(\beta,r)$, with $k_{\beta,r}(A_{\beta,r,0,0})=0$.
This local count does not change the signed fields of $A_{\beta,r,0,0}$.
\end{itemize}
A new block of session $(\beta,r)$ has the form
\[
 B=(\beta,r,k,h,\mathit{type},x,\Hash(B_{\rm parent}),\rho,\xi),
 \qquad \mathit{type}\in\{\tagg{DATA},\tagg{DRAIN}\},
\]
with $k=k_{\beta,r}(B_{\rm parent})+1$. A \defterm{DATA block} carries a
batch $x$ of requests and has $h=h(B_{\rm parent})+1$. A \defterm{DRAIN
block} carries no batch and has $h=h(B_{\rm parent})$. So DATA advances
both positions, and DRAIN advances only the physical position.

\paragraph{L1 references.}
The \defterm{reference} $\rho$ names a finalized L1 block. It holds that
block's height $\lambda(B)$ and hash, and a registry proof for the L1
control records up to that height. It tells validators which fixed
authorities and rules apply. These rules hold:
\begin{itemize}
\item along DATA ancestry, $\lambda$ never decreases;
\item the first DATA block of session $(\beta,r)$ has $\lambda\ge
\lambda^{\rm open}_{\beta,r}$, the L1 lower bound fixed by its boundary
(\cref{sec:witness}; genesis uses \cref{sec:l1});
\item a DRAIN block copies the reference of $T_{\beta,r}$.
\end{itemize}
The timestamp of the L1 block named by $\rho$ is readable from its
header, but it is not a commit time for $B$: recovery may commit $B$
much later (\cref{sec:late-receipt}). The Vote step applies no
receive-time freshness filter to $\rho$.

The L1 reference serves regional validation only. It does not choose
the time of a Force or the Espresso position used for selection.
Regional validity uses the fixed session authority, so a later status
update does not change an old receipt's voting weights.

\paragraph{Espresso selection references.}
The outer field $\xi(B)$ names a finalized Espresso header $J_q$ by
height and hash, together with its authenticated eligibility snapshot
after the prelude. Write $\vartheta(B)=\theta_q$. Its height does not
decrease along DATA ancestry, and the first DATA of session $(\beta,r)$
has $q\ge q_{\beta,r}^{\rm open}$. Terminal DATA uses the same rule as
other DATA. DRAIN copies the selection reference of $T_{\beta,r}$, just
as it copies its L1 reference. The proposal hash and signatures bind
these references and the session context. They do not change the native
mini-block encoding.

Regional full validation authenticates the referenced history, directly
or from its validated cache. Public \tagg{CheckFINAL} establishes the
intrinsic regional validity of the reference. Global acceptance
separately authenticates the required Espresso context and requires the
reference to precede the Espresso block that contains the EXPORT. So no
reference can justify itself through that same EXPORT. Historical
references make no claim about freshness at receipt time. The reference
of $T_{\beta,r}$ starts the regular successor search in \cref{sec:witness}.

\paragraph{Why the selection reference is needed.}
The terminal reference gives every node the same starting point for
choosing the next region. With it, nodes can choose before the terminal
DATA itself is published to Espresso. It is not proof that the DATA is
globally final.

For example, let the region order be A, B, C. At Espresso block $J_{500}$,
B and C are eligible, and A's opening reference is below 500. An announced
L1 status update makes B ineligible at $J_{505}$. A's terminal DATA 110
binds $\xi(110)=(500,\Hash(J_{500}))$. One node receives the regional
receipt when its Espresso view ends at 503; another receives it when its
view ends at 506. If each chose from its own latest status, they would
pick B and C respectively. With the signed reference, both start at 500
and pick B, provided no Force has closed the branch. The historical
status is derived from L1 controls at that Espresso position; Espresso
does not execute the status rule.

The later OFF does not revise this regular boundary. The reference gives
a common choice, not a fresh status. Force is the tool for closing an
already authorized branch. Bounded visibility does not remove the need
for a common reference, because honest nodes can receive a receipt at
different times.

The next regular selection starts from the greater of two heights: the
signed reference height of $T_{\beta,r}$, and the completed session's
opening height plus one. This moves successive selection floors forward,
without making the current session wait for a later header before it can
enter or produce its terminal result. If a later header is already known,
the direct handoff need not wait for terminal publication. Otherwise the
current complete result can be published first; its own inclusion then
supplies a later header. DRAIN keeps the terminal DATA's reference.

\subsection{Session quota, term positions, and DRAIN}
\label{sec:quota}
\paragraph{1. Fix the session endpoint.}
Session $(\beta,r)$ may add $\TB_r\ge1$ DATA blocks after
$A_{\beta,r,0,0}$. Its \defterm{last DATA height} is
\[
 H_{\beta,r}=h(A_{\beta,r,0,0})+\TB_r.
\]
The DATA block at this height is the \defterm{terminal DATA}
$T_{\beta,r}$. Recovery never changes $H_{\beta,r}$. The quota counts the
DATA blocks after $A_{\beta,r,0,0}$ in a proposal's ancestry, committed
or not. DRAIN, resends and ANCHOR votes use no quota.

\paragraph{2. Produce DATA while space remains.}
If $h(A_{\beta,r,g,0})<H_{\beta,r}$, generation $(\beta,r,g)$ is a
\defterm{DATA generation}. Each term begins at its anchor
$A_{\beta,r,g,t}$ and has
\[
 m_{\beta,r,g,t}=\min\{K,\;H_{\beta,r}-h(A_{\beta,r,g,t})\}
\]
new positions, numbered $j=0,\ldots,m_{\beta,r,g,t}-1$: at most $K$ per
term, and never more than the remaining quota. At position $j$, the DATA
block has
\[
 h=h(A_{\beta,r,g,t})+j+1,\qquad
 k=k_{\beta,r}(A_{\beta,r,g,t})+j+1.
\]
A new position is a \defterm{fresh position}, and its block is a
\defterm{fresh block}. The anchor itself is not fresh. The last position
$j=m_{\beta,r,g,t}-1$ is the \defterm{terminal position}, and its QC is
the \defterm{terminal QC}.

If the terminal DATA height of a term is below $H_{\beta,r}$, its QC
admits term $(\beta,r,g,t+1)$, with that terminal block as
$A_{\beta,r,g,t+1}$. If it equals $H_{\beta,r}$, there is no next term.
Only the last DATA term can have fewer than $K$ positions. The initial
generation is always a DATA generation, since $\TB_r\ge1$.

\paragraph{3. Use DRAIN when recovery starts at the endpoint.}
If $g>0$ and $h(A_{\beta,r,g,0})=H_{\beta,r}$, generation $(\beta,r,g)$ is
a \defterm{drain generation}. It has only term $t=0$ and position $j=0$,
with $m_{\beta,r,g,0}=1$. Its one fresh block is a DRAIN, whose parent is
$A_{\beta,r,g,0}$ and whose justification is the generation's ANCHOR QC.
It carries no requests, keeps DATA height $H_{\beta,r}$, and has no next
term.

DRAIN exists to supply a commit timer. ANCHOR voting starts no timer, so
re-certifying $A_{\beta,r,g,0}$ alone does not commit it. At height
$H_{\beta,r}$, a DATA child would exceed the quota. A DRAIN child is
voted on like any proposal, and its timer commits its ancestors,
including $T_{\beta,r}$.
\begin{itemize}
\item A DATA generation that reaches $H_{\beta,r}$ does not append DRAIN.
It waits for $T_{\beta,r}$'s timer or recovers after a quit.
\item A later drain generation may start from an earlier DRAIN.
Each interrupted recovery may add one DRAIN, so their number is not
bounded.
\item The choice depends on the recovery anchor's DATA height, not on
whether a replica has already committed $T_{\beta,r}$.
\end{itemize}
DRAIN stays in the regional log. The external chain contains only DATA,
and the next session inherits $T_{\beta,r}$.

\paragraph{Example.}
Let $h(A_{\beta,r,0,0})=100$, $\TB_r=10$, and $K=4$, so $H_{\beta,r}=110$.
\begin{center}\small
\begin{tabular}{@{}lll@{}}\toprule
View & Anchor height & Fresh DATA heights\\\midrule
$(\beta,r,0,0)$ & 100 & 101--104\\
$(\beta,r,0,1)$ & 104 & 105--108\\
$(\beta,r,0,2)$ & 108 & 109--110\\\bottomrule
\end{tabular}
\end{center}
Recovery from DATA 107 starts a new generation whose first term produces
108--110. Suppose instead that 110 has a QC but a Quit cancels its timer,
and recovery selects 110 again. DATA 111 would exceed the quota, and
ANCHOR voting on 110 starts no timer. So the recovery leader proposes a
DRAIN on top of 110; its timer commits 110. In both cases the endpoint
stays 110.

\subsection{Certificates and signed messages}
\label{sec:certificates}
\Cref{tab:certificates} lists the certificates of a session. Each needs
matching signatures from distinct members of $V_{i_{\beta,r}}$ with total
weight $>F_{i_{\beta,r}}$.

\begin{table}[htbp]\centering\small
\caption{Certificates of a session.}
\label{tab:certificates}
\begin{tabularx}{\linewidth}{@{}>{\raggedright\arraybackslash}p{.25\linewidth}>{\raggedright\arraybackslash}p{.24\linewidth}Y@{}}
\toprule
Certificate & Signed message & Purpose\\
\midrule
\defterm{Proposal QC} $\QC_{\beta,r,g,t}(B)$ & DATA-VOTE or DRAIN-VOTE & Justifies a
child. A terminal QC below $H_{\beta,r}$ also admits the next term.\\
\defterm{ANCHOR QC} & ANCHOR-VOTE & Re-certifies the opening $A_{\beta,r,g,0}$ in
generation $g$, before any fresh proposal.\\
\defterm{BLAME proof} & BLAME & Freezes a generation (\cref{sec:quit}).\\
\defterm{Regional receipt} $\FC_{\beta,r}(B)$ & FINAL & Proves that a correct replica
committed DATA block $B$.\\
\bottomrule
\end{tabularx}
\end{table}

A correct replica casts at most one fresh vote per $(\beta,r,g,t,j)$ and
one ANCHOR vote per $(\beta,r,g)$. DATA and DRAIN votes follow the same
rule but are different signed types. No QC and no ANCHOR vote starts a
commit timer. A QC means only that enough votes exist; it does not mean
that any timer has expired.

\paragraph{Signed context.}
Every signature binds the chain identifier, protocol domain, session
$(\beta,r)$, voting configuration, assigned region and message type.
View messages also bind $(g,t)$ and the L1 authority reference
$\tagg{AUTH}_{r,g}$; their full view is $(\beta,r,g,t)$. The other
signed fields are:
\begin{itemize}
\item \textbf{PROPOSE and fresh votes:} the generation's opening,
position $j$, block type, physical position, DATA height and block hash;
\item \textbf{NEWVIEW and ANCHOR-VOTE:} the chosen opening. NEWVIEW
also binds its justification and entry proof;
\item \textbf{STATUS:} the sender's entry lock and entry proof;
\item \textbf{BLAME:} the view only, with no opening, block hash or
progress count. All complaints for a view sign the same statement;
\item \textbf{FINAL:} the session and DATA fields below, with no $g$ or $t$.
\end{itemize}
Displayed messages include $\beta$ and $r$ but omit the common chain,
protocol, region and configuration fields where they are clear from
context. A verifier also checks the boundary that assigned the pair
$(\beta,r)$; the pair alone is not evidence of authority.

\paragraph{Rank and the base justification.}
Within session $(\beta,r)$, a QC formed in term $(\beta,r,g,t)$ for block
$B$ has \defterm{rank} $(g,t,k_{\beta,r}(B))$, compared lexicographically.
The ANCHOR QC of generation $g$ has rank $(g,0,k_{\beta,r}(A_{\beta,r,g,0}))$,
just below the generation's first fresh QC. A fixed total order on
certified content breaks ties; signer sets and encodings do not matter.
The \defterm{opening identity} of generation $g$ is
$(g,k_{\beta,r}(A_{\beta,r,g,0}),h(A_{\beta,r,g,0}),\Hash(A_{\beta,r,g,0}))$.
As in Part~I, ``opening'' always means the generation's opening block
$A_{\beta,r,g,0}$. QCs are ranked only within one session $(\beta,r)$.

The anchor $A_{\beta,r,0,0}$ also has a \defterm{base justification}
$Q_{\beta,r}^0$: the anchor together with its direct or covering
commitment evidence in the boundary witness
(\cref{sec:coverage,sec:witness}). Its justified block is exactly
$A_{\beta,r,0,0}$, even if the evidence uses a receipt for a higher
block of the old session; that higher block is evidence only.
$Q_{\beta,r}^0$ ranks below every QC of session $(\beta,r)$. Recovery
uses it when no QC of the session is known. A QC of another session is
never a justification in session $(\beta,r)$, so it cannot replace
$A_{\beta,r,0,0}$.

\paragraph{Regional receipts.}
After a replica commits a DATA block $B$ of session $(\beta,r)$, by $B$'s
own timer or as an ancestor, it sends
\[
 \langle\tagg{FINAL},\beta,r,\mathit{cfg}_r,h(B),\Hash(B)\rangle_v
\]
to its region. Matching FINAL messages of weight $>F_{i_{\beta,r}}$ form
the \defterm{regional receipt} $\FC_{\beta,r}(B)$. The receipt carries
$B$ and the ancestry needed to validate it, or refers to a validated
local cache.

On first obtaining a complete result, each correct holder offers it to
the Publish Layer and multicasts the full validated receipt, the DATA
and their dependencies to the fixed regional membership. Duplicates are
removed by session and DATA hash; the outgoing action is kept until it
is dispatched. This rule also covers a receipt used for Completion and
a holder that is Completed or archival. The holder may also send the
result directly to other regions for an early handoff.

\begin{definition}[\defterm{Valid regional receipt}]
\label{def:receipt}
A regional receipt $\FC_{\beta,r}(B)$ is valid when all of the following hold:
\begin{enumerate}[label=(\roman*)]
\item $B$ is a DATA block of session $(\beta,r)$, extends
$A_{\beta,r,0,0}$, and has $h(A_{\beta,r,0,0})<h(B)\le H_{\beta,r}$;
\item distinct members in session $(\beta,r)$'s fixed voting snapshot have
valid signatures on the same
$\langle\tagg{FINAL},\beta,r,\mathit{cfg}_r,h(B),\Hash(B)\rangle$,
with total weight $>F_{i_{\beta,r}}$ and the signed context specified above; and
\item $B$ and its required ancestry pass the block-validation rules of
\cref{sec:admission}.
\end{enumerate}
Regional engines check the full evidence under \cref{sec:admission};
they do not replay other signers' timers. The public interpreter uses
\cref{sec:receipt-kernel} to establish the same intrinsic validity from
the exact DATA and FINAL signatures. The weight threshold guarantees a
correct signer, and that signer's own validation and commit ground the
conclusion.
\end{definition}
Here $\mathit{cfg}_r$ identifies session $(\beta,r)$'s \defterm{fixed
voting snapshot}. $\FC(B)$ without a session subscript means the receipt
from the session that produced $B$. An inherited anchor keeps that
receipt; it does not acquire a new one. The public checker of
\cref{sec:receipt-kernel} establishes condition~(iii) without requiring
any optional admission sidecar from a publisher. The regional engines
still check the full evidence.
\begin{itemize}
\item FINAL binds no generation or term, so FINAL messages for one block
combine across recoveries. The block hash keeps the block's signed
context.
\item If a DRAIN commits $T_{\beta,r}$, replicas send FINAL for
$T_{\beta,r}$, never for the DRAIN.
\item Replicas of session $(\beta,r)$ send no FINAL for $A_{\beta,r,0,0}$.
It keeps its original commitment evidence, which may come from an earlier
session if intermediate sessions added no DATA.
\item A replica may also commit $B$ and its uncommitted ancestors when it
obtains a valid $\FC_{\beta,r}(B)$ of its own session. Completion does
this for $T_{\beta,r}$ (\cref{alg:steady}). Such a commit sends no FINAL,
because the receipt already exists. Only Completion also cancels timers.
Receipts of other sessions serve only as history or anchor evidence;
they never change this session's locks or log.
\end{itemize}

\paragraph{Receipt coverage.}
\label{sec:coverage}
A valid $\FC_{\beta,r}(D)$ also proves commitment of every DATA ancestor
$B$ of $D$ in the same session: at least one correct signer committed
$D$, and committing $D$ commits its ancestors. The \defterm{covering
evidence} for $B$ is the valid receipt for $D$ plus the fully verified
ancestry from $D$ to $B$. It does not claim that the signatures name $B$
directly. For example, a receipt for 110 with the exact path back to 109
proves that 109 was committed.

Coverage never crosses sessions: a receipt of B's session for 111 does
not replace A's evidence for 110. Coverage also does not show that $B$
belongs to the authorized external branch. A later Force may discard $D$
while keeping $B$; the regional evidence for $B$ still holds. External
observers, retained-anchor witnesses and base justifications may use
direct or covering evidence, with no extra FINAL messages or certificate
round. We write ``\defterm{commitment evidence} for $B$'' for either form.

\subsection{Admission, validation, and engine states}
\label{sec:admission}
\paragraph{\defterm{Admission.}}
Replica $v$ admits views as follows.
\begin{itemize}
\item $(\beta,r,0,0)$ is admitted at $v$'s Ready event for the boundary
that opens session $(\beta,r)$ (\cref{def:ready}). The boundary witness is
the \defterm{entry proof}. The NEWVIEW must choose exactly
$A_{\beta,r,0,0}$.
\item $(\beta,r,g,0)$ with $g>0$ is admitted once $v$ holds
$\tagg{AUTH}_{r,g}$ and a quit proof for generation $g-1$. That proof is
the \defterm{entry proof}. The NEWVIEW is justified by a QC of an earlier
generation of session $(\beta,r)$, or by $Q_{\beta,r}^0$.
\item $(\beta,r,g,t)$ with $t>0$ is admitted once $v$ also holds the
terminal QCs of terms $0,\ldots,t-1$, each for DATA below $H_{\beta,r}$.
\end{itemize}
Term zero is admitted before its NEWVIEW arrives, so its progress checks
run even if the initial leader is silent. Fresh proposals need the
generation's ANCHOR QC. This holds in every generation, including
generation $0$ of the first session from genesis; Part~I's genesis
shortcut does not apply. A message is \defterm{admissible} if its view is
admitted and its signatures are valid.

A verifier that checks a remote session's full regional history has no
local admission event for it. There, a view counts as admitted if the
evidence contains its \defterm{admission witness}: the boundary witness
for $(\beta,r,0,0)$, the chain of entry proofs for $(\beta,r,g,0)$, and
the earlier terminal QCs for $t>0$.

\paragraph{\defterm{Complete regional validation.}}
Complete regional validation checks signatures, authority, admission,
block type, quota, both positions, the L1 and Espresso reference rules,
the parent rule (\cref{tab:parent}), parent QCs, hash links and
application validity throughout the required ancestry. Every relevant
ancestor comes with its signed proposal and block data; a chain of
certificate hashes alone is not enough. Dependencies come from the
message or from a validated local cache. Correct senders send
\defterm{complete validation bundles}. A reference to a dependency is
allowed only if the receiver can resolve it within the delay bound of
\cref{ass:network}. Bundles list each dependency once.

Regional Vote, recovery, and processing of own-session evidence and
receipts always use this complete procedure. Ready also checks the
target's admission, exact state and boundary inputs in full; only the
receipt of the inherited anchor, which another session produced, may be
checked with the public kernel of \cref{sec:receipt-kernel}.

\paragraph{Evidence takes priority.}
Before validating a message or discarding a duplicate, inspect all signed
material in it for a quit proof (\cref{sec:quit}). A proof needs its own
valid signatures and admission witnesses. A valid proof counts even if
the message that carries it is invalid or incomplete. Compare incoming
material with all stored material, not only with objects inside the
message.

\paragraph{Processing proposals.}
Unless retirement has closed the signing gate, every admitted fresh
proposal of the current generation that $v$ has not voted for goes
through the Vote step. Closing programs still inspect and forward signed
evidence. This includes proposals inside attached histories, proposals
of earlier terms, and proposals whose blocks are already committed as
ancestors. A proposal's identity is $(\beta,r,g,t,j,\Hash(B))$; the signer
set or encoding of its parent QC does not change it. Duplicates are
removed by this identity after evidence inspection. Old-generation
proposals are stored without new votes or timers. Future-generation
messages wait for entry.

\paragraph{Engine states.}
A replica runs one engine for each session pair $(\beta,r)$ assigned to
its region. Engines of different sessions stay separate, even if they
share a turn index. The engine of session $(\beta,r)$ has three core
phases; the full program lifecycle in \cref{sec:session-lifecycle} adds
Waiting, Ready and Terminated.
\begin{itemize}
\item \textbf{\defterm{Active}:} the current generation is admitted and
unfrozen. The engine proposes, votes and changes terms.
\item \textbf{\defterm{Frozen}:} the current generation has a quit proof.
The engine stops proposals, votes and term changes. It still processes
evidence, old votes, QCs, receipts and recovery.
\item \textbf{\defterm{Completed}:} the replica holds
$\FC_{\beta,r}(T_{\beta,r})$ satisfying \cref{def:receipt}. Production in
session $(\beta,r)$ is over, even if the external handoff is still
pending. The engine only serves data.
\end{itemize}
Committing $T_{\beta,r}$ locally does not complete the engine: other
members may still need this replica's recovery votes before a receipt
forms. Each message goes to the engine named by its signed pair
$(\beta,r)$, so evidence from one session or branch never freezes
another.

These states describe the local engine, not external finality. The
external chain may move to another branch while an old engine still
runs. L1 status updates do not rewrite an engine's state or log. An
effective Force closes the external branch version; the Session Manager
then requests the retirement procedure of \cref{sec:retirement}. A
closing Active or Frozen program keeps that phase until a genuine
Completion or FinishClose.

\begin{table}[htbp]\centering\small
\caption{Parent selection for the next fresh position $j$ of term $(\beta,r,g,t)$.}
\label{tab:parent}
\begin{tabularx}{\linewidth}{@{}l Y Y@{}}
\toprule
Condition & $B_{\mathrm{parent}}$ & $Q_{\mathrm{parent}}$\\
\midrule
$j=0,t=0$ (DATA or DRAIN) & Opening $A_{\beta,r,g,0}$ & Its ANCHOR QC in $(\beta,r,g,0)$\\
$j=0,t>0$ & Anchor $A_{\beta,r,g,t}$, the terminal block of $(\beta,r,g,t-1)$ & Its proposal QC in $(\beta,r,g,t-1)$\\
$1\le j<m_{\beta,r,g,t}$ & Position $j-1$ of $(\beta,r,g,t)$ & Its proposal QC in $(\beta,r,g,t)$\\
\bottomrule
\end{tabularx}
\end{table}

\section{Regional protocols}
\label{sec:protocols}
\Cref{alg:steady,alg:recovery} run inside the engine of one session
$(\beta,r)$. They are Protocols 1 and 2 of Part~I~\cite{part1}, with the
changes listed below. \RR{Red text} marks where they differ from Part~I.
The selection reference $\xi$ is defined in \cref{sec:blocks}.

Two rules connect these protocols to the cross-region machinery. First,
external replay never changes the expiry of a running regional timer;
only a genuine Quit or Completion cancels one. Second, retirement can
block fresh signing, new timers and further entry, through the gate
below.

\paragraph{Notation.}
The full view is $(\beta,r,g,t)$ throughout. Every anchor, receipt and
entry time below names its branch and session explicitly. In these two
protocols, $\Delta$ and $F$ mean $\Delta_{i_{\beta,r}}$ and
$F_{i_{\beta,r}}$. The authority fields $\mathit{cfg}_r$, $\TB_r$ and
$\tagg{AUTH}_{r,g}$ are indexed by $r$ alone: a later branch that reuses
$r$ reads the same immutable authority. Its anchor, endpoint and engine
state can still differ.

\paragraph{Changes from Part~I.}
\begin{enumerate}[label=\arabic*.]
\item \emph{Stake.} Certificates need weight $>F$, not $f+1$ signatures.
\item \emph{Blocks.} A block also carries its session, physical
position, type, L1 reference $\rho$, and Espresso reference $\xi$
(\cref{sec:blocks}). Validity checks all of them.
\item \emph{Quota.} A term has $m_{\beta,r,g,t}\le K$ positions. No term
follows $H_{\beta,r}$.
\item \emph{DRAIN.} A drain generation proposes one DRAIN. Its vote
counts as progress.
\item \emph{Receipts.} Replicas send FINAL after committing DATA. A
receipt for $T_{\beta,r}$ completes the engine (step~5 of \cref{alg:steady}).
\item \emph{Timers.} Completion also cancels timers. A Force never does.
\item \emph{Recovery entry.} After the $\Delta$ wait, a replica enters
when $\tagg{AUTH}_{r,g+1}$ is finalized. $Q_{\beta,r}^0$ can serve as the
entry lock and as the justification.
\item \emph{Genesis.} The initial session $(\beta,0)$ also certifies its
anchor with an ANCHOR QC before its first fresh proposal.
\item \emph{Sessions.} Each message goes to the engine of its session.
The initial session $(\beta,0)$ opens by bootstrap (\cref{alg:bootstrap}).
\end{enumerate}

\paragraph{Retirement gate and core events.}
The following actions require that the Session Manager holds no
authenticated $\tagg{Closed}[\beta]$, that the program's
\texttt{close\_requested} is false, and that the program is not
Terminated: any fresh proposal; any DATA, DRAIN or ANCHOR vote; STATUS
and NEWVIEW; BLAME from a progress check; term entry; and recovery entry.
The same check applies to buffered-message and authority-arrival
callbacks.

The gate never prevents checking, comparing, storing and relaying signed
evidence; processing a genuine Quit; executing a surviving commit timer
and sending its FINAL; or applying a valid receipt and Completion. It
does not overwrite a lock or cancel a timer. \Cref{sec:retirement}
defines when the gate closes and how existing work drains. Entry, Quit,
recovery, Completion and timer-drain changes emit a \texttt{CoreEvent} to
the local Session Manager; that event does not change public $G$.

\paragraph{Message-processing order.}\leavevmode
\RR{Replica $v$ first passes each message to the engine of the session
pair $(\beta,r)$ that the message binds. L1 updates go to \cref{alg:force}. Boundary witnesses and receipts of
other sessions go to \cref{alg:rotation,alg:external}.} The engine handles each received
message in this order:
\begin{enumerate}[label=(\arabic*)]
\item check all signed material for a quit proof;
\item check completeness and validity;
\item process attached DATA \RR{and DRAIN} proposals under Vote, parents
first;
\item process attached terminal QCs under step~4 of \cref{alg:steady},
in increasing term order;
\item process the message itself.
\end{enumerate}
When a message and a timer fall at the same time, $v$ handles the
message first.

\subsection{Steady state}
\label{sec:steady}
\Cref{alg:steady} gives the steady state of a session: proposals, votes,
timers, term rotation, and completion.

\begin{protocol}{alg:steady}{Propose, vote, commit, normal rotation\RR{, and completion}}
Replica $v$ runs these steps \RR{in the engine of session $(\beta,r)$}. Its
current term is $(\beta,r,g,t)$.
\begin{enumerate}[label=\textbf{\arabic*)},itemsep=3pt]
\item \textbf{Propose.} Run this step only if $v=L(\beta,r,g,t)$, $g$ is
unfrozen\RR{, and session $(\beta,r)$ is not complete}.
\RR{In a DATA generation,} for each position
$j=0,\ldots,\RR{m_{\beta,r,g,t}}-1$, in order:
 \begin{itemize}
 \item Wait until $v$ holds $B_{\mathrm{parent}}$ and
       $Q_{\mathrm{parent}}$ for $j$ (\cref{tab:parent}).
 \item Build \RR{$B=(\beta,r,k,h,\tagg{DATA},x,\Hash(B_{\mathrm{parent}}),\rho,\xi)$,
       with $k$ and $h$ from \cref{sec:quota} and the references
$\rho,\xi$ from \cref{sec:blocks}}.
       The batch $x$ follows \cref{sec:requests}.
 \item Broadcast
       $P=\langle\tagg{PROPOSE},\beta,r,g,t,j,B,Q_{\mathrm{parent}}\rangle_v$
       with its validation bundle \RR{to regional replicas. }
 \item Process $P$ under Vote.
 \end{itemize}
 \RR{In a drain generation, propose once, at $(t,j)=(0,0)$. The block is
 a DRAIN with parent $A_{\beta,r,g,0}$, no batch, DATA height $H_{\beta,r}$, and the L1
 and Espresso references of $T_{\beta,r}$. Broadcast and process it as above.}
 Do not wait for the parent to commit. Never sign two different
 blocks for one position. A resend repeats the same $P$.

\item \textbf{Vote.} Run this step on every DATA \RR{or DRAIN} proposal
$P$ of term $(\beta,r,g_P,t_P)$ at position $j_P$. Here $t_P$ may be earlier
than $t$. The step also applies to proposals inside attached histories.
 \begin{enumerate}[label=(\alph*)]
 \item If $P$ or its attached material contains a quit proof for $g$,
       run Quit and stop.
 \item If $g_P<g$, keep $P$ and stop. If $g_P>g$, buffer $P$
       until $v$ enters generation $g_P$.
 \item If $g$ is frozen \RR{or session $(\beta,r)$ is complete}, keep $P$ and
       stop.
 \item If a dependency is missing, keep $P$, request the dependency,
       and stop. Rerun this step when it arrives.
 \item If $P$ is invalid, discard it and stop. \RR{Validity includes the
       block type, the quota, both positions, and the L1 and Espresso references.}
 \item If $v$ has already voted for $P$, stop.
       A duplicate never restarts a timer.
 \item Otherwise, let $\tau_v(P)$ be the current time. Forward $P$ and
       its bundle to all \RR{regional replicas. }
       Broadcast
       $\langle\tagg{DATA\mbox{-}VOTE},\beta,r,g,t_P,j_P,\Hash(B)\rangle_v$
       \RR{(DRAIN-VOTE if $B$ is a DRAIN)}.
       Start the commit timer of $P$. It expires at $\tau_v(P)+2\Delta$.
 \end{enumerate}

\item \textbf{Commit.} When the timer of $P$ expires, commit $B$ and
its uncommitted ancestors. \RR{For each DATA block $D$ of session $(\beta,r)$
that this commits, send
$\langle\tagg{FINAL},\beta,r,\mathit{cfg}_r,h(D),\Hash(D)\rangle_v$ to the regional fixed membership and the Publish Layer collection interface. A committed DRAIN yields FINAL for $T_{\beta,r}$, never for
itself. } Only Quit \RR{and Completion} cancel a timer. Entering a
new term does not cancel it. A QC never starts a timer. \RR{A force
instruction never cancels a timer.}

\item \textbf{Terminal QC and rotation.} Run this step the first time
$v$ obtains a valid terminal QC for a term $(\beta,r,g,t')$.
``Obtains'' means that $v$ assembles the QC from votes or finds it in
any message, including as another proposal's $Q_{\mathrm{parent}}$.
 \begin{enumerate}[label=(\alph*)]
 \item Process each unseen proposal in the QC's history under Vote,
       parents first.
 \item If $g$ is now frozen, stop. \RR{If the QC certifies DATA at
       $H_{\beta,r}$, or a DRAIN, stop: no next term exists.}
 \item Send the QC and its validation bundle to $L(\beta,r,g,t'+1)$.
 \item If $t\le t'$, enter $(\beta,r,g,t'+1)$ and set $e_v(\beta,r,g,t'+1)$ to the
       current time. Cancel the progress checks of all earlier
       terms.
 \end{enumerate}
 A QC never moves $v$ to an earlier term. It never changes a recorded
 entry time.

\item \RR{\textbf{Completion.} Run this step the first time $v$ obtains a
regional receipt $\FC_{\beta,r}(T_{\beta,r})$ satisfying \cref{def:receipt}.}
 \begin{itemize}\color{rrRed}
 \item Multicast the complete validated terminal receipt and its dependencies
       to the fixed regional membership, once, and persist this outgoing action.
 \item Commit $T_{\beta,r}$ and its uncommitted ancestors.
 \item Mark session $(\beta,r)$ Completed. Stop its proposals, votes, progress
       checks, term changes, and recovery. Cancel its unexpired timers.
 \item Offer the complete terminal result to the Publish Layer.
       It may also be sent directly to the next region.
 \item Keep serving data. \Cref{alg:rotation} runs the handoff.
 \end{itemize}
 \RR{A force instruction never runs this step.}
\end{enumerate}
\end{protocol}

\Needspace{9\baselineskip}
The \defterm{acceptance time} $\tau_v(P)$ is when complete validation and the first
vote occur, not an earlier arrival of an incomplete packet. Each replica
records its own entry time for each term. In session $(\beta,r)$, $e_v(\beta,r,0,0)$ is
set at the atomic Ready/Enter event, in step~1 of \cref{alg:bootstrap}. For $g>0$, $e_v(\beta,r,g,0)$ is
set in step~3 of \cref{alg:recovery}. For $t>0$, $e_v(\beta,r,g,t)$ is set in
step~4 of \cref{alg:steady}. A term's progress checks therefore start
when the term is admitted locally, not when its leader's first proposal
arrives. A silent new leader cannot delay them.

A proposal or QC at height $H_{\beta,r}$ does not end the session. Without a
quit, the timer of $T_{\beta,r}$ commits it; no later block is needed. Only a
receipt for $T_{\beta,r}$ completes the engine.

\subsection{Quit proofs and recovery}
\label{sec:quit}
A quit proof freezes a generation. \Cref{alg:recovery} then moves the
session to the next generation.

\begin{protocol}{alg:recovery}{Blame, quit, and recovery}
Replica $v$ runs these steps \RR{in the engine of session $(\beta,r)$}. Its
current term is $(\beta,r,g,t)$.
\begin{enumerate}[label=\textbf{\arabic*)},itemsep=3pt]
\item \textbf{Blame.} \RR{Let $q=m_{\beta,r,g,t}$ in a DATA generation and
$q=1$ in a drain generation.} For each $p=1,\ldots,\RR{q}$, run a
progress check at time
\[
 e_v(\beta,r,g,t)+(2p+4)\Delta.
\]
\RR{Schedule the check for $p=1$ at entry, even if the opening is not yet
known. Schedule the others once $q$ is known. A check whose time has
already passed runs at once.}
At that time, broadcast $\langle\tagg{BLAME},\beta,r,g,t\rangle_v$ if all of
these hold:
 \begin{itemize}
 \item $v$'s current term is still $(\beta,r,g,t)$, $g$ is unfrozen\RR{, and
       session $(\beta,r)$ is not complete};
 \item $v$ has cast DATA \RR{or DRAIN} votes for fewer than $p$ positions
       of term $(\beta,r,g,t)$ (ANCHOR votes do not count);
 \item $v$ has not yet sent BLAME for $(\beta,r,g,t)$.
 \end{itemize}

\item \textbf{Quit.} Run this step the first time $v$ holds a quit proof
for its current generation $g$.
 \begin{itemize}
 \item Freeze $g$. Stop all progress checks, proposals, votes, and
       term changes in $g$. Cancel every unexpired commit timer of $g$.
 \item Set $a_v$ to the current time. Broadcast the proof.
 \item Keep all commitments. Keep collecting votes and QCs of $g$.
       They can raise the entry lock chosen in step~3. \RR{Keep
       processing regional receipts (step~5 of \cref{alg:steady}).}
 \item If the proof is of type (b) or (c), submit it to L1. Do
       not wait for L1 to process it.
 \end{itemize}

\item \textbf{Status and entry.} \RR{Skip this step if session $(\beta,r)$ is
complete.} \RR{Wait until at least $a_v+\Delta$ and until
$\tagg{AUTH}_{r,g+1}$ is finalized. Set $e_v(\beta,r,g+1,0)$ to the time both
conditions hold, and read that authority from L1.} Then:
 \begin{itemize}
 \item Set the entry lock $Q_v$ to the highest-ranked QC
       (DATA\RR{, DRAIN,} or ANCHOR) that $v$ knows from generations
       $0,\ldots,g$\RR{, or to $Q_{\beta,r}^0$ if it knows none}.
       $Q_v$ never changes after this. QCs that arrive later do not
       update it.
 \item Enter $(\beta,r,g+1,0)$ as unfrozen and start its progress checks.
       \RR{Keep $A_{\beta,r,0,0}$, $H_{\beta,r}$, and the voting snapshot of session $(\beta,r)$.}
 \item Send $\langle\tagg{STATUS},\beta,r,g+1,0,Q_v,E\rangle_v$ and its
       validation bundle to $L(\beta,r,g+1,0)$. Here $E$ is the quit proof for
       $g$, which is also the entry proof of $g+1$.
 \item Process any buffered proofs for $g+1$.
 \end{itemize}

\item \textbf{New-view.} Run this step only at $L(\beta,r,g+1,0)$. Wait $2\Delta$
after its own entry. If $g+1$ is still current and unfrozen:
 \begin{itemize}
 \item Let $Q^*$ be the highest-ranked QC from generations $0,\ldots,g$
       that the leader knows, including the QCs in STATUS
       messages\RR{; if it knows none, let $Q^*=Q_{\beta,r}^0$}.
 \item Let $A_{\beta,r,g+1,0}$ be the block that $Q^*$ certifies. \RR{If
       $h(A_{\beta,r,g+1,0})<H_{\beta,r}$, generation $g+1$ is a DATA generation.
       Otherwise it is a drain generation.}
 \item Broadcast
       $\langle\tagg{NEWVIEW},\beta,r,g+1,0,A_{\beta,r,g+1,0},Q^*,E\rangle_{L(\beta,r,g+1,0)}$
       with its validation bundle. Choose only one opening block for
       $g+1$.
 \end{itemize}

\item \textbf{Anchor vote.} On receiving a NEWVIEW for $g+1$:
 \begin{enumerate}[label=(\alph*)]
 \item If it contains a quit proof for $g+1$, run Quit and stop.
 \item Validate its leader, entry proof $E$, opening block $A_{\beta,r,g+1,0}$,
       $Q^*$, and dependencies. If data are missing, request them and
       wait without voting. If the message is invalid, discard it.
 \item Vote only if all of these hold: $g+1$ is current and unfrozen; $v$
       has cast no ANCHOR vote in $g+1$; and $\rank(Q^*)\ge\rank(Q_v)$.
 \item To vote, forward the NEWVIEW to all replicas. Broadcast
       \[\langle\tagg{ANCHOR\mbox{-}VOTE},\beta,r,g+1,0,\Hash(A_{\beta,r,g+1,0})\rangle_v.\]
       Start no timer. Use no fresh position.
 \end{enumerate}
 Matching ANCHOR votes \RR{of total weight $>F$} form the ANCHOR QC.
 When $L(\beta,r,g+1,0)$ obtains it, the leader starts Propose in
 \cref{alg:steady} at $j=0$\RR{; in a drain generation, this is the
 DRAIN}. Every replica, including one that cast no ANCHOR vote,
 processes DATA \RR{and DRAIN} proposals backed by a valid ANCHOR QC
 under Vote.
\end{enumerate}
\end{protocol}

\paragraph{\defterm{Complete quit proofs.}}
For an admitted generation $g$ \RR{of session $(\beta,r)$}, a proof consists of
either:
\begin{enumerate}[label=(\alph*)]
\item distinct signatures \RR{of total weight $>F$} on the same
$\langle\tagg{BLAME},\beta,r,g,t\rangle$ for an admitted term, with that term's
admission witness;
\item two different admissible DATA \RR{or DRAIN} blocks signed by
$L(\beta,r,g,t)$ at the same $(\beta,r,g,t,j)$, with their admission \RR{and quota}
witnesses; or
\item two valid NEWVIEWs signed by $L(\beta,r,g,0)$ selecting different opening
blocks for generation $g$.
\end{enumerate}
\RR{A DATA block beyond the quota, or a DRAIN outside a drain generation,
is invalid. It is rejected and never counts for (b).}
Different openings count even when their blocks are comparable, because
they fix different starting heights $h(A_{\beta,r,g,0})$. The height formula then
maps the same position to different heights.

Different QCs for the same opening block are not equivocation. Evidence
from any admitted term can freeze its generation, including a term not
yet locally reached. Old-generation proofs are archival and may be
submitted to L1. \RR{Evidence from another session never freezes this
engine.} Future-generation proofs are buffered until entry. Neither
resets an unrelated recovery.

The Publish Layer uses these DATA and FINAL receipts. No publication
acknowledgement is required for regional commitment or Completion.
Participants may forward regional bundles directly; no external send
or Espresso confirmation delays a regional vote or timer.

The entry-lock test governs ANCHOR voting only. A replica that has
committed $T_{\beta,r}$ locally still takes part in recovery until it holds
$\FC_{\beta,r}(T_{\beta,r})$, because other members may need its votes. This continued-participation
rule applies until a verified RequestClose closes the fresh-action gate.

\section{Regional and global finality}
\label{sec:external}
Regional commitment and global acceptance are separate events. A node
follows a production branch $R_u$ and a continuous global DATA prefix
$G_u$. The branch may extend several sessions beyond $G_u$. A Force
closes that extension, never $G_u$ itself.

This section first states what a regional receipt proves
(\cref{sec:regional-finality}) and gives the public check of a receipt
(\cref{sec:receipt-kernel}). It then defines the global chain $G$ and
global finality (\cref{sec:global-chain}), and the rule by which every
verifier computes $G$ from the Espresso log and L1
(\cref{sec:exports}).

\subsection{Regional finality}
\label{sec:regional-finality}
\begin{definition}[\defterm{Regional finality}]
\label{def:regional-finality}
A DATA block $B$ has regional finality on a specified production branch
when its boundary history assigns it to its claimed region and session,
it has valid direct or same-session covering commitment evidence, and
the DATA path to it is valid and continuous. Each session on that path
supplies its own authority and commitment evidence.
\end{definition}
The region is derived from the boundary before its receipt is checked.
A valid receipt therefore cannot authorize another region's DATA or a
different session. Current live/off state does not replace historical
authority. When a Force closes an unconfirmed suffix, its receipts remain
evidence about the old regional engine, but they no longer let that
suffix extend the global prefix.

The progress statements in the remark below concern sessions that run
normally, with no RequestClose during the interval considered. Closing
sessions follow \cref{sec:retirement}; there a replica may keep a local
commitment for which no complete receipt ever forms.

\paragraph{Remark (Commitment by one replica and by all replicas).}
Here ``all replicas'' means all correct voters of the session that
produced $B$, with recovery authorities set in advance.
\begin{enumerate}[label=(\arabic*),itemsep=4pt]
\item \textbf{What a receipt proves.} A valid direct or covering receipt
proves that at least one correct replica has committed $B$. Other correct
replicas may still have running timers, or may have canceled their timers
on Quit. Regional finality does not claim that they have all committed.

\item \textbf{When all replicas commit.} Regional safety preserves $B$
through recovery: the selected anchor is $B$ or a descendant of $B$.
ANCHOR voting re-certifies that block but starts no commit timer.
Suppose a correct replica votes for and forwards a fresh DATA or DRAIN
proposal for a descendant $D$ at time $s$. If no correct replica freezes
that generation through $s+3\Delta_{i_{\beta,r}}$, all correct replicas
commit $B$ by that time. Replicas that have not completed the session
vote by $s+\Delta_{i_{\beta,r}}$; their timers commit $D$ and its
ancestors at most $2\Delta_{i_{\beta,r}}$ later. Completion may cancel a
timer, but it first commits $T_{\beta,r}$ and its ancestors, including
$B$. In particular, a DRAIN need not itself commit at a replica that
completes through the receipt for $T_{\beta,r}$.

\item \textbf{Why one honest leader is not enough.} An honest leader may
finish its term before its timers expire. A Byzantine successor can then
equivocate and freeze the generation. What is needed is the uninterrupted
commit window of (2). Under regional safety and the progress conditions
of \cref{sec:progress}, including permanent equivocation exclusion and
fair leader schedules, the engine eventually obtains such a window while
it keeps running, or completes through the terminal receipt. Earlier
recoveries do not revoke $B$. These conditions give no fixed bound on
when the window begins.

\item \textbf{When a replica can know.} A correct regional replica that
directly commits $D$ by its timer at time $c$ can conclude that all
correct regional replicas have committed $B$ if it stays Active in the
same generation through $c+2\Delta_{i_{\beta,r}}$, after processing the
messages due at that time. Other replicas' timers are due by
$c+\Delta_{i_{\beta,r}}$. A Quit that canceled one of those timers would
reach this replica by $c+2\Delta_{i_{\beta,r}}$ and freeze its
generation. A replica whose timer was canceled by Completion has already
committed $B$. Normal term rotation does not affect this test.

\item \textbf{What an external observer can know.} A receipt alone does
not show that the uninterrupted window occurred. The delivery bounds
cover complete submitted results, not every regional BLAME or quit event,
so the absence of a published quit proof does not prove that no replica
quit. Espresso evidence shows that the result is publicly available. A
bound on when every replica runs its local commit step would also need
the relevant online and receipt-processing conditions. This protocol
protects the external prefix against Force through Espresso acceptance;
it uses no external all-replicas-committed test.
\end{enumerate}

\subsection{Public receipt verification}
\label{sec:receipt-kernel}
This subsection answers one question: given a DATA block $B$, can anyone
check from public data that $B$ is regionally final
(\cref{def:regional-finality})? The answer is the check
$\tagg{CheckFINAL}$. It needs three things: the context of the session
that produced $B$, which the verifier derives itself; FINAL signatures
on a root block $D$; and the exact bytes that link $D$ back to $B$. It
does not need the proposal or recovery history that a publisher may
attach, which can name recovery authorities the verifier does not have.

\paragraph{Who uses it.}
The public interpreter (\cref{sec:exports}) and handoff verifiers use
$\tagg{CheckFINAL}$ for receipts that another session produced, such as
a historical anchor or a REGULAR terminal receipt. A handoff verifier
need not fetch the producer's optional recovery-admission path, but it
still checks the exact inherited DATA and path, the quota and the
canonical boundary and selection. Everything else keeps full
validation: own-session Vote, recovery and receipt-driven Completion;
Ready, which validates the target boundary, authority and exact
executable state; and admission in the target session.

\paragraph{The session context $C$.}
Every check is made against the \defterm{session context} $C$ of the
producing session $(\beta,r)$. $C$ consists of two authenticated
objects: the session's incoming boundary (the handoff record that opened
it, \cref{sec:witness}) and its initial authority $\tagg{AUTH}_{r,0}$.
They fix the following fields.
\begin{center}\small
\begin{tabularx}{\linewidth}{@{}>{\raggedright\arraybackslash}p{.36\linewidth}>{\raggedright\arraybackslash}p{.2\linewidth}Y@{}}\toprule
Field & Source & Example: A's $(7,20)$\\\midrule
chain identifier and protocol domain & boundary & this chain\\
producing pair $(\beta,r)$ & boundary & $(7,20)$\\
assigned region $i$ & boundary & A\\
anchor $A_{\beta,r,0,0}$: height, hash, original producer & boundary & DATA 100\\
fixed voting snapshot $\mathit{cfg}_r$: members, weights, $F_i$ & $\tagg{AUTH}_{r,0}$ & $\mathit{cfg}_{20}$\\
DATA quota $\TB_r$, hence $H_{\beta,r}$ & $\tagg{AUTH}_{r,0}$ & 10, so $H=110$\\\bottomrule
\end{tabularx}
\end{center}
The verifier derives $C$ from public history, never from the sender of a
record. The global finality rule (\cref{alg:global-rule}) installs a
regular boundary when it accepts the predecessor's terminal DATA, a
forced boundary at a Force, and fixes the region in
$\tagg{ResolveBoundary}$; the
initial checkpoint supplies the genesis boundary; a handoff verifier
derives the same boundary with $\tagg{VerifyHandoffHistory}$
(\cref{sec:verify-handoff}). $C$ contains no recovery authority
$\tagg{AUTH}_{r,g}$ with $g>0$ and no generation or term: FINAL does not
bind them, and the voting snapshot is the same in every generation of
the session.

\paragraph{The check.}
Write
\[
 \tagg{CheckFINAL}(C,B,K),\qquad K=(D,\Sigma,\pi).
\]
Its inputs are the following.
\begin{center}\small
\begin{tabularx}{\linewidth}{@{}>{\raggedright\arraybackslash}p{.05\linewidth}Y>{\raggedright\arraybackslash}p{.15\linewidth}>{\raggedright\arraybackslash}p{.22\linewidth}@{}}\toprule
Input & What it is & Supplied by & Example: A's $(7,20)$\\\midrule
$C$ & the session context above & derived by the verifier & anchor DATA 100, quota 10\\
$B$ & the candidate DATA to check & publisher & DATA 103\\
$D$ & the root DATA that carries the FINAL certificate: $B$ itself, or a same-session descendant of $B$ & publisher & DATA 105\\
$\Sigma$ & the FINAL signatures on $D$ & publisher & FINAL on DATA 105\\
$\pi$ & the exact hash-linked path from $D$ back to $B$, with full bodies & publisher & $105\to104\to103$\\\bottomrule
\end{tabularx}
\end{center}
The objects in $K$, read against $C$, form the \defterm{kernel} for $B$.
The check has two parts.
\begin{enumerate}[label=(\arabic*),itemsep=2pt]
\item \textbf{Signatures.} Write $\tagg{RawFINAL}_C(D,\Sigma)$ when the
bytes $D$ have the named hash and height, and $\Sigma$ contains valid
FINAL signatures from distinct members of $\mathit{cfg}_r$, with total
weight $>F_i$. The signed context must match $C$ in full: chain,
protocol domain, message type, producing pair, assigned region and
configuration, even where the displayed FINAL notation omits those
fields. ``Raw'' means a cryptographic and type check only. It does not
validate the block, include \cref{def:receipt}(iii), read a generation
or term, or check any attached admission witness.
\item \textbf{Path.} Each body on $\pi$ is an actual DATA or DRAIN body
with the correct hash, parent link and position; $\pi$ ends at $B$; $D$
and $B$ belong to the same session $(\beta,r)$; and
$h(B)\le H_{\beta,r}$. For a direct receipt, $D=B$ and $\pi$ is just $B$.
A hash without its payload is not enough, and a later session's receipt
never covers an earlier session's DATA.
\end{enumerate}
The check returns \emph{true} if both parts pass, \emph{false} if one
fails, and \emph{incomplete} if a required object is absent. The caller
handles \emph{incomplete} under the missing-object rule; it is not an
instruction to fetch optional sidecars. Required bytes must be supplied
or resolved from established public objects.

In the example, part (1) checks that the FINAL signatures on DATA 105
have weight $>F_A$ under $\mathit{cfg}_{20}$. Part (2) checks that
$105\to104\to103$ is correctly linked inside $(7,20)$ and that
$103\le110$. The check returns \emph{true}: DATA 103 is regionally final.

\paragraph{Why it suffices.}
\begin{proposition}[Soundness of the public FINAL kernel]
\label{prop:final-kernel-soundness}
Let $C$ be the authentic incoming context of session $(\beta,r)$, and
suppose $\tagg{RawFINAL}_C(D,\Sigma)$ holds. Then some correct replica of
that session committed exactly this DATA $D$ through its own commit
timer, and before doing so it fully validated $D$ and the inputs $D$
depends on. Hence $D$ is a validly admitted regional input under its
anchor, including the $\rho$ and $\xi$ conditions, whatever witness a
publisher happens to attach. The same holds for every same-session DATA
ancestor of $D$ on its hash-linked path.
\end{proposition}
This relies on the standing assumptions: the fault bound of the fixed
snapshot, the key-compromise bound, the cryptographic assumptions, sound
validation interfaces, and the regional rules for full validation and
for sending FINAL. The intuition is short. Weight above $F_i$ contains a
correct signer; a correct replica signs FINAL only after its own timer
commits the block; and it starts that timer only after full validation.
Commits caused by a receipt send no FINAL, so they cannot be the source
of the first correct signature. The precise statement and proof are in
Appendix~\ref{app:final-kernel-proof}.

So if $\tagg{CheckFINAL}$ returns \emph{true}, part (2) places $B$ on
$D$'s validated path and the proposition makes $B$ regionally final.
The proposition shows that a valid witness \emph{exists}; it does not
certify the one a publisher attaches. Because FINAL binds no generation
or term, a publisher cannot force a fetch of $\tagg{AUTH}_{r,g}$ by
naming a different $g$, and optional archival validation never gates
acceptance of a complete kernel.

\paragraph{What it does not check.}
A \emph{true} result says that $B$ is regionally final, and nothing more.
\begin{itemize}[itemsep=2pt]
\item \textbf{Global acceptance.} Whether $B$ exactly extends the
current $G$, the current branch permission and export duty, the strict
precedence of Espresso selection references, and the cutoff after a
Force are checked by the caller (\cref{sec:exports}). The interpreter
must still obtain the complete L1 control history; FINAL cannot vouch
for it.
\item \textbf{Uniqueness.} The threshold $>F_i$ alone is not a
quorum-intersection argument. Regional safety, proved separately under
its premises, makes competing valid commitments compatible.
\item \textbf{Execution.} The kernel certifies only the hash-bound
input history. Unbound native context, state objects, execution output
and later $C_{\rm pub}$ publications keep their own execution and DA
checks; a receipt is not a state object, so Ready still requires the
exact native state. A native claim bound by the DATA hash and checked by
the regional validity interface is covered, and only that claim.
\item \textbf{Delivery.} The kernel is a subset of a complete result.
It does not relax the delivery and retention rules, adds no deadline,
and proves nothing about request inclusion.
\end{itemize}

\subsection{The global chain}
\label{sec:global-finality}
\label{sec:global-chain}
The \defterm{global chain} $G$ is the sequence of DATA blocks that every
verifier accepts. It starts at the genesis anchor $A_0$ and passes
through a sequence of sessions. Each session on the way contributes one
contiguous \defterm{segment}, and the chain records which session owns
each segment and why that session was opened.

\paragraph{Session records.}
A \defterm{session record} is
\[
 \sigma=(\beta,\ r,\ i,\ A,\ H,\ q^{\rm open},\ \mathit{cause}),\qquad
 \mathit{cause}\in\{\tagg{GENESIS},\ \tagg{REGULAR},\ \tagg{FORCE}(q_f)\}.
\]
Here $(\beta,r)$ is the session, $i$ its assigned region, $A$ its
inherited anchor, $H=h(A)+\TB_r$ its last DATA height, and
$q^{\rm open}$ the Espresso height at which its region was selected. The
cause says how the session was opened; for a Force, $q_f$ is the Espresso
block whose prelude executed it. A session record fixes the session
context $C_\sigma$ of \cref{sec:receipt-kernel}: the pair, region and
anchor are fields of $\sigma$, and $\mathit{cfg}_r$ and $\TB_r$ come from
$\tagg{AUTH}_{r,0}$. The L1 bound $\lambda^{\rm open}_{\beta,r}$ is not
part of $\sigma$; only regional validation uses it.

\begin{definition}[\defterm{Global chain}]
\label{def:global-chain}
A global chain is a sequence
\[
 G=A_0\cdot s_0\cdot s_1\cdots s_m
\]
together with session records $\sigma_0,\ldots,\sigma_m$, where
$\sigma_0$ has cause \tagg{GENESIS} and $r_0=0$, and each segment $s_k$ is
a possibly empty sequence of DATA blocks such that:
\begin{enumerate}[label=(\roman*),itemsep=2pt]
\item \textbf{Continuity.} $A_k$ is the last block of
$A_0\cdot s_0\cdots s_{k-1}$. Every block of $s_k$ has the block before it
in $G$ as its parent and the next DATA height.
\item \textbf{Ownership.} Every block of $s_k$ is produced by session
$(\beta_k,r_k)$, has height at most $H_k$, and is regionally final under
$C_{\sigma_k}$ (\cref{def:regional-finality}).
\item \textbf{Boundaries.} $r_{k+1}=r_k+1$ if $s_k$ is nonempty and
$r_{k+1}=r_k$ if it is empty. If $\sigma_{k+1}$ is \tagg{REGULAR}, then
$s_k$ ends at height $H_k$ and $\beta_{k+1}=\beta_k$. If it is a
\tagg{FORCE}, then $\beta_{k+1}>\beta_k$. Every segment $s_k$ with $k<m$
that ends below $H_k$ is followed by a \tagg{FORCE}.
\end{enumerate}
\end{definition}
A segment stops below its endpoint only because a Force closed its
version. A Force can also hit a finished segment whose successor is not
yet selected; the next session then keeps the index $r_k+1$ but has
cause \tagg{FORCE}. The version can rise by more than one between two
records when several Forces act before a region is selected.

The rule of \cref{sec:exports} fixes which global chain each verifier
holds at each Espresso height. In the running example the chain
contains, among others, the two entries
\begin{center}\small
\begin{tabular}{@{}llll@{}}\toprule
Entry & Session record $\sigma$ & Cause & Segment $s$\\\midrule
$k$ & $(7,20,\text{A},\ 100,\ 110,\ q_{\rm A})$ & \tagg{REGULAR} & $101,\ldots,108$\\
$k+1$ & $(8,21,\text{B},\ 108,\ 108+\TB_{21},\ q_s)$ & \tagg{FORCE}$(q_f)$ & $109',110',\ldots$\\\bottomrule
\end{tabular}
\end{center}
A's segment stops at 108 because the Force closed version 7 before 109
was accepted. B's regular session $(7,21)$ and its DATA 111--115 do not
appear, because $G$ never passed through them. The production branch
$R$ may run ahead of $G$; the global chain records only what was
accepted.

\begin{definition}[\defterm{Global finality}]
\label{def:global-finality}
A DATA block $B$ is \defterm{globally final at Espresso height $q$} when
$B\in G_q$, the global chain output by the rule of \cref{sec:exports} at
$q$. Observer $u$ treats $B$ as globally final once it has authenticated
the finalized Espresso history through $J_q$ and the complete L1 control
history that applies to it, and has applied the rule through a step that
appends $B$.
\end{definition}

The rule only appends, so $G_{q-1}$ is a prefix of $G_q$ for every $q$
(\cref{prop:global-prefix}). No later Force removes a globally final
block or changes its session record, and no extra waiting period
follows. Espresso inclusion alone, or a receipt alone, does not make a
block globally final: an old-version block may be finalized in Espresso
as raw bytes and still never enter $G$.

This is finality of the input prefix. Espresso neither executes the
native transactions nor proves their resulting state; the regional
validity contract and the execution-verification interface are
responsible for that. Finality of a state claim or withdrawal on L1 is a
later, separate stage. Publication lag never permits skipping globally
final input.

\subsection{The global finality rule}
\label{sec:exports}
The global finality rule computes the global chain from public data. It
is a deterministic function of its input, so every verifier that reads
the same history computes the same chain.

\paragraph{Input.}
\label{sec:logical-carrier}
\begin{enumerate}[label=(I\arabic*),itemsep=2pt]
\item \textbf{The Espresso log.} The finalized blocks $J_1,J_2,\ldots$,
each with its timestamp $\theta_q$, its hash, and the set $M_q$ of
regional objects in this chain's namespace. $M_q$ is what the EXPORT
records of $J_q$ carry (\cref{sec:delivery}). The rule does not depend on
how records encode or group the objects. It uses two kinds:
 \begin{itemize}[itemsep=1pt]
 \item block bodies
 $B=(\beta,r,k,h,\mathit{type},x,\Hash(B_{\rm parent}),\rho,\xi)$ of type
 DATA or DRAIN, identified by their hash;
 \item FINAL signatures: a signature of a member on the hash of a block
 $D$, in the signed context of $D$'s session.
 \end{itemize}
Write $\mathcal O_q=M_1\cup\cdots\cup M_q$ for everything published
through $J_q$.
\item \textbf{L1.} The authorities $\tagg{AUTH}_{r,0}$ and the complete
ordered STATUS and MSIG FORCE records with their activation thresholds.
They give the eligibility snapshot $\chi_E(q)$ after the prelude of $J_q$
(\cref{sec:l1}).
\item \textbf{Start.} The genesis checkpoint of \cref{sec:l1}: the chain
identifier, the region order $\mathcal L_R$, the genesis anchor $A_0$ and
the first replay position $q_{\rm init}$. A later trusted checkpoint
supplies an earlier output of the rule instead.
\end{enumerate}
Nothing else is an input. Proposals, QCs, recovery authorities
$\tagg{AUTH}_{r,g}$ with $g>0$ and native sidecars may also appear in
$M_q$; the rule ignores them. The outer fields of an EXPORT, such as its
expected base, interval and version, are routing hints and are ignored
too. A malformed or irrelevant object grants nothing and blocks nothing.

The rule uses one test derived from the input. For a session record
$\sigma$ and a DATA block $B$, write $\mathrm{Fin}_q(\sigma,B)$ when
$\tagg{CheckFINAL}(C_\sigma,B,(D,\Sigma,\pi))$ returns \emph{true} for
some block $D$ and path $\pi$ whose bodies all lie in $\mathcal O_q$, with
$\Sigma$ the set of FINAL signatures on $D$ in $\mathcal O_q$. The
signatures may come from different records and different blocks; only
their total weight counts.

\paragraph{Output.}
For every Espresso height $q\ge q_{\rm init}$ the rule outputs
\[
 \mathrm{Out}_q=\big\langle(\sigma_0,s_0),\ldots,(\sigma_m,s_m)\big\rangle
 \circ\mathit{tail},\qquad \mathit{tail}\in\{\varepsilon,\ \omega\}.
\]
The list is a global chain (\cref{def:global-chain}), and
$G_q=A_0\cdot s_0\cdots s_m$. The tail says who owes the next DATA.
\begin{itemize}[itemsep=2pt]
\item $\mathit{tail}=\varepsilon$: $\sigma_m$ is the \defterm{export
session}, and $s_m$ may still grow. Its accepted count is $a=|s_m|$.
\item $\mathit{tail}=\omega=(\mathit{cause},\beta,r',i_{\rm src},A,q_0)$:
a \defterm{pending boundary}. The last segment is finished. The next
session will be $(\beta,r')$ and continues after $A=\operatorname{last}(G_q)$,
but its region is not yet selected; the search starts at Espresso height
$q_0$, after source region $i_{\rm src}$.
\end{itemize}
The current branch version is the $\beta$ of the tail when the tail is
pending, and $\beta_m$ otherwise. Before the first session is selected,
the list is empty and $G_q=A_0$.

\begin{protocol}{alg:global-rule}{Global finality rule: from $\mathrm{Out}_{q-1}$ to $\mathrm{Out}_q$}
Initially $\mathrm{Out}=\langle\rangle\circ(\tagg{GENESIS},0,0,\bot,A_0,q_{\rm init})$.
For each finalized block $J_q$ in order, start from $\mathrm{Out}_{q-1}$,
apply Step~1 once, and then apply Steps~2--4 while one of them applies.
The result is $\mathrm{Out}_q$. Below, $\beta$ is the current branch
version and $\operatorname{last}(G)$ the last block of the current chain.
\begin{enumerate}[label=\textbf{\arabic*)},itemsep=3pt]
\item \textbf{Force} (prelude). Among the valid MSIG FORCE records with
$q(\theta_a)=q$, take the first in $(\theta_a,c)$ order whose target
version is $\beta$. If there is one:
 \begin{itemize}[itemsep=1pt]
 \item if the tail is $\varepsilon$, close $s_m$ and set the tail to
 $\big(\tagg{FORCE}(q),\ \beta+1,\ r',\ i_m,\ \operatorname{last}(G),\ q\big)$,
 with $r'=r_m$ if $a=0$ and $r'=r_m+1$ otherwise;
 \item if the tail is a pending $\omega$, keep its $r'$, $i_{\rm src}$ and
 $A$, and set its cause to $\tagg{FORCE}(q)$, its version to $\beta+1$ and
 its $q_0$ to $q$.
 \end{itemize}
All other MSIGs due at $J_q$ have no effect.
\item \textbf{Select} ($\tagg{ResolveBoundary}$). If the tail is
$\omega=(\mathit{cause},\beta,r',i_{\rm src},A,q_0)$ and some
$q'\in[q_0,q]$ has $\tagg{NextLive}(i_{\rm src},\chi_E(q'))\ne\bot$ under
$\tagg{AUTH}_{r',0}$, let $q_s$ be the least such $q'$ and
$i=\tagg{NextLive}(i_{\rm src},\chi_E(q_s))$. Replace $\omega$ by the
entry $\big((\beta,r',i,A,h(A)+\TB_{r'},q_s,\mathit{cause}),\ \langle\rangle\big)$;
the tail becomes $\varepsilon$.
\item \textbf{Accept.} If the tail is $\varepsilon$ and $\mathcal O_q$
contains a DATA block $B$ such that
 \begin{enumerate}[label=(\alph*),itemsep=1pt]
 \item $B$ is signed in this chain's domain by session $(\beta_m,r_m)$;
 \item $\mathrm{parent}(B)=\Hash(\operatorname{last}(G))$ and
 $h(B)=h(\operatorname{last}(G))+1\le H_m$;
 \item $\mathrm{Fin}_q(\sigma_m,B)$ holds; and
 \item $\xi(B)$ names a block $J_{q'}$ of the log, with its hash and
 eligibility snapshot, and $q^{\rm open}_m\le q'<q$,
 \end{enumerate}
then append $B$ to $s_m$. If several blocks qualify, take the one with
the least hash.
\item \textbf{Complete.} If the tail is $\varepsilon$ and $s_m$ has
reached height $H_m$, set the tail to
$\big(\tagg{REGULAR},\ \beta_m,\ r_m+1,\ i_m,\ \operatorname{last}(G),\ q_0\big)$
with $q_0=\max\{\operatorname{height}(\xi(\operatorname{last}(G))),\ q^{\rm open}_m+1\}$.
\end{enumerate}
\end{protocol}

\Cref{alg:global-rule} is the public core of \cref{alg:force} and
\cref{alg:external}. Step~1 is the Force of \cref{alg:force}; the STATUS
records of that prelude are already folded into $\chi_E(q)$. Step~2 is
the historical selection of \cref{sec:witness}, shared by REGULAR and
MSIG boundaries. \Cref{sec:force,sec:observers} explain these steps and
the waiting rules around them.

\paragraph{Why each check.}
\begin{itemize}[itemsep=2pt]
\item \textbf{Export session only} (a). Only the session that owes the
next height may extend $G$. A closed version has no export session, so
after a Force its DATA never enters $G$. A regular successor that runs
ahead becomes the export session only after its predecessor's terminal
DATA is accepted.
\item \textbf{Continuity} (b). $G$ has no gap and no fork. DRAIN never
enters $G$, because it is not DATA and does not advance the height.
\item \textbf{Regional finality} (c). By
\cref{prop:final-kernel-soundness}, a correct replica of the session
committed $B$ after fully validating it. The rule needs no proposal or
recovery history.
\item \textbf{Selection reference} (d). A reference cannot justify itself
through the block that carries it. Because $q^{\rm open}\ge q_f$ for a
session opened by a Force at $J_{q_f}$, no new-version DATA is accepted
in $J_{q_f}$ itself. For the same reason, when the block that accepts
terminal DATA is also the selection height $q_s$ of the successor, the
successor's DATA cannot be accepted in that block.
\end{itemize}

\paragraph{Determinism and accumulation.}
Each step reads only the current output and the input, so
$\mathrm{Out}_q$ is a function of the input through $J_q$. Under regional
safety at most one block qualifies in Step~3, since two regionally final
DATA of one session at one height are identical; the hash tie-break only
makes the rule total. The set $\mathcal O_q$ only grows. An object stays
usable after the block that carried it: if DATA 109 is published before
108 is accepted, Step~3 accepts 109 right after 108, in the same or a
later block, without republication. A wrong hint or an invalid neighbor
in the same record removes nothing.

For example, at the Force of the running example the tail is
$\varepsilon$, $\sigma_m$ is A's $(7,20)$ and $s_m=101,\ldots,108$. Step~1
closes $s_m$ and sets the tail to
$(\tagg{FORCE}(q_f),8,21,\text{A},108,q_f)$, since $a=8>0$. Step~2
selects B at the first eligible $q_s\ge q_f$. B's new DATA $109'$ carries
a reference at least $q_s$, so Step~3 can accept it from
$J_{q_f+1}$ on. A's old 109 is still in $\mathcal O_q$ but fails (a).

\paragraph{Waiting and partial progress.}
\label{sec:subpositions}
\label{sec:terminal-commit}
The rule is stated per Espresso block, but its steps are atomic, and an
observer may stop between them. Steps~1 and~2 need L1 inputs: the
applicable control history and $\tagg{AUTH}_{r',0}$. If these are not yet
retrievable, the observer waits at that step; missing evidence never
counts as \emph{no eligible region}. With complete evidence and no
eligible region, the boundary stays pending and the next block is
processed. Step~3 without a candidate is not a wait either. Step~4 runs
before Step~2, so terminal DATA is accepted before the successor's
authority is needed.

The chain held after any step within $J_q$ is a prefix of $G_q$ and may
be exposed at once. An implementation persists the block and the last
completed step, its \defterm{replay subposition}, so that a restart
neither repeats nor skips a step. A duplicate delivery of a source block
resumes at that subposition.

\paragraph{What the transport must supply.}
\label{sec:transport-scope}
The rule can accept only what $\mathcal O_q$ contains. The Publish Layer
must therefore put the objects of a complete regional result into
Espresso: the DATA bodies, the FINAL signatures and any covering path.
The delivery bound $T_{\rm pub}$ of \cref{ass:services} applies to
complete results; a hash, a READY message or a receipt without its
payload is not one. The encoding of EXPORT records is the Publish Layer's
choice, provided every reader recovers the same $M_q$ from the namespace
content. An encoding that hides required objects, or a private cache,
does not supply them.

\subsection{What block references do not tell}
\label{sec:late-receipt}
A block's L1 reference and Espresso selection reference are fixed when it
is proposed. A later recovery or DRAIN may commit the block long
afterwards. Neither reference shows global finality or authorizes a
Force. Regional timers use the real-time clock of \cref{ass:network}.
Controls take effect at the Espresso threshold of an explicit L1 record.
For example, a block that references an Espresso block with time 100 may
be committed much later after a recovery; it was not committed at time
100.

\section{Cross-region rotation}
\label{sec:xprotocols}
This section describes what each regional node does across sessions:
how it learns about handoffs and Forces, how it prepares and enters a new
session, and how it shuts down an old one. The public rules that decide
handoffs and Forces are in \cref{sec:force,sec:observers}; this section
consumes their results.

\paragraph{What B does at the same time.}
Return to the running example. A's last DATA is 110. B holds its regional
receipt and produces in $(7,21)$ from 111, while $G$ is still 108. L1 now
announces a future Force that targets branch version 7. B can do two
things at once:
\begin{center}\small
\begin{tabular}{@{}lll@{}}\toprule
Program & What it does now & Continues after\\\midrule
old $(7,21)$ & already entered; runs the regional BFT & 110 (fixed)\\
possible new $(8,21)$ & cannot produce; fetches data and state & 108 for now\\\bottomrule
\end{tabular}
\end{center}
If $G$ grows from 108 to 109 before the Force takes effect, B starts
preparing state after 109 instead. The old program's anchor stays 110.
When the Force takes effect and B is still selected: branch version 7
can no longer append to $G$; the new program's anchor is fixed to the
current $G$; B checks that its prepared material matches this anchor;
and if so, B enters the new program while the old one drains its timers,
evidence and data duties. Preparing early never grants the right to
produce early.

\subsection{Cross-region architecture}
\label{sec:runtime-modules}
Each correct regional replica runs a \defterm{Session Manager}. It reads
finalized L1 and Espresso history, receives regional results, and manages
independent \defterm{session programs}, one per pair $(\beta,r)$. A
program holds the regional engine of its session together with its
preparation and storage state. The Session Manager is local logic
replicated at every node, not a trusted controller. A region can run an
old program while preparing a possible successor and handing completed
data duties to storage and publication workers.

There are three control roles. The Multisig Controller detects health
problems and submits L1 controls. The Session Manager reconciles
verified public changes with the local programs. Each program's
Ready/Enter gate checks its exact authorized start. Shared source
verification and \cref{alg:force,alg:external} produce the public state
that all three use.

\begin{table}[htbp]\centering\small
\caption{Cross-region responsibilities.}
\label{tab:runtime-modules}
\begin{tabularx}{\linewidth}{@{}p{.22\linewidth}Y Y@{}}\toprule
Executor & Inputs & Outputs and authority limit\\\midrule
Multisig Controller & Verified progress, regional health and operations
reports & Future L1 ON/OFF and MSIG; no choice of retained DATA.\\
Common interpreter & Complete L1/Espresso history and public result
evidence & Eligibility, fixed boundaries and $G$; no health diagnosis.\\
Session Manager & Verified public changes, receipts and local callbacks
& Prepare, enter, request close and transfer data duties; no independent Force.\\
Ready/Enter gate & Fixed handoff, exact anchor state and authority
& Waiting or atomic initial entry; preparation grants no early voting right.\\\bottomrule
\end{tabularx}
\end{table}

Operational regions keep watching both public streams while idle. Their
verified cursor, candidate programs and storage tasks stay separate. A
program's local phase does not decide whether DATA may advance $G$, and
a monitoring report has no direct control effect.

\begin{figure}[p]\centering
\resizebox{\linewidth}{!}{%
\begin{tikzpicture}[
 x=1cm,y=-1cm,font=\small,
 box/.style={draw=rrBlue!65,rounded corners=3pt,fill=rrBlue!3,
   align=center,inner sep=7pt,text width=5.3cm,minimum height=1.25cm},
 shared/.style={box,draw=rrTeal,fill=rrTeal!4,text width=11.6cm},
 prog/.style={box,text width=3.75cm,minimum height=3.15cm},
 flow/.style={-{Stealth[length=2.2mm]},semithick},
 lab/.style={fill=white,inner sep=3pt,align=center,font=\footnotesize}
]
\path[use as bounding box] (0,0) rectangle (14.8,19.0);
\node[font=\large\bfseries] at (7.4,.45) {Explicit control; concurrent local sessions};
\node[box] (gov) at (3.4,2.1) {\textbf{Multisig Controller}\\
 Observe health and accepted progress.\\Submit future controls on L1.};
\node[box] (l1) at (11,2.1) {\textbf{L1}\\
 Immutable AUTH and ordered\\ON / OFF / MSIG records.};
\draw[flow] (gov.east)--node[lab,above]{signed control}(l1.west);

\node[box] (pub) at (3.4,4.45) {\textbf{Publish Layer}\\
 Leader first; backup participants.\\Deliver complete regional results.};
\node[box] (esp) at (11,4.45) {\textbf{Espresso}\\
 Ordered records and finality.\\Authenticated timestamps\\and data retrieval.};
\draw[flow] (pub.east)--node[lab,above]{complete result}(esp.west);

\node[shared] (replay) at (7.4,7.2) {\textbf{Shared validation and replay: Protocols 5--6}\\
 Verify both histories. Apply due controls before payload.\\
 Accept continuous DATA into $G$; derive handoffs and eligibility.\\
 Health diagnosis stays with the Multisig Controller.};
\draw[flow] (l1.east)--(14.55,2.1)--(14.55,6.7)--(replay.east);
\draw[flow] (esp.south)--(replay.north -| esp.south);
\draw[flow] (replay.west)--(.25,7.2)--(.25,2.1)--(gov.west);
\node[lab,rotate=90] at (.25,4.9) {verified progress};

\node[shared] (manager) at (7.4,10.35) {\textbf{Session Manager: Protocol 3}\\
 Serialize control and lifecycle events; run fetches concurrently.\\
 Reconcile fixed handoffs, future candidates and version-wide closure.\\
 Forward complete results and preserve evidence across worker exit.};
\draw[flow] (replay.south)--node[lab,right]{verified state and closure}(manager.north);

\node[prog,draw=rrTeal,fill=rrTeal!5] (old) at (2.6,13.5)
 {\textbf{Old session $(\beta,r)$}\\
 Active or Frozen.\\After Force, gate new work,\\drain timers and retire safely.};
\node[prog] (new) at (7.4,13.5)
 {\textbf{Next session $(\beta',r')$}\\
 Waiting: exact anchor,\\state and authority.\\Ready: atomic entry.};
\node[prog,draw=black!45,fill=black!3] (archive) at (12.2,13.5)
 {\textbf{Completed / Terminated}\\
 Keep DATA, receipts,\\source and closure proofs.\\Transfer publication duties.};
\draw[flow] (manager.south)--(7.4,11.65)--(2.6,11.65)--(old.north);
\draw[flow] (manager.south)--(new.north);
\draw[flow] (manager.south)--(7.4,11.65)--(12.2,11.65)--(archive.north);

\node[shared,minimum height=1.35cm] (gate) at (7.4,17.1)
 {\textbf{Ready / Enter and the original regional engine: Protocols 4, 1--2}\\
 Fixed handoff + matching executable state $\to$ atomic Enter.\\
 STATUS, $2\Delta$ wait, NEWVIEW, ANCHOR QC, then DATA / DRAIN.\\
 EL supplies payloads and state; client-request resubmission remains in EL.};
\draw[flow] (new.south)--node[lab,right]{handoff still applicable}(gate.north);
\end{tikzpicture}%
}

\caption{Governance control and concurrent regional programs. Only L1
control records can authorize a Force. The same public replay rules run
at regional nodes, observers and certifiers.}
\label{fig:runtime-modules}
\end{figure}
\FloatBarrier

\subsection{Multisignature governance interface}
\label{sec:runtime-controls}
The multisignature service decides whether to intervene. It can observe
regional health, the current export session, accepted DATA progress and
operational reports. Tests for missing blocks or slow delivery belong to
this policy; they are not consensus transitions or certificate
predicates. Its interface is
\[
 \text{verified observations and reports}
 \ \longrightarrow\
 \text{future L1 STATUS and MSIG records}.
\]
Other participants may run monitors and send reports. Only finalized,
authorized L1 controls can change public eligibility or cause a Force.
A report delivered to one region cannot change that region's replay
result.

To exclude a failed region and force a handoff, the controller schedules
OFF and MSIG. To force a handoff without changing eligibility, it
schedules MSIG alone. ON restores eligibility for future selections.
Restoring a region, serving every region, and fairness among regions are
choices of the controller, not obligations of the protocol.

The certifier checks a control's authorization, source completeness,
activation position and target version. It does not check whether an
internal missing-block threshold was met. An instruction means the same
thing regardless of the report that motivated it. The committee cannot
supply a preferred retained block; replay uses the actual $G$ at the
start of the activating block.

Detection thresholds, report quality, signing policy, response time and
intervention frequency are left to a separate governance design. Hard
Forces may follow one another without limit, even during preparation or
before the new session has produced DATA. Such interference can stop
progress, but it cannot change accepted output.

If a region is offline and no valid MSIG is submitted, the consensus
layer does not force a switch. The L1 service bound starts at actual
submission and excludes detection, decision and signature collection.
Recovery liveness therefore needs the committee-response and
service-opportunity conditions of \cref{sec:progress}.

\subsection{The Session Manager}
\label{sec:session-manager}
The Session Manager processes events that change its local control
state one at a time. Data retrieval and state restoration may run in
parallel in the background. \Cref{tab:manager-inputs} lists its inputs.

\begin{table}[htbp]\centering\small
\caption{Session Manager inputs.}
\label{tab:manager-inputs}
\begin{tabularx}{\linewidth}{@{}>{\raggedright\arraybackslash}p{.24\linewidth}Y@{}}\toprule
Event & Validation and routing\\\midrule
\texttt{L1Finalized} & Authenticate the complete control prefix and
authorities; queue future controls without activating them.\\
\texttt{EspressoFinalized} & Authenticate complete ordered history;
consume the state transition produced by \cref{alg:external,alg:force}.\\
\texttt{RegionalFinalized} & Validate the DATA, complete FINAL receipt
and required ancestry; route to the producing pair and offer the result
to the Publish Layer.\\
\texttt{MaterialAvailable} & Check the target pair, the anchor's producing
pair, height and hash, the requested object and the verified result;
update only the matching preparation or fixed handoff.\\
\texttt{CoreEvent} & Accept an event from the named local engine;
record its view, quit, recovery, Completion or retirement progress.\\\bottomrule
\end{tabularx}
\end{table}

\paragraph{The Closed map.}
The Session Manager keeps $\tagg{Closed}[\beta]$, with its authenticated
proof, even if it has no program for that branch version. This record is
public source evidence, not a new committee certificate. Deleting a
program cannot delete the map or its pending control-routing duties.
Before processing a complete old-session proposal, ReadyBundle, authority
packet or archived message, the dispatcher checks the map. If closure
forbids entry or a fresh action, the Session Manager serves and relays
the full closure, using the packet's verified historical membership
context, as specified in \cref{sec:retirement}. It never silently drops
the packet. This response signs no BFT vote and does not revive a
terminated engine.

\paragraph{Receipts.}
One FINAL signature is not a \texttt{RegionalFinalized} event. Existing
participants collect matching signatures into the receipt of
\cref{def:receipt}; this is not a consensus round. A correct holder
delivers the complete result to the Session Manager and to the entered
source engine, multicasts the complete receipt and dependencies to the
fixed regional membership, and offers it to the Publish Layer. The source
engine applies the receipt and Completion rules of \cref{alg:steady}. A
program that has not entered only caches the result: a receipt cannot
create an engine or skip handoff checks. $T_{\rm pub}$ starts when a
complete result becomes available.

\paragraph{Outputs.}
The outputs are:
\texttt{Prepare/Fetch}, naming the target and the exact requested
material; \texttt{Activate}, invoking \cref{alg:bootstrap} once;
\texttt{RequestClose}, carrying a verified closure cause;
\texttt{OfferResult}, preserving the original DATA and receipt fields;
and \texttt{Persist/Release}, recording state and transferring resource
responsibilities. A cached callback is material, not an activation
command. A late result can still be useful to storage even if it does not
match the program's current preparation target.

\begin{protocol}[H]{alg:rotation}{Handle cross-region events and manage sessions}
Each correct regional replica runs this handler independently.
\begin{enumerate}[label=\textbf{\arabic*)},itemsep=3pt]
\item \textbf{Validate the event.} Route authenticated public input,
regional results and local callbacks by their stated context. Keep
pending material-fetch requests and source evidence. Reject a claimed
completion whose material or session does not match its request.
\item \textbf{Consume public replay.} For new public history, run or
consume \cref{alg:external}, including \cref{alg:force} in its canonical
order. Use its resulting $\beta$, $G$, eligibility snapshots, duty and
boundary records. This handler never increments $\beta$, chooses a
different cutoff or modifies $G$ on its own.
\item \textbf{Reconcile handoffs.} From valid terminal receipts, derive
regular boundaries while their branch version stays open. Consume fixed
Force boundaries from replay. Verify them by \cref{sec:witness}. Prepare
selected local targets, and send complete handoff material to the
selected members when available. A future Force does not block a current
regular handoff.
\item \textbf{Reconcile candidates.} Recompute possible future targets
from the current export context, $G$, the eligible regions and known
future controls. Update only unfixed preparations. A different target
pair uses a different record. Pause speculative work that is no longer
useful. Terminate an unentered program only when verified facts make its
target impossible or assign its fixed handoff to another region.
\item \textbf{Handle closure and results.} Store every verified closure
in the durable $\tagg{Closed}[\beta]$ map. Gate new entry and fresh
actions in every phase. For each known session snapshot of that branch
version, including unentered and archived programs, run the complete
closure relay and idempotent \texttt{RequestClose} of
\cref{sec:retirement}. If the snapshot is not yet known, keep a pending
relay obligation. Any later old-session packet with the needed context
triggers this control response before ordinary processing. Keep signed
routing, historical evidence and publication duties. Apply terminal
receipts to entered source engines under the Completion rule.
\item \textbf{Check and enter.} For each selected local Waiting program
with a fixed handoff, check \cref{def:ready}. In the same local event,
recheck public validity and the durable entry flag, record Ready, and
invoke \cref{alg:bootstrap} Enter atomically. A failed check leaves the
program Waiting or applies a justified cancellation; it never leaves a
partly initialized engine.
\item \textbf{Persist and dispatch.} Durably record state changes, entry
and exit records, and required outgoing actions before exposing them.
Dispatch asynchronous work with its full target and anchor context.
Repeated events resume outstanding actions without creating another
logical entry.
\end{enumerate}
\end{protocol}

The public interpreter has one mutation path: \cref{alg:force} defines
control and Force transitions, and \cref{alg:external} authenticates and
applies them and accepts DATA from EXPORTs. \Cref{alg:rotation} consumes
those results. \Cref{alg:steady,alg:recovery} update only the named
session's regional engine. An implementation may place these modules in
one process without changing this division of authority.

\subsection{Session lifecycle}
\label{sec:session-lifecycle}
The Session Manager keeps the following record for each local program:
\[
\begin{aligned}
\tagg{Program}[\beta,r]=(&\mathit{phase},\mathit{preparing\_plan},
 \mathit{fixed\_handoff},\mathit{anchor\_block},\\
 &\mathit{execution\_state},\mathit{authority\_and\_bounds},
 \mathit{core\_state},\mathit{pending\_fetches},\\
 &\mathit{entered},\mathit{close\_requested},\mathit{closure\_evidence},\\
 &c_v,d_v,\mathit{publication\_and\_retention\_handles}).
\end{aligned}
\]
Its six phases describe the local program, not public finality.
\begin{center}\small
\begin{tabularx}{\linewidth}{@{}>{\raggedright\arraybackslash}p{.15\linewidth}Y@{}}\toprule
Phase & What the program is doing\\\midrule
Waiting & Preparing a possible or selected handoff. It may fetch data,
recover state and cache evidence. It has not entered its regional
engine, signed a regional vote, or started a regional commit or progress
timer.\\
Ready & A fixed handoff, its exact anchor state and all entry
dependencies have passed the gate below. Enter follows in the same
serialized event, with no extra scheduling delay.\\
Active & Entered an unfrozen regional generation. Initial STATUS, the
$2\Delta$ wait, NEWVIEW and ANCHOR certification all happen here. The
engine runs \cref{alg:steady,alg:recovery}, restricted only by the
retirement interface of \cref{sec:retirement} once closure is requested.\\
Frozen & The current generation has a valid quit proof, and the evidence,
receipt and recovery rules of \cref{alg:recovery} apply. A Force alone
never produces a quit proof or this phase.\\
Completed & The engine holds a valid terminal receipt and has run
Completion: regional production has stopped and timers are canceled.
Data delivery and storage duties may remain.\\
Terminated & The program has exited. It runs no BFT timer and issues no
BFT signature. Storage and publication workers keep its required evidence
and duties. An exit record stops late messages from restarting it.\\\bottomrule
\end{tabularx}
\end{center}

\paragraph{Transitions.}
The normal path is Waiting $\to$ Ready $\to$ Active. Active becomes
Frozen on Quit; Frozen becomes Active again on recovery admission, with
$g$ advanced and $t=0$; Active or Frozen becomes Completed on terminal
Completion. Waiting or Ready may terminate when its target is definitely
invalid. Completed may terminate after its responsibilities are
transferred; if RequestClose began while the program was Active or
Frozen, the drain deadline of \cref{sec:retirement} still applies. Active
or Frozen may terminate only through that retirement contract. While a
close is pending, the program keeps its core phase and sets
\texttt{close\_requested}; closing is not a phase of its own.

\paragraph{Fixed and unfixed plans.}
An unfixed \texttt{preparing\_plan} can change when $G$, the source duty or
candidate eligibility changes. The fixed handoff cannot. Its target pair,
assigned region, initial anchor, source context, selection position,
authority and reference bounds never change. In particular, a later OFF
does not replace a historical regular selection. Local locks, votes, view
entry times and timers belong to one pair only. After a crash, the
process resumes this durable record; it does not clear its votes or run
initial Enter again. Storage and task references carry the full pair and
anchor identity, so no other session identifier is needed.

Several regions may hold unfixed Waiting preparations for the same future
pair; only the region selected by the fixed handoff can enter. A
temporary loss of eligibility may pause work without terminating a
preparation that could become eligible again. A retired record stays as a
tombstone even after its resources are released.

\paragraph{Shared preparation survives a retry.}
The Session Manager keeps a shared preparation cache keyed by the
anchor's producing pair and hash, the native-state format, the
configuration and the region. Cancelling a program removes one consumer
of the cache, not already verified native state or useful unfinished
preparation for the same anchor. No consensus signature, timer or lock is
shared through this cache. The node-level public evidence store of
\cref{sec:source-retention} is likewise independent of its consumers:
branch-version changes, a growing $G$ and worker exit do not delete its
referenceable raw history, and do not move it to an archive that is only
eventually accessible.

It also keeps an authority cache $\tagg{AuthCache}[r,g]$ and a background
task that fetches, in order, the recovery authorities of the active
turn index. The task obtains and validates each next immutable recovery
authority, stores it, and advances its fetch cursor. A branch-version
change does not reset the cursor, drop an already learned future
authority, or cancel the task. Correct regional members relay complete
authority evidence. While the turn index stays fixed, every finite
required prefix of authorities is eventually cached. The cursor only
fetches evidence; it never enters a generation or substitutes a leader.
A new engine still starts from $g=t=0$ with its own admission and
recovery proofs.

An Active old program may already lack permission to extend $G$.
Conversely, the existing results of a Completed or Terminated program may
still be accepted, if their branch version and interval satisfy the
public acceptance rules. The phase neither grants nor revokes that
permission.

\subsection{Handoff records, anchors and Ready}
\label{sec:witness}
A handoff record, or \defterm{witness}, answers three questions that the
successor and every verifier must agree on: why the handoff happened,
who takes over, and which block the new session continues after. It is
not an extra round of consensus.

\paragraph{Choosing the next region: NextLive.}
$\tagg{NextLive}(i,\chi)$ searches the other eligible regions after $i$
in the fixed cyclic order $\mathcal L_R$ and returns the first one. It
returns $i$ itself only if no other region is eligible and $i$ still is.
Otherwise it returns $\bot$. A synthetic source $\bot$ searches from the
beginning of the order. Every call uses the designated target authority.

Eligibility means the effective L1 flag is ON and that authority gives
the region positive total weight. REGULAR handoffs and MSIG Forces both
use NextLive. A Force does not disable its source region. A physical
region may be selected again, but its new pair has fresh locks, votes and
timers.

\paragraph{The target index of a Force.}
\label{sec:retry-slot}
For an assigned export duty, let $(\beta_b,r_b)$ and $i_b$ be its
producing context, and let $a_b$ be its accepted DATA count so far. This
count is the number of the duty's DATA accepted into $G$; it is not a
regional height or a number of raw publications. The target index is
\[
 r'=
 \begin{cases}
 r_b,&a_b=0,\\
 r_b+1,&a_b>0.
 \end{cases}
\]
In words: if the export session has put nothing into $G$, the Force
retries the same authority index; otherwise it moves to the next one.
This rule chooses only the authority index. It does not decide whether a
Force is allowed, and the new region may be a different physical region
even when $a_b=0$. Both cases change the branch version to $\beta+1$ and
start their new engine at $g=t=0$. The retained anchor is exactly $G$ at
the start of the activating block, with its original producing pair,
height, hash and commitment evidence.

In the running example, $G=108$ and the next DATA 109 belongs to A's
$(7,20)$, which has already put 101--108 into $G$. So $r'=21$, even though
B's regular $(7,21)$ already exists and B has produced up to 115. The
export session is chosen by what $G$ needs next, not by the region that
runs furthest ahead.

An unassigned boundary stores both its source context and its target
index. A later Force changes the branch version and the cause but keeps
that target index; it does not recompute it as $r_b+1$. A regular
boundary stores source $(i_b,\beta_b,r_b)$ and target $r_b+1$. Genesis
stores source $(\bot,\bot,-1)$ and target 0. This handles repeated Forces
before any new duty is assigned, without creating empty sessions.

Regular completion opens $(\beta,r+1)$ from exactly $T_{\beta,r}$. A
zero-output retry may reuse the number $r$, even at another physical
region, but never rebinds $\tagg{AUTH}_{r,g}$. For a fixed region and
$(r,g,t)$, the recorded leader does not depend on $\beta$: the branch
version separates message and engine identities, not the schedule. The
fixed snapshot and quota stay those of the record. No old Quit, vote,
lock or timer becomes part of a new engine's state.

\paragraph{Historical selection.}
All nodes must search for the next region from the same historical
position; otherwise nodes that learn a handoff at different times could
pick different regions. For a REGULAR boundary, the search starts at
\[
 q_0=\max\{\operatorname{height}(\xi(T_{\beta,r})),
                         q^{\rm open}_{\beta,r}+1\}.
\]
For a Force, it starts at the triggering Espresso height. Select the
first verified historical position at or above $q_0$ for which
$\tagg{NextLive}(i_b,\chi_E(q))$ returns a region. REGULAR and MSIG use
the same selector. The eligibility snapshots after each prelude are
immutable replay products. The selected height $q_s$ fixes the region and
the opening reference bound. Missing source or authority evidence causes
a wait; it never causes a candidate to be skipped or an empty set to be
assumed.

A REGULAR handoff does not need the terminal DATA to be globally
accepted. It needs only a suitable later finalized header, which can be
empty or carry unrelated records. If none exists yet, the current session
can still form and offer its terminal result using its known opening
header; publication then supplies a later block. The Ready gate never
waits for this future header.

If the block that contains the terminal result is the first possible
selection position, the terminal DATA can be accepted there while the
successor opens there. The successor's DATA cannot be in that same
EXPORT, because its minimum reference would equal the containing block.
EXPORTs that span several sessions are still possible when the successor
was selected at an earlier finalized position. A speculative regular
handoff does not itself append DATA to $G$.

\paragraph{Handoff packet.}
A handoff packet contains the fields of \cref{tab:handoff-packet}.
References may name authenticated cached objects or objects still to be
fetched; a reference alone does not establish Ready.
\begin{table}[htbp]\centering\small
\caption{Handoff packet fields.}
\label{tab:handoff-packet}
\begin{tabularx}{\linewidth}{@{}>{\raggedright\arraybackslash}p{.26\linewidth}Y@{}}\toprule
Field & Type and meaning\\\midrule
Domain and kind & Chain/protocol domain; REGULAR or MSIG; genesis uses its trusted initial boundary.\\
Source context & Predecessor pair $(\beta_b,r_b)$ and source region
$i_b$, including the accepted count when assigned and the stored target
index when unassigned.\\
Target & Target pair $(\beta',r')$, the reason for its index, and the
selected region $i_{\rm next}$.\\
Anchor & Producing pair, DATA height, hash, full DATA or reference,
and original direct or covering commitment evidence.\\
Authority & Finalized $\tagg{AUTH}_{r',0}$, its fixed configuration,
quota, and publication-height proof.\\
Selection and bounds & $q_0,q_s$, authenticated eligibility,
$q^{\rm open}_{\beta',r'}$ and $\lambda^{\rm open}_{\beta',r'}$.\\
Cause and evidence & Mode-specific cause; trusted or validated replay
checkpoint, complete required L1/Espresso history and regional proofs.\\\bottomrule
\end{tabularx}
\end{table}

For the Force in the running example, the packet says: the cause is a
valid MSIG effective at a given Espresso block; it closes version 7 and
opens version 8; the export session was A's $(7,20)$, owing DATA 109; the
target is B's $(8,21)$, selected by the historical rule; the anchor is
DATA 108 with its exact hash, data and receipt, still produced by
$(7,20)$; and the authority is the finalized $\tagg{AUTH}_{21,0}$. A
regular packet has the same shape, but its cause is the terminal receipt
of the previous session, and its anchor is that terminal DATA.

Two checks use the packet. \tagg{VerifyHandoffHistory} asks whether the
handoff is valid \emph{as history}. \tagg{CanEnter} asks whether a node
may enter \emph{now}. For example, after a Force closes version 7, DATA
108 of $(7,20)$ is still valid history and can anchor $(8,21)$, but that
does not let anyone re-enter $(7,20)$.

\paragraph{Handoff verification.}
\label{sec:verify-handoff}
The pure procedure $\tagg{VerifyHandoffHistory}(C,\mathcal H,P)$ takes an
authorized replay checkpoint $C$, the required source history $\mathcal H$
and a packet $P$. It returns the historical fixed fields, \emph{invalid},
or \emph{needs evidence}. It authenticates what the boundary established.
It does not require the branch version to be open now, and it grants no
entry permission. Neither procedure changes public replay.
\begin{enumerate}[label=(\roman*),itemsep=2pt]
\item Check the domain, field types, and source and target names. Replay
authenticated compatible history from $C$ through all positions needed by
the cause, the historical selection and the incoming boundaries. Verify
the complete applicable controls, not only the named control. Consult
validated caches when available. A historical receipt from a producing
session may use \tagg{CheckFINAL} for intrinsic input validity; this check
does not fetch that producer's optional recovery-admission path. It still
authenticates each fixed initial snapshot and the exact DATA ancestry. The
target's executable state and its own regional admission are verified
separately.
\item For REGULAR, validate the predecessor's fixed boundary, initial
authority, exact quota and terminal DATA
$h(T_{\beta_b,r_b})=h(A_{\beta_b,r_b,0,0})+\TB_{r_b}$. Validate its terminal
FINAL kernel, actual ancestry and signed $\xi$, including the separate
historical source and boundary checks. Require $\beta'=\beta_b$ and
$r'=r_b+1$. The anchor must equal the terminal DATA. Validate the
predecessor's own boundary recursively in the same historical mode;
current closure does not erase its original evidence.
\item For MSIG, verify its authorized finalized L1 record, control
sequence and first-crossing header $J_{q(\theta_a)}$. Replay the full
prelude and check that this is the first due instruction whose target
matches the branch version at the start of the block. A stale or already
consumed instruction cannot form another boundary. No proof of missing
blocks is required. Take the captured source context, the old and new
branch versions, the retained prefix and the boundary from this unique
transition. Require $\beta'=\beta_{\rm entry}+1$, the target index of
\cref{sec:retry-slot}, and anchor equality with $G$ at the start of the
block, including original provenance. Verify the zero or positive
accepted count directly; do not infer it from the anchor's producer. For
an unassigned boundary, verify its stored target index and any Forces in
between. A later Force removes entry permission without invalidating the
historical evidence of the earlier boundary.
\item Recompute the first eligible $q_s$ from $q_0$ and all required
historical snapshots, using NextLive under the designated target
authority. Check the target region and the immutable target authority.
Set $q^{\rm open}_{\beta',r'}=q_s$, and set $\lambda^{\rm open}_{\beta',r'}$
to the maximum of the anchor's L1 reference height and the initial
authority's publication height. Validate the anchor's original direct or
covering commitment evidence under \cref{sec:coverage,sec:receipt-kernel},
together with its complete block and required ancestry. Do not relabel
it with the target pair.
\item Compare the derived fixed fields with $P$ and with any existing
fixed record for its target pair. Equal duplicates share that record;
different fixed fields for the same pair are rejected. A packet cannot
replace an installed handoff. This step records historical facts, not a
claim that the target may still enter.
\end{enumerate}

\paragraph{Permission to enter now.}
\label{sec:can-enter}
$\tagg{CanEnter}(S,P)$ takes the Session Manager's current verified
public state $S$ and a packet whose historical check has passed. It
requires that:
\begin{itemize}
\item the target branch version is current and open, and the Session
Manager holds no authenticated $\tagg{Closed}[\beta]$ for it;
\item the packet is the fixed boundary that still applies, and this
replica's region is the selected one;
\item the local record has no earlier initial entry, close request or
Terminated record.
\end{itemize}
A Closed record blocks entry even before a complete per-session $c_v$ can
be set; context resolution and the closure relay still run before the
refusal. Ready also checks the exact anchor state and the dependencies
below. A duplicate packet for an already entered record is serviced
without another Enter.

All initial entry paths use both checks. Receipt, ancestry and
retained-anchor validation use historical validation only, and global
EXPORT acceptance separately checks current append permission. A cache
miss goes through the same historical checks as a cache hit; no earlier
cached verdict is assumed.

A Force without an eligible successor still has a fixed cause, anchor,
source context and new branch version. The interpreter keeps that
pending boundary. It has no complete selected packet, and no program can
become Ready until the first eligible historical position fixes the
target region. A later Force can supersede it while keeping the same
stored predecessor context. The synthetic genesis boundary is verified
against the initial checkpoint and the selection rule of \cref{sec:l1};
it claims no REGULAR or MSIG cause.

\begin{definition}[\defterm{Ready}]
\label{def:ready}
A local Waiting program for $(\beta,r)$ is Ready exactly when:
\begin{enumerate}[label=(\roman*)]
\item historical handoff verification fixes this target pair and selects
the replica's region, or the authenticated genesis rule does so; and
the separate current-entry check $\tagg{CanEnter}$ also passes;
\item its exact anchor's original producing pair, height, hash, complete
DATA and commitment evidence have been validated;
\item the native execution state needed to continue after this anchor
has been recovered and verified against that same anchor;
\item the matching finalized authority, fixed quota, separate reference
lower bounds, the opening Espresso header and all validation
dependencies are available; a header after the opening is not required;
\item the Session Manager's processed public history has not closed or
superseded the handoff; and
\item the durable record has never performed initial Enter and has no
Terminated tombstone.
\end{enumerate}
\end{definition}
In the running example, the new B must: confirm that it is selected and
version 8 is open; hold the exact DATA 108, with its producer $(7,20)$,
hash and receipt; have the execution state after 108 (state after 110, or
after another branch, does not qualify); hold $\tagg{AUTH}_{21,0}$, the
quota and the opening references; and confirm that no later Force has
replaced the handoff and that it has never entered $(8,21)$ before.

State recovery uses the native validation interface. Cached state for
DATA 110 cannot make a program anchored at 112 Ready. An announced but
not yet effective Force permits only Waiting and preparation.

\paragraph{How long preparation takes.}
The preparation bound concerns one question: once public replay has
assigned the next export session to a region, how long until that region
is ready?

\begin{assumption}[Bounded preparation for a selected live region]
\label{ass:ready-service}
Let $J_o$ be the canonical block in which replay opens an assigned
export session, and let $t_0=t_E(o)$ be its finalization time.
Suppose the selected region stays operational under its fixed snapshot:
its correct members stay online, the regional complete-message contract
holds, and the declared input, load and state-service conditions hold.
Suppose also that the selected handoff stays applicable throughout
preparation. Then by $t_0+T_{\rm ready}$, each correct target member has
either already performed initial Enter, or obtained and validated the
exact anchor DATA, commitment evidence, matching executable state,
finalized authority and public context needed for its guarded initial
Ready/Enter.
\end{assumption}
The budget includes observing the opening block, retrieval, validation
and state installation; its clock does not wait until authority or data
has arrived. A regular successor that already entered through direct
delivery meets the obligation earlier. A Force that supersedes the
handoff ends the obligation for that target; not every Force leaves time
for a full preparation interval.

The bound applies to the declared, continuously served configuration and
state interface, not to arbitrary offline catch-up. It does not, by
itself, bound entry skew among members, require new output, or restrict
the next Force. $T_{\rm ready}$ is measured in real time, while controls
use Espresso timestamps and the advance-announcement rule, so no new
two-sided Espresso clock bound is involved. Where $t_0$ lies for each
kind of handoff is stated after \cref{alg:bootstrap}.

\paragraph{A complete regional ReadyBundle.}
\label{sec:ready-bundle}
The members of a selected region pass the entry material among themselves
in a ReadyBundle. It makes the transferred state explicit. It is not a
READY certificate and does not change the public opening of an export
duty.

For target session $(\beta,r)$, let $A=A_{\beta,r,0,0}$ and
$\Delta=\Delta_{i_{\beta,r}}$. A \tagg{READYBUNDLE} is an authenticated
regional message with the typed fields of the table below. Every required
object is present as bytes, once, in its normalized dependency set. A
URL, a state root, an unresolved hash or a promise to fetch later is not
a substitute.
\begin{table}[htbp]\centering\small
\begin{tabularx}{\linewidth}{@{}>{\raggedright\arraybackslash}p{.25\linewidth}Y@{}}\toprule
Field & Required value\\\midrule
Domain and target & Chain and protocol domain, encoding/native-state
format identifier, target $(\beta,r)$, selected region and sender identity.\\
Fixed handoff & The complete handoff of \cref{sec:witness}, including
its cause, predecessor/source context, selection and fixed target.\\
Exact anchor & $A$'s original producing pair, height, hash, full DATA,
original direct or covering commitment evidence and required ancestry.\\
Executable state $S_A$ & The full, already constructed application-state
object from which execution after $A$ can continue. It includes all state,
code and continuation metadata required by the native execution interface;
if partial execution is supported, its exact continuation position is
included. It carries no regional locks, votes or entry timers.\\
State evidence $\Pi_A$ & Native validation evidence connecting this exact
$S_A$ to this exact $A$, with all proof inputs and dependencies, under the
sound application-validation interface. A sender's assertion or a bare
state hash is not enough.\\
Authority and bounds & Finalized $\tagg{AUTH}_{r,0}$, fixed configuration
and quota, $Q^0_{\beta,r}$, $q^{\rm open}_{\beta,r}$ and
$\lambda^{\rm open}_{\beta,r}$, and all authority evidence needed to
validate the inherited history.\\
Public evidence & An authorized replay checkpoint, with its trust and
authorization dependencies, and the complete finalized L1
control/authority and Espresso header/namespace histories needed to
verify the handoff, references and authorization context. It contains the
finalized opening header and its evidence, but need not contain a later
header: the successor's selection floor advances only after the current
terminal result forms.\\
Dependency set & Full bytes for the transitive validation dependencies
of the fields above. A receiver need not fetch an object, replay
execution to rebuild $S_A$, or obtain a new authority before processing
the bundle.\\\bottomrule
\end{tabularx}
\end{table}

The native interface determines how $\Pi_A$ verifies the state at $A$.
If the block format does not commit to a state root directly, the bundle
supplies the execution evidence and authenticated starting state that the
interface requires; this document does not invent such a commitment. A
serialized snapshot that still needs restoration or conversion is not yet
the directly executable state required here. Recipients validate the
packet themselves; they do not trust the sender to authorize a handoff or
vouch for state.

\paragraph{Sending and receiving a ReadyBundle.}
Each correct selected member builds its complete bundle before its first
Ready/Enter event for the pair. Gathering histories, fetching across
regions, rebuilding state and forming the bundle all happen before that
event, and the model gives them no fixed time bound. At the event, the
member installs its verified state, executes initial Enter once, and
sends the complete bundle to every member of the fixed selected-region
snapshot. Since every correct member does this at its own first event,
the first correct entrant sends without needing to know that it is first.
Building and dispatching the bundle are not postponed past the recorded
entry time.

The receiver processes the bundle through its serialized Session Manager:
\begin{enumerate}[label=(\roman*),itemsep=2pt]
\item Authenticate the domain, sender and target, and inspect supplied
evidence under the evidence-first rule. Validate the historical handoff,
the exact anchor, all required authorities, references and native state
evidence. Incomplete or invalid material does not authorize entry. A
compatible existing cache can avoid duplicate validation, but the sender
never relies on an unknown receiver cache to omit required bytes.
\item Apply $\tagg{CanEnter}$ against the receiver's processed public
state. A closed or superseded target, a conflicting fixed handoff, a
close request or an exit record prevents initial Enter. The packet cannot
replace a fixed anchor or override a newer verified closure. If the
target is closed, run the Session Manager's phase-independent closure
relay using the packet's verified historical snapshot before refusing
entry. This also applies at a Terminated record and after a restart; the
worker need not be recreated. If the same pair has already entered, keep
any useful evidence without resetting its state or entry time.
\item For an unentered valid target, validate and adopt the received
$S_A$, check every remaining condition of \cref{def:ready}, and then
execute the Ready/Enter step of \cref{alg:bootstrap} in this same
message-processing event. The durable entry flag, the attached state and
the entry time are recorded atomically before signatures are dispatched.
Relay the complete bundle as above. No Ready queue, state-fetch task or
post-Ready installation job lies between this event and Enter.
\end{enumerate}

To \emph{adopt} the state means that the validated received state object
becomes the new program's executable application state. The program works
on that object; it does not first rebuild a different database from it.
This is a contract on representation and on the message handler, not a
claim that database import, disk copying or runtime initialization take
zero time. No receiver can pass Ready while installation work remains.

\paragraph{The timing consequence.}
Let $e$ be the first correct initial Enter time for this pair. The send
rule places a complete correct-to-correct bundle send at $e$. Suppose the
fixed handoff stays activatable throughout and no closure or other valid
entry prohibition intervenes. If validation, adoption and the atomic
Enter transition together form the message-processing step covered by
\cref{ass:network}, then
\[
 e_v(\beta,r,0,0)\le e+\Delta
 \qquad\text{for every correct selected member }v.
\]
Indeed, the first correct sender supplies the full executable object and
all validation dependencies. By the delivery-and-processing bound, each
correct recipient finishes that handler by $e+\Delta$. Its last guarded
step is Enter, unless it has entered earlier. No further
state-restoration premise is needed.

This conclusion depends on the abstract-message reading above. The
network premise covers delivery and processing of complete validation
dependencies; it does not bound an external native import or storage
operation merely because the protocol calls it ``installation''. If such
work remains after bounded message validation, let $u_v\ge0$ be its
actual residual delay. Then only
\[
 e_v(\beta,r,0,0)\le e+\Delta+u_v
\]
holds, when entry eventually completes and stays authorized; the network
assumption gives no uniform bound on $u_v$, and calling that work one
event does not make $u_v=0$. In such an implementation, the $e+\Delta$
claim and the bootstrap bounds derived from it do not hold. A concrete
native adapter must show that it meets this contract; sound state proofs
alone do not establish a timing property. \Cref{ass:ready-service} bounds
total preparation from public assignment, but it does not prove $u_v=0$
or replace this handler obligation. The network bound also does not
bound construction of the first complete packet or historical catch-up.

A ReadyBundle only transports existing state and evidence within a
region. It adds no quorum of READY signatures, no Espresso transaction,
no publisher guarantee and no public activation certificate, and it does
not execute a Force. Public assignment and the current-entry check remain
those of \cref{alg:force,alg:external,sec:ready}.

\subsection{Activation and bootstrap}
\label{sec:ready}
\begin{protocol}[H]{alg:bootstrap}{Enter and bootstrap the selected session}
Replica $v\in V_{i_{\beta,r}}$ runs these steps for one fixed target
$(\beta,r)$. Let $\Delta=\Delta_{i_{\beta,r}}$.
\begin{enumerate}[label=\textbf{\arabic*)},itemsep=3pt]
\item \textbf{Enter.} In the Ready event, atomically recheck
$\tagg{CanEnter}$ and the Closed map, then check and set the durable
initial-entry flag. Install $A_{\beta,r,0,0}$, $\tagg{AUTH}_{r,0}$,
$\TB_r$ and both opening reference bounds. Set the entry lock to
$Q^0_{\beta,r}$. Enter $(\beta,r,0,0)$ unfrozen, change Ready to Active,
record $e_v(\beta,r,0,0)$ once, and start the progress checks of
\cref{alg:recovery}. Send STATUS with this lock and the fixed handoff to
the initial leader, and process buffered proofs for this session.
Persist entry before dispatching signatures.
\item \textbf{Declare the fixed anchor.} The initial leader waits
$2\Delta$ from its own entry. If the view is still current and unfrozen
and the retirement contract permits the action, it sends NEWVIEW choosing
exactly $A_{\beta,r,0,0}$, justified by $Q^0_{\beta,r}$ and the fixed
handoff. The leader does not choose the anchor; the handoff has fixed it.
\item \textbf{Certify and produce.} Members run the ANCHOR-vote step of
\cref{alg:recovery}. On the ANCHOR QC, the leader starts
\cref{alg:steady} at DATA height $h(A_{\beta,r,0,0})+1$ and physical
position 1. Local failures use the same recovery rules. All these steps
occur in Active, unless a Quit, Completion or retirement transition
intervenes.
\end{enumerate}
\end{protocol}

For example, with anchor 108, the first DATA is 109. ANCHOR voting itself
starts no commit timer.

Let $e=\min_v e_v(\beta,r,0,0)$ over correct members. Under the
complete-handler timing contract above, if the handoff stays activatable
and bootstrap is not interrupted, all correct members enter by $e+\Delta$.
With a correct initial leader and no Quit:
\begin{center}\small
\begin{tabular}{@{}ll@{}}\toprule
Event & Bound\\\midrule
all correct members have entered & $e+\Delta$\\
NEWVIEW is sent & $e+3\Delta$\\
its ANCHOR QC reaches the leader & $e+5\Delta$\\
the first DATA reaches all correct members & $e+6\Delta$\\\bottomrule
\end{tabular}
\end{center}
These bounds start at the first actual Enter, not at public assignment or
candidate creation. The last one says the DATA has arrived, not that it
is committed or globally accepted. If native installation is not part of
bounded message processing, these bounds have not been established; the
extra delay is not silently absorbed into $\Delta$. \Cref{ass:ready-service}
alone does not prove this entry-skew bound.

\paragraph{Public assignment and the preparation clock.}
Public replay assigns the next export session independently of local
program phases. The \emph{public opening position} of that assignment is:
\begin{itemize}
\item for a normal successor that is selectable at once, the position of
the EXPORT that delivers the terminal DATA;
\item for a forced successor that is selected at once, the triggering
block;
\item if selection found no eligible region, the first later position at
which one is eligible.
\end{itemize}
A late proof of an earlier selection does not move this historical
position to the local retrieval time.

The real finalization time of the public opening block is the origin
$t_0$ of \cref{ass:ready-service}. A direct regular successor may already
be Active before its export assignment opens. An unassigned boundary has
no selected target, so no preparation obligation applies to it. There is
no public delivery deadline: download completion, Ready, restart and
silence never generate a cross-region Force.

\subsection{Concurrent candidates and races}
\label{sec:runtime-races}
Suppose A's session is $(7,20)$ with terminal DATA 110, B already produces
in $(7,21)$, and $G$ ends at 108. A future MSIG targets branch version 7.
B can keep its old Active program while preparing candidate $(8,21)$ from
108. If $G$ advances to 109, that preparation can switch to 109; the old
program's anchor stays 110. If the Force then fixes B as successor with
anchor 109, all further preparation is for that fixed anchor. B can enter
$(8,21,0,0)$ while $(7,21)$ closes under \cref{sec:retirement}. Their
locks, votes, entry times and timers are distinct.

The candidate can also change region and index, as
\cref{fig:runtime-candidates} shows:
\begin{center}\small
\begin{tabularx}{\linewidth}{@{}>{\raggedright\arraybackslash}p{.33\linewidth}>{\raggedright\arraybackslash}p{.33\linewidth}Y@{}}\toprule
New fact & Force now would open & Preparation\\\midrule
$G=108$; 109 still belongs to A & B's $(8,21)$ after 108 & B prefetches 108\\
$G=109$ & B's $(8,21)$ after 109 & update the new program; old B unchanged\\
A's 110 accepted; B has no DATA in $G$ & export is B's $(7,21)$ with zero output; next region after B is C, index 21 & C prefetches 110\\
OFF(C) takes effect & D's $(8,21)$ & pause unneeded work; D prepares\\
B's 111, 112 accepted & D's $(8,22)$ after 112 & D prepares state at 112 and authority 22\\
Force takes effect & the fixed handoff & check Ready only against it\\\bottomrule
\end{tabularx}
\end{center}
All rows except the last are what-if planning; they grant no right to
produce. In the row with zero B output, the export context is B's
$(7,21)$, so the candidate region moves past B while the target index is
still 21. Once B's DATA enters $G$, the positive-output rule moves the
index to 22. If the Force finally retains B's 112, the fixed target is D's
$(8,22)$ with the anchor still produced by $(7,21)$. D must verify the
state at 112 and authority 22 before producing DATA 113. Earlier
preparation of $(8,21)$, or of state at 109, grants no entry. OFF cannot
rewrite a regular handoff that is already fixed.

A terminal receipt received before the Force can prepare or start a
regular successor in the old branch version. Once the same node verifies
the Force, it cancels an invalid unentered target or requests retirement
of an entered one. If the Force arrives first, the late receipt is kept
as regional evidence and cannot revive the closed target. Nodes can
briefly hold different speculative programs, but the same canonical
replay subposition, with complete evidence, fixes the same public result
and Force boundary. Local caches and future control queues need not agree.

Every fetch callback keeps its original target and anchor identity. A
reply for $(7,21)$ at 110 cannot make $(8,21)$ at 108 Ready. An old
candidate reply may be cached and validated for another use, but it
cannot overwrite a fixed anchor. Duplicate receipts, packets and
callbacks cannot run initial Enter twice. A stale MSIG that follows an
earlier MSIG does not close the new branch version. Independent STATUS
records keep their own effects. Reconciliation drops an obsolete
preparation basis without invalidating a different fixed handoff.

\begin{figure}[p]\centering
\includegraphics[width=.98\linewidth,height=.87\textheight,keepaspectratio]{figures/cross_region_rotation_v4/candidate_ready.pdf}
\caption{A changing Force candidate while the old program continues.
The finalized boundary fixes the target and exact anchor; only then can
verified state at that anchor enable Ready and Enter.}
\label{fig:runtime-candidates}
\end{figure}
\FloatBarrier

\subsection{Graceful close and termination}
\label{sec:runtime-close}
Closing a branch version publicly and shutting down a local program are
separate events.
\begin{itemize}
\item \emph{Public closure is immediate.} Once a Force takes effect at its
Espresso block, \cref{alg:force} stops the closed branch version from
appending anything new to $G$.
\item \emph{Local shutdown takes time.} Each node, once it has verified
the closure, stops new proposals, votes, BLAME and term or recovery
changes for that branch version. It does not use the Force to cancel
running commit timers. Running timers, FINAL messages, receipts and
evidence are processed as before.
\end{itemize}
\Cref{alg:rotation} stores the verified closure and disseminates it for
every known session snapshot, in every program phase. Entered programs
also drain through \texttt{RequestClose}. The request is idempotent and
sets \texttt{close\_requested}; it is neither a Quit nor terminal
Completion. \Cref{sec:retirement} defines the exact limits on new work,
the treatment of running timers and messages, and when
\texttt{FinishClose} may run. Its drain threshold comes from the regional
delay bound; transferring resources has no time bound. A new session
need not wait for the old program to finish closing.

For an unentered program, a definitive invalidation cancels pending entry.
It may release preparation tasks once the durable closure, pending-relay
and late-message duties have been transferred. It has no regional timers
or votes to drain, but it must not stop the Session Manager's closure
dissemination. A Completed program has already stopped its engine in
Completion. For an Active or Frozen program, only the retirement contract
can authorize \texttt{FinishClose}. An empty local timer list, a
different current public branch version, or a locally observed terminal
block is not such an authorization.

In every case, releasing resources requires durable exit and
responsibility records. Committed DATA, receipts, regional evidence,
anchors and recoverable state must stay accessible under the storage
rules. Outstanding complete results and publication retries are handed
to the Publish Layer, and pending requests in discarded suffixes stay
available for reproposal. Storage and publication workers acknowledge
the duties they accept before the program releases its last required
copy. Later messages can update those services without reactivating the
engine. Termination does not undo regional commitments, remove DATA from
$G$, or authorize stake withdrawal.

\subsection{The retirement contract}
\label{sec:retirement}
This subsection gives the exact rules for shutting down the programs of a closed branch version. Their goal is that closure is never silent: every correct member of an old session learns of it within a bounded time, running timers finish, and no evidence is lost.

Fix a session $(\beta,r)$ and write $\Delta=\Delta_{i_{\beta,r}}$,
using its immutable voting snapshot. Closure is a duty of the Session Manager,
whether or not a program exists or has entered.
The Session Manager persistently stores authenticated $\tagg{Closed}_v[\beta]$
and its complete source evidence, not just a hash or retrieval promise.
The regional control owner keeps those bytes available for the receive
handler below; a pending storage copy does not suspend that owner.
This branch-wide record immediately
prohibits every old-$\beta$ initial entry and fresh core action.
It is separate from a program's six phases and from the later
per-session complete-close event defined below.

For each known fixed session snapshot in which $v$ is a member, the
manager retains a control record even in Waiting, Ready, Completed or
Terminated, or when no worker exists. The record includes the historical
fixed boundary, source region, immutable authority and membership,
\texttt{close\_requested}, any first complete-close time $c_v$, and
pending relays. Active/Frozen programs keep their core phase while
draining; only genuine Quit or terminal Completion changes that phase.
Force is neither a quit proof nor a terminal receipt.

\paragraph{Complete-receipt dissemination.}
\label{sec:retirement-receipt-relay}
The following obligation applies during normal running as well as
retirement. The first time a correct regional member or its local
manager/collector holds a complete valid receipt of its own session,
it multicasts the receipt and all validation dependencies to the
session's fixed membership. In particular, a terminal receipt is
multicast no later than its local Completion event. Deduplication by
session and DATA hash suppresses repeated encodings or signer sets,
not this first complete relay. The relay and any pending delivery are
preserved across Completed and Terminated by the persistent owner.
Recipients process an own-session terminal receipt under Completion
even while Frozen or closing, independently of its generation and
before a simultaneous timer. A program that has not entered caches
it; a Terminated program archives it without restarting. This rule
forwards existing signatures and introduces neither a new vote nor a
certificate. The intra-region complete-message bound applies to the
actual relay; it is not a bound on first forming a receipt.

\paragraph{Closure handling and \texttt{RequestClose}.}
The following handler runs at the Session Manager or persistent control owner,
not only at an entered engine.
\begin{enumerate}[label=(\roman*),itemsep=2pt]
\item \emph{Record public closure.} Verify complete evidence that the
canonical prelude of an authenticated finalized Espresso block closed
$\beta$. Retain its triggering header, applicable complete control and
namespace history from a validated checkpoint, and closure and
retained-prefix context. Check the explicit MSIG rule; a pending
MSIG, OFF, unfinalized header or bare timestamp is insufficient.
Atomically persist $\tagg{Closed}_v[\beta]$, prohibit fresh actions for
that version and notify all its local control records, regardless of
phase. This consumes an established closure; it does not execute Force
or select an anchor a second time.
\item \emph{Disseminate for every known snapshot.} For each known fixed
session $(\beta,r)$ of this region, combine the retained closure proof
with its immutable authority, membership and historical admission
context. This is a \defterm{complete close bundle}: another member can
validate both closure and the exact recipient snapshot without a prior
program, live-entry permission or an unprovided cache. On first complete
processing, perform RequestClose atomically: set
\texttt{close\_requested}, record $c_v$ as the current time, and multicast
the bundle to the fixed membership. This includes Waiting, Ready,
Active, Frozen, Completed and Terminated records. A recipient validates
the bundle, persists Closed and performs the same idempotent operation;
no worker or engine entry is created. Repeated bundles keep the first $c_v$.
\item \emph{Keep unresolved dissemination pending.} If the fixed snapshot or
its dependencies are not yet known, retain the authenticated Closed
record and a pending relay obligation. The entry gate is already shut,
but $c_v$ is not assigned and no $4\Delta$ clock is invented.
Whenever any old-pair Ready, proposal, authority, history or other packet
provides the missing authenticated snapshot context, the manager or
archive completes and multicasts the close bundle in that receive
handler, and replies with it to the packet's sender. Do this before
ordinary BFT processing, rejecting a live-entry request, discarding a
duplicate, or returning a Terminated/tombstone result. Historical
context is validated without requiring that the old version remain
open. Complete proposals from correct voters carry or resolve this
context within the complete-message bound. This handler is
mandatory even if the old worker was never created, its ReadyBundle
is refused, or a delayed callback arrives after exit or restart.
\item \emph{Gate new work.} Closed prohibits every old-version
DATA/DRAIN proposal, NEWVIEW, DATA/DRAIN/ANCHOR vote, new commit timer,
progress check producing BLAME, term entry and generation entry.
The checks cover initial Enter, delayed authority/recovery callbacks,
terminal-QC rotation and buffered-message processing. Neither a new
STATUS nor recreation of a local worker bypasses the persistent gate.
Already dispatched messages are not withdrawn. New-session preparation
and entry use their own pair and need not wait for retirement.
\item \emph{Drain running work.} Keep running commit timers and
their deadlines; observing Closed alone does not cancel a
timer or release its evidence-processing owner. Continue complete
validation, comparison with retained signed material, and evidence-first
processing. A genuine current-generation quit proof still runs Quit,
broadcasts its proof, cancels its timers and preserves commitments;
it cannot reopen the gate. A surviving timer commits its block and
ancestors. For each DATA thereby committed, its FINAL from \cref{alg:steady} goes
both to the fixed regional membership and to the publication collector,
preserving \cref{def:receipt}. A partial FINAL does not invoke the
complete-result guarantee of \cref{ass:services}(i).
\item \emph{Keep evidence duties.} Continue collecting old votes, QCs,
FINAL messages and receipts. Multicast complete own-session receipts
as specified above, and offer every complete result to the Publish
Layer. A valid terminal receipt still runs genuine Completion,
possibly changing the phase to Completed and canceling timers; a
drain deadline that is already set still applies. Neither receipt arrival nor
later ON reopens the branch. Keep data, control-relay and publication
responsibilities until their durable transfer.
\end{enumerate}
All closure relays forward authenticated source evidence and
signatures that already exist, not new BFT votes or certificates. The manager/control owner
continues transporting them after worker termination. A known snapshot
is never forgotten merely because its program exits. Recovery from a
crash restores Closed, session close times, pending relays and owners
before dispatching callbacks or signatures. Messages, including complete
closure, quit and terminal-receipt evidence, are processed before a
simultaneous timer, entry or retirement event. No receipt, lock, vote,
committed block or provenance is erased by this handler.

\paragraph{\texttt{FinishClose}.}
For a program Active or Frozen when it first learns closure, first
require its complete per-session RequestClose event. Set
\[
 d_v=c_v+4\Delta.
\]
An earlier Closed flag without a complete-close event is not a substitute
for $c_v$ and never authorizes early release of timer/evidence processing.
FinishClose may mark it Terminated only at a time $t\ge d_v$ when:
\begin{enumerate}[label=(\roman*),itemsep=2pt]
\item every commit timer that was running has executed or was canceled by
genuine Quit or Completion, and no queued action can pass a fresh-action
gate;
\item persistent storage owns the exit, Closed and per-session close records,
committed log, required signed evidence and validation dependencies,
receipts, anchor data and recoverable state, without releasing the last
required copy;
\item a persistent owner has accepted pending closure-context discovery,
complete-close relays and replies, complete-receipt multicasts,
complete-result offers, publication retries and data-service duties, and each locally complete result has been offered to the
Publish Layer; and
\item the dispatcher routes later messages for $(\beta,r)$ to those
owners. They must run the manager-level closure handler even for an
unrecognized old pair, and may collect or relay existing signatures,
validate late receipts, serve data and continue publication, but cannot restart the
engine, execute its BFT timers or issue fresh BFT signatures for it.
\end{enumerate}
The ownership transfer may be local to the same correct replica; it
requires no additional agreement or certificate. The receiving control
owner must already hold the required evidence and receive route before
the old owner releases them; an unresolved retrieval task is not such
a handoff. Only after these
guards hold may the production/recovery worker and its timer machinery
be released. A late terminal receipt records archival evidence and does
not change a Terminated program back to Completed or Active.

A definitively invalid Waiting or Ready program has no regional timers
to drain. It may release its worker without waiting $4\Delta$ only after
its durable exit record and all manager/control duties above have an
owner. In particular, learning Force before its first ReadyBundle never
permits silent cancellation: known snapshots trigger immediate
RequestClose; unknown snapshots leave pending relays that the refused
ReadyBundle or a later proposal must resolve. Already-Completed programs
have already canceled their timers in Completion and may release workers
after the same responsibility transfer. They likewise retain all
control-relay duties. A temporary candidate-eligibility change alone
need not terminate preparation. Termination never deletes a commitment,
removes DATA from $G$, or authorizes stake withdrawal.

\begin{proposition}[No silent closure in a vote window]
\label{prop:retirement-observable-close}
Let a correct member vote for and completely forward a valid proposal
of $(\beta,r)$ at time $s$. If a correct member's matching vote or entry
is suppressed by its authenticated Closed record during processing of
that proposal by $s+\Delta$, a correct member processes and multicasts
a complete close bundle for this session by $s+\Delta$. Thus, writing
$c=\min_{v\ {\rm correct}}c_v$ over all phase-independent complete-close
events (with $c=+\infty$ if none occurs), such suppression implies
$c\le s+\Delta$.
\end{proposition}
\begin{proof}
The forwarded proposal and its historical authority/admission dependencies
reach the member's manager or archive within $\Delta$, even if it has
no worker or an exit tombstone. These dependencies identify the fixed
regional snapshot. If Closed was already recorded, the receive handler
combines them with its retained closure evidence and sends the complete
bundle before refusing ordinary processing. If closure is learned after
that packet but by the time it suppresses a vote, the snapshot is already
known and closure's own handler performs the multicast immediately.
Either way some $c_v\le s+\Delta$. No new initial entry, executable-state
restoration or remote membership discovery is required for this control
relay. Restart and duplicate paths use the same persistent handler.
\end{proof}

\begin{proposition}[Closure propagation and honest-message drain]
\label{prop:retirement-drain}
Under \cref{ass:network}, let $c=\min_{v\ {\rm correct}}c_v$ be the first
complete per-session RequestClose time among this snapshot's correct
members in any phase, including its manager/archive owners. Then every correct member has closed by $c+\Delta$.
Every surviving honest commit timer expires by $c+3\Delta$, and the
regional FINAL messages generated by those timer events reach every
correct member by $c+4\Delta$. Consequently $d_v=c_v+4\Delta$ is a
conservative local drain threshold, not a deadline for finishing all
retirement duties.
\end{proposition}
\begin{proof}
The first correct relay includes closure and the fixed snapshot's full
validation context. Its delivery and processing are covered by the
regional bound even when no receiving worker exists. Thus
$c\le c_v\le c+\Delta$ and every old-version fresh-action gate is shut
by $c+\Delta$. A member may have set Closed still earlier; that can only
stop its votes earlier. A correct DATA/DRAIN vote at time
$s=\tau_v(P)$ starts its only timer, with expiry $s+2\Delta$.
The close gate gives $s\le c_v\le c+\Delta$; no callback starts a
replacement timer. Thus its expiry is at most $c+3\Delta$.
Quit and Completion can cancel a timer but cannot postpone it.
The FINAL messages actually generated at surviving expiries are
delivered regionally within another $\Delta$. Finally $c_v\ge c$
implies $d_v\ge c+4\Delta$. Message-first processing covers equality.
\end{proof}

The bound starts at complete per-session closure processing and its
actual regional relay. It does not start at an earlier Closed observation
whose session snapshot was unknown, a control signature, the triggering Espresso timestamp, or the
start of unbounded historical catch-up. It does not add a new source
visibility assumption. Three local $\Delta$ intervals cover the last
honest timer event but not its FINAL delivery: a member can vote just
before $c+\Delta$, commit just before $c+3\Delta$, and deliver FINAL
only just before $c+4\Delta$.

Four $\Delta$ intervals do not bound subsequent receipt relays,
Byzantine message release, a storage copy, $T_{\rm pub}$, $T_E$, or
successful EXPORT acceptance. Their outstanding duties pass to the
persistent owners; no bounded storage size or transfer latency is
claimed. Elapsed grace alone therefore never authorizes FinishClose.

\begin{proposition}[Conditional regional safety during retirement]
\label{prop:retirement-safety}
Assume the stated regional model, the regional validation
and DATA/DRAIN soundness premises, and the complete-message Ready/Enter
installation contract of \cref{sec:ready}. With the phase-independent manager closure handler,
complete-receipt multicast and retirement rules above, correct regional committed logs
remain prefix-comparable and irrevocable. The external prefix argument
of \cref{sec:safety-argument} is preserved. This does not assert a
protecting QC for every late surviving timer, old-session progress,
or complete-receipt formation.
\end{proposition}
\begin{proof}
We verify the prerequisites of the timer/lock argument explicitly;
Part~I's safety theorem is not used to infer an unproved universal-QC property.
Fix a correct direct timer for DATA or DRAIN block $B$ in generation
$g$, voted at $s$, that actually expires at $s+2\Delta$ without
cancellation. All comparisons below concern this same session.
Let $e_0$ be its first correct initial-entry time. The specified Ready
contract enters every nonclosed correct member by $e_0+\Delta$.
An own-session terminal receipt requires a correct FINAL after a timer,
whose earliest possible expiry is $e_0+2\Delta$. Hence such a receipt
cannot evade Completion by arriving only at an unentered cache of a
correct member that will run the chosen timer.

\emph{Quit and Completion cannot suppress its early vote window.}
If any correct member first quits $g$ at $a\le s+\Delta$, its complete
quit proof reaches the committer by $a+\Delta\le s+2\Delta$ and
cancels the timer. Thus any first correct Quit satisfies
$a>s+\Delta$. If any correct member Completes at $b\le s+\Delta$,
its complete terminal receipt is multicast by $b$ and reaches the
committer by $b+\Delta\le s+2\Delta$. The committer then executes
Completion: it commits the terminal DATA and ancestors and cancels
this timer, including a DRAIN timer. Again this contradicts the chosen
surviving direct timer. No correct Completion can therefore occur at
or before $s+\Delta$. This is an operational cancellation argument;
it does not assume that receipt-triggered commitments are already
known to be prefix-comparable.

\emph{The proposal is processable within its regional window.}
For $g=0$, the self-contained Ready relay and installation contract
in \cref{sec:ready} supplies entry by $s+\Delta$; this assertion is
conditional on that precise contract, not a bound on state restoration
before any member is Ready. For $g>0$, let $a^-$ be the first correct
Quit of $g-1$, and let $e$ be the first correct entry into $g$.
The recovery rule gives $e\ge a^-+\Delta$ and $s\ge e$.
The complete quit relay reaches every correct member by $a^-+\Delta$;
its own mandatory wait has therefore ended by
$a^-+2\Delta\le e+\Delta\le s+\Delta$. The forwarded complete
proposal supplies the finalized authority and remaining admission
history. Hence a receiver with no Quit, Completion or close preventing
entry can process it in $g$ by $s+\Delta$. Applying the same
complete-proof argument inductively to earlier generations supplies
any predecessor entry needed to process the quit proof. This uses
the recovery waits and complete validation bundles, not a
new simultaneous authority-visibility time.

\emph{Early and late windows.}
Let $c$ be the first complete per-session RequestClose event among
correct manager/control owners in any phase. A prior Closed flag or
unentered/Terminated record is not excluded merely by definition of $c$.
Instead, \cref{prop:retirement-observable-close} shows that if it
suppressed a matching vote or entry by $s+\Delta$, then $c\le s+\Delta$.
Consequently, if $s+\Delta<c$, no closure suppresses that vote or entry.
The preceding cancellation arguments also exclude Quit and Completion
through that window. Earlier-term processing and prior indirect
commitment do not suppress that vote either. An incompatible retained
history would expose a genuine quit proof in this window and contradict
timer survival. Every correct member therefore has cast, or now casts,
the matching DATA/DRAIN vote by $s+\Delta$. Their total weight is at least
$W-F>F$, so the votes form $B$'s protecting QC by $s+2\Delta$ at
all members that can still take a later recovery entry. A member that
Completed after the vote window cannot take such an entry; it need
not be assigned an unused recovery lock.

Any correct entry into $g+1$ must obey
\[
 e_v(\beta,r,g+1,0)\ge a_v+\Delta
 \ge a+\Delta>s+2\Delta.
\]
In the early case, every such entrant therefore knows the protecting
QC before its lock snapshot. The highest-QC rule gives the
recovery-lock invariant of Part~I. Closure preserves collecting those already
sent votes and QCs, and may only remove a later entry. In the late
case $s+\Delta\ge c$, the timer expires at or after $c+\Delta$.
The complete snapshot-bearing relay has closed every correct member's
entry gate by then, including previously unentered and archived members; the strict displayed
inequality prohibits every later-generation entry and thus its
correct signatures. A later-generation QC needs a correct signer,
so no recovery can exploit a missing protecting QC in this case.

\emph{Direct timers and receipt-based commits.}
In one generation, complete signed histories either extend one
another or expose an admissible equivocation proof.
A correct voter forwards its complete history. A correct signer of
an incompatible history either processes $B$ first and cannot vote,
or forwards its history before processing $B$, no later than
$s+\Delta$. That history reaches the committer by expiry and exposes
the proof. Genuine Quit, or Completion before expiry, would cancel
the timer. Thus no conflicting same-generation correct direct timer
can survive; the same argument excludes a conflicting QC, including
one assembled later from old votes.

Across generations, a surviving lower-generation direct timer cannot
be in the late case if a correct higher-generation vote exists.
Its early-case QC is known before every later live lock snapshot.
The ranked-QC descent of Part~I now has its actual prerequisites:
correct ANCHOR voters meet those locks; DATA/DRAIN QCs inherit a
lower-ranked parent justification. It excludes a later QC that does
not extend the committed block, and the parent of a later direct
proposal consequently extends it. These are precisely the
protecting-QC, recovery-lock and unique-extension steps of Part~I,
with the explicit vote-window checks above and the
DATA/DRAIN validation rules. Direct commitments are comparable, and
ancestor commitments inherit that property.

Finally a valid receipt contains a correct FINAL signer. That signer
sent FINAL only after a timer committed its named DATA, directly or
as an ancestor; committing from another receipt emits no FINAL.
Consequently every receipt-based commit, including Completion, traces
to a causally earlier direct timer commitment of that DATA or a
descendant. It adds an already protected prefix. This provenance
argument does not take the safety of receipt-based commits as its own
premise. A timer canceled by Completion contributes no additional
direct commitment, particularly no commitment of its unexpired DRAIN.
No retirement step removes a prior commitment.

The public interpreter still preserves all of $G$ at its canonical
Force boundary and rejects further old-version DATA. Retirement and
receipt relay neither select a retained anchor nor authorize an
append, so they do not affect those external transitions.
\end{proof}

The universal protecting-QC lemma from Part~I cannot simply be copied:
its proof requires matching votes from all correct members, whereas a
closed member keeps a proposal without voting. For example, with
three unit-weight members and $F=1$, let correct $H_0$ close at $c=0$
and correct $H_1$ receive closure at $\Delta$. A Byzantine leader gives
a valid proposal to $H_1$ at $0<s<\Delta$. $H_1$ votes, forwards and
later commits; $H_0$ keeps it without voting. If the Byzantine member
withholds its vote and FINAL, only $H_1$'s FINAL exists. Its weight is
not greater than $F$: neither a proposal QC nor a complete receipt need
ever form. Retirement preserves this local commitment and evidence
without calling it Completion. The Publish Layer's complete-result
premise does not repair the missing signature.

\paragraph{Force before the first ReadyBundle.}
Suppose $H_0$ enters at $e_0$ and relays its complete ReadyBundle, while
$H_1$ learns Force before that bundle arrives. $H_1$ immediately stores
Closed. If its snapshot is already known it multicasts closure at once;
otherwise the ReadyBundle supplies the historical snapshot by
$e_0+\Delta$ and triggers the same multicast before its entry refusal.
The complete close evidence therefore reaches $H_0$ by $e_0+2\Delta$,
regardless of $H_0$'s later direct Espresso retrieval. Control processing
wins equality with its initial leader wait. If a proposal is the first
packet to reveal the snapshot instead, \cref{prop:retirement-observable-close}
gives the required $c\le s+\Delta$ and the late-window argument blocks
later recovery. Worker termination, a stale callback and restart cannot
remove this relay path. This rules out a silently absent member being
used to classify an unprotected timer as an early-window commitment.
\subsection{Node storage and source relay}
\label{sec:source-step}
\label{sec:source-retention}
A node must keep the exact bytes of old public objects. The object set
$\mathcal O_q$ of \cref{alg:global-rule} only grows, so the rule may
append a DATA whose body or FINAL signatures arrived many blocks
earlier; a handoff or recovery may also refer to old objects. Keeping
only a hash is not enough. This subsection states what each regional member stores, how it starts storing, and how it relays new source data to its region.

\paragraph{The public store.}
Each correct regional member keeps a durable \tagg{PublicStore} shared
by its source reader, interpreter, manager, engines and publication
workers. It stores content, not only cursors and indexes. The store
has the following distinct components.
\begin{table}[htbp]\centering\small
\begin{tabularx}{\linewidth}{@{}>{\raggedright\arraybackslash}p{.24\linewidth}Y@{}}\toprule
Component & Stored content\\\midrule
Raw Espresso source & Every acquired complete relevant namespace,
including rejected or malformed records; headers, finality/linkage
proofs, authenticated absence, and canonical source order.\\
Object set $\mathcal O_q$ & The block bodies and FINAL signatures
published through $J_q$, with their bytes, indexed by hash and session,
and the source block that carried each. An index without the bytes is
not a stored object.\\
L1 source objects & Every obtained finalized authority, control,
header and registry/completeness proof, with its authentication inputs;
pending controls as well as applied controls.\\
Regional witnesses & Full validation bundles actually used by correct
regional actions, committed DATA and receipt dependencies. Possession
does not put a private object into $\mathcal O_q$.\\
Prepared anchors & Exact executable continuation objects and their
validated provenance, separately keyed by anchor, configuration,
region and native format.\\
Replay state & Verified checkpoint, the current output $\mathrm{Out}$ of
\cref{alg:global-rule} with its pending boundary, the replay subposition,
source-availability cursors, and typed pending fetches.
The source-availability cursor is distinct from the interpreter cursor.\\\bottomrule
\end{tabularx}
\end{table}

At each source-availability event, store the entire authenticated raw
namespace and its evidence before advancing the availability cursor,
releasing the input buffer, or making a relay depend on it. An object
that \cref{alg:global-rule} has not used, or never uses, is stored all
the same. Likewise store each acquired L1
object before advancing its authority-fetch cursor. Regional validation
stores the witness it used, so a later required complete send does not
reconstruct that witness from an archive. Correct FINAL and complete
receipt sends obey the full-bundle rule of \cref{sec:admission}; a
reference is used only when its receiver can resolve it within the
regional delay bound.

A new branch version, session closure, advancement of $G$, checkpoint
compaction and worker exit do not delete these objects. Cancellation
removes a consumer, not the store's retention responsibility. Transfer
outstanding sends and ownership durably before releasing resources.
Restart restores the exact store and replay subposition; loss of required
bytes requires hydration before the prepared timing claim can resume.
These are protocol storage rules, not an additional storage service.

\paragraph{Finite initial hydration and pruning.}
\label{sec:source-hydration}
Choose a fixed authenticated source cut $q_0$ when beginning this
retention discipline. Hydrate the raw source interval still legally
referenceable after that cut and all exact objects designated by the
starting checkpoint. Because $\mathcal O_q$ is cumulative, this
includes relevant raw namespaces back to the authorized source origin,
including objects that were never appended. Storing only accepted
ancestry is insufficient. Also prepare the fixed anchor's provenance and the
authority and state dependencies of the assigned handoff.
Store every newly available namespace beyond $q_0$ while hydration
runs. Neither a new branch version nor a later source height restarts the fixed
old task or changes its target to a moving cut.

The interval and each admitted payload are finite. Eventual historical
retrieval gives a finite completion time for that fixed set; the anchor
and finite AUTH prefix have their separate eventual-availability
premises. Taking a finite maximum imposes no bound on arbitrary archive catch-up.
The selected-region preparation service is stated separately in
\cref{ass:ready-service}. It does not include all future
L1 objects that an arbitrary sidecar might name.

A replay checkpoint may compact computation without deleting
referenceable bytes. Deduplication and equivalent lossless storage are
allowed while exact content and source-occurrence mappings remain
available to the modeled handler. A hash or an only-eventually-readable
archive is not a replacement for those bytes. Pruning requires a
protocol proof that no future supported reference, checkpoint,
anchor, pending validation or publication duty can need the object.
Advancing $G$ or terminating its producing engine does not prove this.
\Cref{alg:global-rule} itself never appends an object of a closed
version, or a DATA at or below the height of $\operatorname{last}(G)$
that is not in $G$; other duties may still need such objects, so this
alone does not justify pruning them.
The present rule supplies no finite-storage bound.

\paragraph{What the rule reads from the store.}
\label{sec:record-completeness}
\Cref{alg:global-rule} reads only $\mathcal O_q$ and the L1 inputs of
\cref{sec:logical-carrier}. $\mathcal O_q$ consists of the block bodies
and FINAL signatures carried by the authenticated namespace content of
$J_1,\ldots,J_q$, together with the objects that a protocol-authorized
checkpoint establishes. A private cache, a newer private checkpoint, a
hash or name without its object, or an unresolved sender assertion adds
nothing to $\mathcal O_q$. L1 objects keep their own authentication
interface and need not first occur in Espresso.

For a candidate $B$ in Step~3, the rule needs exactly the body of $B$,
the bodies of a covering path to a root $D$, the FINAL signatures on
$D$, and the session record $\sigma_m$. A publisher-selected PROPOSE,
recovery $\tagg{AUTH}_{r,g}$ path, L1 reference proof sidecar or unbound
native claim is not an input; its optional verification never delays the
rule. The correct signer's full regional checks justify intrinsic input
validity by \cref{prop:final-kernel-soundness}; they do not turn a
private object or a sidecar assertion into a public input. Check typed
identities, actual hashes, the signed domain, the current pair and the
fixed snapshot before verifying FINAL signatures and applying
\tagg{CheckFINAL}. The outer fields of an EXPORT are not checked.

An absent required object or an invalid signature only means that the
candidate does not qualify yet. Other candidates are unaffected, and
since $\mathcal O_q$ only grows, the candidate can qualify in a later
block once the missing object is published. An old object stays in
$\mathcal O_q$ even if nothing was appended when it arrived. It never
grants append permission to a closed version. In contrast, a missing
complete namespace, applicable control prefix or $\tagg{AUTH}_{r',0}$
needed by Steps~1--2 pauses the rule at that step. It is never zero
weight or an empty candidate set. A cache miss for an object already in
$\mathcal O_q$ requires retrieval and resumption at the same
subposition; the prepared retention invariant below rules out such an
Espresso byte miss in the timed cohort. Optional regional evidence is
stored and routed to its own engine; a complete genuine Quit keeps its
evidence-first handling.

Step~4 precedes Step~2 ($\tagg{ResolveBoundary}$), so the terminal DATA
is appended, persisted and exposed before any wait for its successor's
authority. Packet preparation cannot invert this order.

\paragraph{Immediate regional source relay.}
A \tagg{SOURCESTEP} is an authenticated message to an installed fixed
regional snapshot. It carries its chain/protocol domain and sender,
the source height and full header, complete relevant namespace or
absence, finality/linkage evidence and canonical record order. It
identifies the maintained verified predecessor checkpoint, supplies
applicable complete L1 controls and the available fixed-boundary
inputs, and supplies required exact public bytes not already resolvable
by the recipient under the complete-message contract. The
recipient never trusts a sender-selected new checkpoint.

On acquiring a complete raw source input, each correct member stores
and sends it in that availability event, before waiting for its own
\cref{alg:external} result. Add its regional objects to the store
without first fetching optional witness targets. This relay
need not know which region replay will select. Native preparation does
not delay it. Missing mandatory boundary evidence remains an explicit
wait and is supplied in a later complete dependency event. An initial
raw relay is not falsely described as a complete boundary operation
when that evidence is absent. Equal packets are deduplicated by source
position and commitment without suppressing the first correct relay.

A sender includes full required bytes unless recipient resolution is
already covered by the regional full-bundle contract. Its own cache hit
alone does not justify omitting bytes at a recipient. Acquisition of an
additional exact dependency immediately resumes its saved operation;
there is no independent polling tick or replay job.

\paragraph{Receive, commit and activate.}
The receiver authenticates and stores the input, then processes the
contiguous dependency-complete prefix from its maintained checkpoint.
Out-of-order source inputs wait for their predecessors. For each ready
block it executes \cref{alg:external}, that is, the steps of
\cref{alg:global-rule}. Each step records its change to the output and
its subposition, and exposes appended DATA before any later wait. A
duplicate cannot repeat a completed step, count or boundary. Local
receive times do not replace canonical timestamps.

The same complete event delivers the verified transition to the Session
Manager. When the exact anchor continuation and remaining full Ready
inputs are already prepared, the Session Manager checks the fixed handoff and
live activation predicate and performs Ready/initial Enter once. It
attaches the already executable object; it does not import a database
or reconstruct native state here. No old lock, vote or entry time is
inherited. A Force-created version still cannot export in its triggering
block.

On closure, persist the version-wide \tagg{Closed} gate and control
routing obligations atomically. Every entry, fresh action and callback
checks that gate by its own branch version. Worker scans, archive copies and
retirement completion are not prerequisites to gating the old version
or entering a valid new pair. A complete old-session context still
invokes the complete-close handler and multicast under
\cref{sec:retirement}, even for Waiting, Completed or Terminated records.
Unresolved snapshot context remains pending. The per-session
$c_v$, the $4\Delta$ drain and the evidence duties of \cref{sec:retirement} still apply; recording
\tagg{Closed} alone neither starts that drain clock nor permits exit.

\begin{proposition}[Stored public source bytes]
\label{prop:retained-public-bytes}
After fixed-cut hydration, a continuously online history-following
member that executes the retention rules keeps every legally
referenceable acquired Espresso object and its authentication evidence.
If $J_q$ first becomes final at real time $a$, the required Espresso
bytes through $q$ are available to that member by $a+T_E$, without a
new archive read for a previously acquired object.
\end{proposition}
\begin{proof}
Hydration establishes the claim through the fixed cut. At each later
availability event, insertion of the entire raw namespace precedes the
availability-cursor advance and any use by \cref{alg:global-rule}.
Earlier objects of $\mathcal O_q$ are already stored; new objects reside
in this namespace. An absent
object is not fabricated. No retry, checkpoint or worker transition
removes referenceable content. Induction over source positions preserves
the claim. Every predecessor through $q$ became final no later than $a$,
so the stated visibility premise supplies its new source bytes and
Espresso evidence by $a+T_E$. Older hydrated bytes remain present.
This is a statement about Espresso bytes, not a deadline for new L1 authorities.
\end{proof}

\paragraph{Prepared regional processing.}
\label{sec:source-step-timing}
When source bytes, authority and exact executable continuation are
resident, the complete receive handler validates them, performs the
canonical public transition and checks the guarded initial entry.
No historical fetch or native import is silently charged to the
regional communication bound.

The complete-message contract bounds delivery and processing of its
declared inputs. Initial preparation has the separate scope and origin
of \cref{ass:ready-service}. Returning-node catch-up outside that scope
remains eventual. The source-availability cursor, canonical interpreter
cursor and local Ready event can differ; none supplies authority to
execute an unannounced Force.
\section{The Force handler}
\label{sec:force}
\Cref{alg:force} is the only cross-region Force transition. It executes
explicit L1 instructions. \Cref{alg:rotation} consumes the verified
closure and retires local programs separately. Silence, a health report
or a local timeout never runs this transition.

\subsection{What an instruction says}
\label{sec:force-stages}
An MSIG binds four things: the chain, the \emph{target} branch version
$\beta_{\rm target}$, the activation threshold $\theta_a$, and a unique L1
control sequence number $c$:
\[
 (\mathrm{chain},\beta_{\rm target},\theta_a,c).
\]
It names no retained height and implies no status change. Where the new
branch starts is decided when the instruction executes: if the committee
saw $G=108$ when it signed but $G=112$ when the instruction takes effect,
112 is kept. The committee submits the instruction under
\cref{ass:committee,ass:services}. L1 finalization makes it known; its
effect begins at $J_{q(\theta_a)}$. Until then, regular handoffs and
valid old-version output continue. A regular handoff changes $r$, not
$\beta$.

\paragraph{Why the target version is checked.}
Suppose instruction X says ``at time 120, close branch version 7''. If
in the meantime A only handed off to B by a regular handoff, the version
is still 7, so X executes at 120 and closes all of version 7, including
B's regular successor. But if another Force has already moved the
version to 8, X does not execute: its target is 7, and it never turns
into ``close whatever version is current''. So a late or duplicate
instruction cannot cut off a newer branch version. We call such an
instruction \defterm{stale}.

\paragraph{The export session.}
The \defterm{export session} is responsible for the next DATA after $G$.
Its source context, fixed quota and accepted count decide the target
index of a Force (\cref{sec:retry-slot}). It has no cross-region failure
timer. Its program can be Waiting while this public responsibility
exists, and a speculative later program may already be producing.

STATUS changes future eligibility only. A Force without OFF can select
the same physical region when it is the only eligible one. If no region
is eligible, the new boundary keeps its anchor and target index without
opening a regional engine. A later ON can enable selection, and another
valid MSIG can supersede the unassigned boundary.

\begin{protocol}{alg:force}{Apply L1 controls and execute an explicit Force}
Before the payload of authenticated finalized block $J_q$, let
$\theta=\theta_q$ and record the branch version $\beta_0$ at the start
of the block.
\begin{enumerate}[label=\textbf{\arabic*)},itemsep=3pt]
\item \textbf{Apply independent status.}
From the complete applicable L1 control prefix, apply every due
independent STATUS not yet processed, in $(\theta_a,c)$ order. Persist
each effect. Repeated source delivery or equal timestamps never apply a
record twice.

\item \textbf{Check due MSIGs.}
Inspect the unprocessed due MSIGs in the same order. Only the first
valid instruction whose target equals $\beta_0$ executes. Mark every
other due instruction processed without a Force. A version created in
this prelude is never a fresh matching target. If none matches, leave
the branch version open and continue to step~4. No output, Ready,
elapsed-time or health condition is checked.

\item \textbf{Close and keep $G$.}
Record the assigned export session's source context, index and accepted
count, or, if unassigned, the stored source context and target index.
Keep all of $G$ as it stands at the start of $J_q$, with its original
provenance. Record the MSIG cause, the triggering header, the retained
anchor and the selection lower bound $q$. Close $\beta_0$ and set
$\beta=\beta_0+1$, exactly once. Fix the target index by
\cref{sec:retry-slot}. Archive the closed version's unconfirmed suffix
and regular successors. The new engine, when entered, starts with fresh
core state and $g=t=0$. No new-version DATA is accepted in $J_q$ itself.

\item \textbf{Snapshot and resolve.}
Save the eligibility snapshot $\chi_E(q)$ after all control effects.
For a new or pending boundary, find its first eligible historical
NextLive position under the designated authority. Missing evidence
pauses at this canonical operation; it does not prove that no candidate
exists. With complete evidence but no candidate, keep the boundary
unassigned and continue source replay. When selection succeeds, fix the
handoff and open its export session with accepted count zero. No delivery
deadline is created. Then return to ordered payload processing.
\end{enumerate}
\end{protocol}

The cutoff and the retained anchor are fixed even if target evidence
arrives later. Resume the original canonical operation; never rerun a
persisted status change or Force. A control can have a fixed effect while
successor selection still waits for evidence. Step~4 also runs when the
block executes no Force: if an ON makes a waiting boundary selectable,
the boundary is resolved there, without another Force.

Every prelude executes at most one Force. Forces in different blocks
need no gap between them, even when the blocks have equal timestamps.
Each instruction must still satisfy its own authorization,
advance-announcement and target-version conditions.

\subsection{Status, stale controls and the cutoff}
An independent OFF beside a stale MSIG still takes effect. It does not
revoke a handoff that was already fixed at an earlier eligible position.
The stale MSIG neither closes the current version nor selects a new
anchor. The controller must keep these separate meanings in mind when it
schedules operations.

All control effects come before the payload of the triggering block.
DATA accepted in earlier blocks stays in $G$. An old-version result in
the triggering block cannot cancel a Force by arriving with late
progress. It may be kept as regional evidence but cannot append its
suffix. A newer version starts only from the fixed retained anchor.

For example, take threshold $d=120$ and Espresso blocks with timestamps
119, 122 and 125. The block at 119 accepts valid DATA of the current
version. The block at 122 is the first to reach 120, so the Force runs
first and old-version DATA in that block can no longer be accepted;
nor can any later old-version DATA. New-version DATA can be accepted only
from the block after 122.

A proof of closure contains the actual MSIG and the complete applicable
source context. A bare timestamp, an unavailable RPC response, an
unfinalized control or a monitoring alarm is not such a proof. After
learning of the closure, regional Session Managers apply
\cref{sec:retirement}; the cutoff does not wait for their local cleanup.

\section{The external observer}
\label{sec:observers}
To \emph{replay} means to start from a verified state and apply
\cref{alg:global-rule} to the finalized Espresso blocks in order, which
yields the same control results and the same $G$. Any correct verifier that processes, say, an
ON(B), then a due MSIG closing version 7, then an old-version DATA 109,
in that order, reaches the same eligibility, the same closure and the
same acceptance decision. There is no hidden private decision, and the
result does not depend on who sent a handoff message to whom first.

An observer keeps the production branch, the output $\mathrm{Out}$ of
\cref{alg:global-rule} (the global chain with its session records and
tail), the object set $\mathcal O_q$ and its source-verification
cursors. It also persists its replay subposition. It stores no per-block
relay clocks and no committee-protected prefix. The same deterministic
logic can run inside a certifier.

\begin{protocol}{alg:external}{Verify Espresso output and advance the prefix}
\begin{enumerate}[label=\textbf{\arabic*)},itemsep=3pt]
\item \textbf{Authenticate the next source input.} Obtain the next
finalized Espresso block, its complete namespace, and the complete
finalized L1 control context that applies, linked to the trusted prior
prefix. Resume a durable subposition if this block is partly executed.
Resolve the authority $\tagg{AUTH}_{r',0}$ needed by a pending boundary.
Regional objects that are not yet published are not fetched: the rule
reads only $\mathcal O_q$ (\cref{sec:record-completeness}), and optional
proposal, recovery-authority and native sidecars do not delay it. Do not
require the completion of a later successor before appending the
current DATA. Missing genuine source inputs cause a wait, never an
empty-history interpretation.

\item \textbf{Run the prelude once.} Execute \cref{alg:force}, which
contains Step~1 of \cref{alg:global-rule}, before any DATA is appended,
and persist its stage and historical snapshots.
Deliver a verified closure and its source proof to the Session Manager.
Correct regional nodes store and propagate the version-wide gate for
known snapshots, including Waiting and archived programs; unresolved
contexts keep durable relay obligations under \cref{sec:retirement}.

\item \textbf{Apply the global finality rule.} Add the block's regional
objects to $\mathcal O_q$ (\cref{sec:logical-carrier}) and apply
Steps~2--4 of \cref{alg:global-rule} while one of them applies.
Durably expose each appended DATA, its count and the subposition at
once. Invalid or incomplete objects do not block another DATA or
reverse an append. No duplicate DATA is counted again.

\item \textbf{Resolve after a terminal commit.} After Step~4 of
\cref{alg:global-rule} installs a REGULAR boundary, Step~2
($\tagg{ResolveBoundary}$) obtains the target authority and verifies the
historical NextLive selection. A genuine evidence wait persists here
without withholding the terminal DATA that is already accepted. Resume
exactly once at this subposition. With complete evidence but no
candidate, keep the boundary pending and continue source replay.

\item \textbf{Finish or resume.} When no step applies, the block is
done; advance to the next block unless a step is waiting for evidence.
Keep the block's snapshots taken in the prelude unchanged. Persist every
completed step; later source or boundary processing cannot roll it back.
Expose provisional production separately. A Force archives only the
unconfirmed suffix and keeps $G$.
\end{enumerate}
\end{protocol}

Missing evidence and having no candidate are different. If B's authority
has not been downloaded, the node does not know whether B is eligible and
must fetch the evidence. Only when all authorities and eligibility are
known and no region qualifies does the boundary stay pending, waiting
for a later ON. A download failure is never recorded as ``no candidate'',
and late evidence never moves a control's effective position to the
present. A missing regional object is different again: it is not a wait
at all. Step~3 simply finds no candidate, the block finishes, and the
object is used in whichever later block publishes it.

The outer fields of an EXPORT, such as its expected base, are not inputs
to the rule. Step~3 derives the real base from $G$ and checks the signed
parent link. A certifier attestation binds the complete actual previous
and next replay checkpoints: the output $\mathrm{Out}$, the replay
subposition, pending controls and the source position. It can attest an
appended prefix before a later successor selection finishes; it cannot
replace a pending boundary with an invented empty candidate set. The certifier
checks that issued controls were executed correctly. It does not judge,
on the committee's behalf, whether a region was missing blocks for too
long.

\subsection{Execution-layer requests and stored data}
\label{sec:requests}
The execution layer supplies admissible payloads and manages client
queues, ordering policy, persistence and resubmission. The consensus
layer may accept empty DATA under its normal quota rules. It does not
promise that an individual request survives a discarded unconfirmed
suffix or is eventually included. For example, if transaction $x$ is in
an unconfirmed 109 that a Force discards, the user or builder may
resubmit $x$. If $x$ is already in $G$, no Force can remove it.

This does not weaken the storage of block data and proofs. Keep the DATA,
receipts and exact state evidence needed for recovery, handoff and
checkpoint verification. Before Espresso finalization, keep complete
results available to the Publish Layer and direct readers. Afterward,
authenticated history and the node-level store provide the required
objects. Global acceptance never authorizes deleting the only needed
copy. The accepted $G$ binds every publication path.

\paragraph{Faithful publication.}
The proposer or batcher publishes the next complete interval after
$C_{\rm pub}$, even if a Force has changed the current region. The
interval must belong to $G$, keep order and multiplicity, and match the DA
payload and native execution evidence. An unavailable accepted payload
causes waiting, not an invented replacement. The native adapter is
specified separately: a Force cannot move back from a partial-EVM anchor
just because the next signer cannot continue it. It must obtain a legal
continuation or wait.

\section{Examples and guarantees}
\label{sec:examples}
\subsection{Regular rotation and a late DRAIN}
\label{sec:regular-example}
Let $h(A_{\beta,r,0,0})=100$, $\TB_r=10$ and $K=4$. The DATA terms produce
101--104, 105--108 and 109--110. A terminal QC below 110 moves to the next
term. A receipt for DATA 110 opens the next region, which starts DATA 111
after Ready and bootstrap. Step by step:
\begin{enumerate}
\item A commits 110 under its regional BFT rules, if necessary through a
DRAIN during recovery.
\item A's members collect the complete regional receipt for 110 and send
it to B.
\item B verifies that it is selected and obtains the data of 110, the
execution state after 110 and the target authority.
\item B enters $(7,21,0,0)$, runs the ANCHOR bootstrap and produces 111.
\item A's result is published to Espresso in parallel; $G$ may reach 110
later.
\end{enumerate}
Step 2 needs the receipt. A QC for 110 alone is not enough for B to take
over.

\begin{figure}[htbp]\centering
% Protocols 3 and 4: a terminal receipt transfers the turn from A to B.
% Representative arrows only; time runs downward and is not to scale.
\begin{tikzpicture}[
  x=1cm,y=-1cm,font=\fontsize{9}{10.5}\selectfont,
  party/.style={draw,rounded corners=2pt,fill=white,align=center,
    inner sep=3pt,minimum width=2.2cm},
  lane/.style={draw=black!30,densely dashed},
  msg/.style={-{Stealth[length=1.7mm]},semithick,draw=rrBlue},
  direct/.style={msg,densely dashed},
  tag/.style={fill=white,inner sep=1.5pt,font=\fontsize{8}{9}\selectfont},
  note/.style={anchor=north west,align=flush left,text width=5.35cm,
    inner sep=0pt,font=\fontsize{8.5}{10}\selectfont},
  phase/.style={anchor=west,fill=white,inner sep=2pt,
    font=\fontsize{9}{10.5}\selectfont\bfseries},
  dot/.style={circle,fill=black,inner sep=1.1pt},
  timer/.style={line width=2.6pt,draw=rrTeal!75},
  wait/.style={line width=2.6pt,draw=rrBlue!70!black,
    dash pattern=on 2.5pt off 1.5pt},
  entry/.style={draw=rrBlue,line width=1.2pt,dotted}
]
  \def\xA{1.2}\def\xR{3.7}\def\xL{6.2}\def\xB{8.7}\def\xn{9.65}
  \path[use as bounding box] (-.1,-.65) rectangle (15.05,15.75);
  \node[party] at (\xA,0) {A: leader /\\receipt holder};
  \node[party] at (\xR,0) {A: other\\replicas};
  \node[party] at (\xL,0) {B: initial leader\\session $(\beta,r+1)$};
  \node[party] at (\xB,0) {B: other\\replicas};
  \foreach \x in {\xA,\xR,\xL,\xB}
    \draw[lane] (\x,.55)--(\x,15.35);

  \node[phase] at (0,.85) {1.\ A commits terminal DATA in session $(\beta,r)$};
  \draw[msg] (\xA,1.40)--node[tag,above,sloped] {$T_{\beta,r}=110$}(\xR,1.70);
  \draw[msg] (\xR,2.15)--node[tag,above,sloped] {DATA votes}(\xA,2.45);
  \foreach \x/\s/\e in {\xA/1.40/3.80,\xR/1.70/4.10} {
    \draw[timer] (\x+.20,\s)--(\x+.20,\e);
    \node[dot] at (\x,\s) {};
    \draw[thick] (\x-.12,\e)--(\x+.45,\e);
  }
  \node[note] at (\xn,1.35) {Replicas validate and vote for the last DATA
    block. Each starts its own $2\Delta_{\rm A}$ commit timer.};
  \node[tag,anchor=east] at (\xA-.15,2.90) {QC$(110)$};
  \node[note] at (\xn,2.80) {The terminal QC does not hand the turn to B.
    The quota permits no DATA 111 in A's session.};
  \draw[msg] (\xR,4.25)--node[tag,above,sloped] {FINAL$(110)$}(\xA,4.55);
  \node[dot] at (\xA,4.55) {};
  \node[note] at (\xn,3.95) {Timers expire; replicas commit 110 and send
    FINAL. Matching signatures of weight $>F_{\rm A}$ form
    $\FC_{\beta,r}(T_{\beta,r})$. Receipt holders complete the old engine.};

  \node[phase] at (0,5.55) {2.\ Direct transfer of the regular boundary};
  \node[note] at (\xn,5.50) {Verify the session and quota. Select B from
    $\chi_E(q_s)$. Normal handoff keeps $\beta$ and changes $r$ to $r+1$.
    No L1 call is needed.};
  \draw[direct] (\xA,6.35)--node[tag,pos=.55,above,sloped]
    {REGULAR witness + bundle}(\xL,7.35);
  \draw[direct] (\xA,6.35)--(\xB,7.55);
  \node[note] at (\xn,6.90) {Direct transfer can precede Espresso publication. Global
    confirmation of 110 proceeds separately; this handoff does not wait
    for it.};

  \node[phase] at (0,8.05) {3.\ B opens $(\beta,r+1)$ and starts DATA};
  \draw[entry] (\xL-.6,8.45)--(\xL,8.45)--(\xB,8.65)--(\xB+.6,8.65);
  \foreach \x/\y in {\xL/8.45,\xB/8.65} \node[dot] at (\x,\y) {};
  \draw[msg] (\xL,8.75)--node[tag,above,sloped] {Ready relay}(\xB,8.95);
  \draw[msg] (\xB,9.25)--node[tag,above,sloped] {STATUS}(\xL,9.50);
  \draw[wait] (\xL+.20,8.45)--(\xL+.20,10.85);
  \node[note] at (\xn,8.40) {Ready requires the boundary, anchor 110,
    history and authority. Enter view $(\beta,r+1,0,0)$; relay the bundle
    and send STATUS.};
  \node[note] at (\xn,9.95) {The initial leader waits $2\Delta_{\rm B}$
    from its own entry. Progress checks already run.};
  \draw[msg] (\xL,10.90)--node[tag,above,sloped] {NEWVIEW$(110)$}(\xB,11.20);
  \draw[msg] (\xB,11.65)--node[tag,above,sloped] {ANCHOR votes}(\xL,11.95);
  \node[dot] at (\xL,11.95) {};
  \node[note] at (\xn,11.05) {NEWVIEW must choose the inherited anchor.
    Weight $>F_{\rm B}$ gives an ANCHOR QC; these votes start no timer.};
  \draw[msg] (\xL,12.55)--node[tag,above,sloped] {DATA 111}(\xB,12.85);
  \foreach \x/\s/\e in {\xL/12.55/14.95,\xB/12.85/15.25} {
    \draw[timer] (\x+.20,\s)--(\x+.20,\e);
    \node[dot] at (\x,\s) {};
    \draw[thick] (\x-.12,\e)--(\x+.45,\e);
  }
  \draw[msg] (\xB,13.35)--node[tag,above,sloped] {DATA votes}(\xL,13.65);
  \node[note] at (\xn,12.55) {The first fresh block has DATA height 111
    and physical position 1. Its votes start new commit timers. A QC
    permits the next proposal before 111 commits.};
  \node[note] at (\xn,14.60) {If uninterrupted, the new timers commit 111;
    replicas send FINAL for B's session.};
\end{tikzpicture}

\caption{Regular cross-region rotation keeps the DATA, FINAL and ANCHOR
sequence. The historical Espresso selection snapshot fixes the successor.
Regional production and Espresso publication proceed in parallel. Neither
the terminal receipt nor the new region waits for global confirmation.}
\label{fig:regular-cross-region}
\end{figure}
\FloatBarrier

Now suppose 110 has a QC but a Quit cancels its timer. Recovery opens at
height 110 and proposes a DRAIN, which supplies a new timer without using
DATA quota. When the DRAIN commits the terminal DATA, replicas send FINAL
for 110, not for the DRAIN. The regular successor still inherits 110. The
old L1 and Espresso references of 110 say nothing about when it was
committed or exported.

Suppose $G$ is only at 100 while A has committed 110 and B has committed
111. The Publish Layer can export 101--111 with the required A and B
receipts and the regular witness. Once those DATA are accepted from
finalized Espresso history, the whole interval is global. A record that
declares only 111 can still supply the predecessors first if it also
carries their DATA and receipts. It cannot skip ancestors that are truly
missing.

\subsection{A multisignature Force and the retained anchor}
\label{sec:force-example}
This is the running example of \cref{sec:overview}. The region order is
A, B, C. A's session $(7,20)$ produces DATA 101--110. Its terminal receipt
lets B start $(7,21)$ before all of A's DATA is globally accepted. B has
produced through 115, while $G$ ends at 108. The export session is still
A's: the next required DATA is 109.

The committee submits an independent OFF(A) and an MSIG targeting branch
version 7, both with activation threshold 120, and both satisfying
advance announcement. Before the first finalized Espresso block that
reaches 120, valid old output may still extend $G$.

If $G$ still ends at 108 at that block, replay applies OFF(A), checks the
target version and closes version 7 before the payload. It keeps 108. Old
A:109--110 and B:111--115 can no longer be appended. Their receipts stay
as historical evidence of regional commitment.

\begin{center}\small
\begin{verbatim}
Production before the Force:
  A: 101 ... 108 109 110
  B:                     111 ... 115
Kept in G at the Force:
     101 ... 108
New B after the Force:
                new 109, new 110, ...
\end{verbatim}
\end{center}

If B is eligible, the forced target is $(8,21)$, because A has already
contributed accepted DATA. The new engine enters $(8,21,0,0)$ from 108 and
may propose DATA 109 after Ready and the ANCHOR bootstrap. It does not
wait for Espresso results to be registered on L1. New-version DATA is
accepted only in a later Espresso block.

If another valid MSIG had already moved version 7 to 8 before 120, the
instruction targeting 7 would be stale. The independent OFF(A) would still
take effect. If instead A had completed its export and B had become
responsible through regular rotation, the version would still be 7: the
instruction would execute, keep the newer $G$, and select after the
actual export region B.

With quota ten, old B $(7,21)$ starts after 110 and ends at 120. New B
$(8,21)$ starts after 108 and ends at 118. They use the same immutable
authority index but different engines and anchors.

\begin{figure}[p]\centering
% Explicit MSIG control and the existing regional bootstrap.
% beta is the numeric branch version, and (beta,r) identifies a session.
% Vertical spacing illustrates event order, not elapsed time.
\begin{tikzpicture}[
  x=1cm,y=-1cm,font=\fontsize{8.3}{9.8}\selectfont,
  party/.style={draw,rounded corners=2pt,fill=white,align=center,
    inner sep=3pt,text width=1.22cm,minimum height=.66cm},
  lane/.style={draw=black!28,densely dashed},
  msg/.style={-{Stealth[length=1.6mm]},semithick,draw=rrBlue},
  read/.style={msg,densely dashed},
  force/.style={msg,draw=rrRed},
  tag/.style={fill=white,inner sep=1.5pt,font=\fontsize{7.7}{9}\selectfont},
  note/.style={anchor=north west,align=flush left,text width=5.55cm,
    inner sep=0pt,font=\fontsize{8.3}{9.8}\selectfont},
  phase/.style={anchor=west,fill=white,inner sep=2pt,
    font=\fontsize{9}{10.5}\selectfont\bfseries},
  dot/.style={circle,fill=black,inner sep=1.1pt},
  wait/.style={line width=2.5pt,draw=rrTeal!75,
    dash pattern=on 2.5pt off 1.5pt},
  rulebox/.style={draw=rrRed!55,rounded corners=2pt,fill=rrRed!3,
    align=center,text width=7.9cm,inner sep=5pt}
]
  \def\xC{.75}\def\xI{2.3}\def\xE{3.85}
  \def\xN{5.4}\def\xL{6.95}\def\xB{8.5}\def\xn{9.45}
  \path[use as bounding box] (-.05,-.55) rectangle (15.05,16.05);
  \node[party] at (\xC,0) {Multisig};
  \node[party] at (\xI,0) {L1};
  \node[party] at (\xE,0) {Espresso};
  \node[party] at (\xN,0) {Node /\\certifier};
  \node[party] at (\xL,0) {B: initial\\leader};
  \node[party] at (\xB,0) {B: other\\replicas};
  \foreach \x in {\xC,\xI,\xE,\xN,\xL,\xB}
    \draw[lane] (\x,.52)--(\x,15.65);

  \node[phase] at (0,.88) {1.\ Announce MSIG Force on L1};
  \draw[force] (\xC,1.55)--(\xI,1.85);
  \node[tag,anchor=south] at (1.52,1.48) {submit $F$ at $s$};
  \node[note] at (\xn,1.30) {$F$ names target version $\beta$ and future
    time $\theta_a$. It applies only if $\beta$ is still current.
    Independent STATUS keeps its own effect if the Force is stale.};
  \node[dot] at (\xI,2.65) {};
  \draw[read,draw=rrRed] (\xI,2.65)--node[tag,above,sloped]
    {finalized control + complete history}(\xN,3.30);
  \node[dot] at (\xN,3.30) {};
  \node[note] at (\xn,2.85) {By $s+T_{\rm L1}$, all correct participants
    can verify the finalized control. $T_{\rm L1}$ includes dissemination
    and verification. This precedes the Espresso crossing.};
  \node[tag,align=center] at (5.35,4.15)
    {While version $\beta$ remains current: output stays eligible;\\
     normal regional handoffs may continue.};
  \node[phase] at (0,5.25) {2.\ Finalize the Espresso boundary; derive one Force};
  \draw[force] (\xE,5.95)--(\xN,6.25);
  \node[tag,anchor=south] at (4.62,5.90) {$J_{q^*}$ finalized};
  \node[note] at (\xn,5.75) {For MSIG Force,
    $q^*=\min\{q:\theta_q\ge\theta_a\}$.
    Interpret the complete history through $J_{q^*-1}$ first.};
  \node[rulebox] at (4.65,7.30) {At the start of $J_{q^*}$:\\[2pt]
    due independent STATUS $\;\to\;$ first matching MSIG;\\
    no matching instruction means no cross-region Force;\\
    only then process EXPORT records.};
  \node[note] at (\xn,7.05) {Health detection belongs to multisig policy. A report or missing
    DATA cannot execute Force. No Ready, output or minimum-gap check
    limits a valid instruction.};
  \node[dot,fill=rrRed] at (\xN,8.55) {};
  \node[tag,anchor=east,text=rrRed] at (\xN-.12,8.55)
    {anchor $:=G$ before payload};
  \node[note] at (\xn,8.70) {Close version $\beta$ for appends; keep $G$.
    Advance to $\beta+1$. Use NextLive to select eligible B, or wait.
    Old engines close separately.};
  \draw[read] (\xN,9.60)--node[tag,pos=.55,above,sloped]
    {complete boundary}(\xL,10.00);
  \draw[read] (\xN,9.60)--(\xB,10.15);
  \node[note] at (\xn,10.15) {Duty $(\beta_b,r_b)$ has accepted count $a_b$:
    $r'=r_b$ if $a_b=0$, else $r'=r_b+1$.
    Session $(\beta+1,r')$ can append only in a later Espresso block.};

  \node[phase] at (0,11.10) {3.\ B uses the existing ANCHOR bootstrap};
  \node[dot] at (\xL,11.65) {};
  \node[dot] at (\xB,11.85) {};
  \draw[msg] (\xB,12.20)--node[tag,above,sloped] {STATUS}(\xL,12.45);
  \draw[wait] (\xL+.18,11.65)--(\xL+.18,13.15);
  \node[tag,anchor=east] at (\xL-.12,12.85) {wait $2\Delta_{\rm B}$};
  \node[note] at (\xn,11.75) {B verifies the boundary, anchor and state.
    At Ready, enter $(\beta+1,r',0,0)$. The leader waits
    $2\Delta_{\rm B}$, then sends NEWVIEW.};
  \draw[msg] (\xL,13.20)--node[tag,above,sloped] {NEWVIEW}(\xB,13.45);
  \draw[msg] (\xB,14.00)--node[tag,above,sloped] {ANCHOR votes}(\xL,14.25);
  \node[dot] at (\xL,14.25) {};
  \node[note] at (\xn,13.65) {An ANCHOR QC with weight $>F_{\rm B}$
    permits DATA after $G$. ANCHOR votes start no commit timer.
    DATA and DRAIN retain their rules.};
  \draw[msg] (\xL,15.00)--node[tag,above,sloped]
    {DATA at $h(G)+1$}(\xB,15.30);
  \node[note] at (\xn,15.20) {The forced handoff preserves $G$;
    it does not wait for $G$ to be registered on L1.};
\end{tikzpicture}

\caption{Explicit Force: L1 supplies the instruction, Espresso time fixes
its cutoff, and the next region inherits exactly the $G$ at the start of
the activating block. There is no automatic trigger and no second
Espresso-to-L1 wait.}
\label{fig:forced-cross-region}
\end{figure}
\FloatBarrier

\subsection{Concrete view updates}
\label{sec:force-view-examples}
The following examples assume complete authority and an eligible target.
Every Force has an explicit valid MSIG cause.

\paragraph{A has accepted output.}
A is in $(7,20,2,3)$, producing 101--110, while $G=108$. Its export
session is 20 with a positive accepted count. Selecting B gives
\[
 (7,20,2,3)\xrightarrow{\mathrm{Force}}(8,21,0,0).
\]
B continues after 108. Only the new engine's $g$ and $t$ start at zero;
the retained DATA keeps its original provenance.

\paragraph{Regular rotation has already reached B.}
A first hands off
\[
 (7,20,0,2)\xrightarrow{\mathrm{regular}}(7,21,0,0).
\]
B reaches $(7,21,0,1)$ and contributes DATA through 113 to $G$. The next
DATA, 114, belongs to B's session 21. A Force selecting C gives
\[
 (7,21,0,1)\xrightarrow{\mathrm{Force}}(8,22,0,0).
\]
C continues with DATA 114. A scheduled MSIG targeting version 7 still
matches, because regular rotation did not change $\beta$. Even if the
committee saw only $G=108$ when it signed, it cannot make the chain go
back to 108.

\paragraph{Speculative production is ahead of the export session.}
A's $(7,20)$ has regional DATA 101--110; B's $(7,21)$ has 111--120; C runs
$(7,22,1,2)$. If $G=108$, the next required DATA is still A's 109. A Force
selecting B opens $(8,21,0,0)$, not $(8,23,0,0)$. It closes the entire old
version, including C's speculative program. Using a lower turn index than
the latest speculative one removes no globally accepted DATA.

\paragraph{A stale instruction changes no view field.}
One MSIG targets version 7 at threshold 120. A different valid MSIG has
already moved the version to 8 at an earlier position. Suppose the new
engine is now in $(8,21,1,2)$. Processing the stale instruction leaves
\[
 (8,21,1,2)\xrightarrow{\mathrm{stale\ MSIG}}(8,21,1,2).
\]
Any independent STATUS is still evaluated on its own terms. To force
again, the committee needs another valid instruction that targets
version 8.

\subsection{Zero output, an offline region and explicit exclusion}
\label{sec:retry-example}
A's $(7,20)$ opens after DATA 100 but gets no DATA accepted. The committee
observes the failure and submits OFF(A) and MSIG. At their activation, if
$G$ is still 100, the target index stays 20. NextLive can choose B, which
enters $(8,20,0,0)$ with a fresh engine. Reusing the number 20 does not
mean selecting A again.

If B also adds no DATA, another explicit MSIG can choose C in
$(9,20,0,0)$. Neither B's Ready nor a first DATA is required before that
instruction. If A's DATA 101 had been accepted before the first Force,
the new index would be 21 with anchor 101.

The monitoring report itself has no effect.
\begin{center}\small
\begin{tabularx}{\linewidth}{@{}>{\raggedright\arraybackslash}p{.45\linewidth}Y@{}}\toprule
Situation & What the protocol does\\\midrule
A misses blocks for a long time, no MSIG & $\beta$ does not change\\
only OFF(A) & later selections skip A; the current fixed program is not closed\\
only MSIG, A the only eligible region & NextLive may select A again\\
OFF(A) and a valid MSIG, B eligible & hand off to B\\
every region OFF at the Force & keep $G$ and the target index; wait for an explicit ON\\\bottomrule
\end{tabularx}
\end{center}

\subsection{Repeated Force and an unassigned boundary}
\label{sec:repeated-force}
A valid MSIG can close a newly assigned version while its target is
still Waiting, or after Ready but before any DATA. There is no minimum
interval between Forces. Each new instruction still needs a matching
target and a valid advance announcement.

If closing session $(7,20)$, which has accepted output, finds no eligible
region, replay stores target index 21 and anchor $G$ in version 8. It
opens no regional engine. A later valid MSIG can move the version to 9
while keeping target index 21 and the same $G$. When ON(B) eventually
enables selection, B enters $(9,21,0,0)$, not $(9,22,0,0)$: a session
that never started does not count as a completed turn. If no ON is ever
issued, the wait can last forever.

\subsection{Strict cutoff, equal timestamps and concurrent preparation}
\label{sec:cutoff-example}
An MSIG targets version 7 at threshold 300. B may start regular session
$(7,21)$ from A's terminal receipt before that threshold. It may also
prepare a possible Force successor, updating the candidate state as $G$
advances. No candidate grants early entry.

At the first finalized block with timestamp at least 300, the prelude
checks the still-current version. A match closes every session of
version 7 and keeps the exact preceding $G$. Old DATA in that block's
payload cannot revoke the control. Candidate state for an earlier anchor
must be updated to the actual fixed anchor before Ready.

If Espresso timestamps jump across the threshold, the first crossing
block is still the activation block. A later block with the same
timestamp does not repeat an already processed instruction. A distinct
valid instruction can execute at a later position if its own target
matches. A single prelude never chains several Forces.

\subsection{Conditional safety and lifecycle invariants}
\label{sec:safety-argument}
The two propositions below cover the public interpreter and the local
programs. They are safety statements only; progress is discussed in
\cref{sec:liveness}. They use the model and verification interfaces of
\cref{sec:model}. In plain words: everyone who reads the same complete
history accepts the same chain; DATA in $G$ is never removed by a Force;
and no program enters twice or mixes state with another session.

\begin{proposition}[Common interpretation and preservation of global DATA]
\label{prop:global-prefix}
Let two correct observers start from the same authenticated interpreter
checkpoint and apply \cref{alg:global-rule} under the stated history,
announcement and verification premises. Their global DATA sequences are
prefix-comparable: one is a prefix of the other. Replaying to the same canonical subposition with all applicable source
evidence and L1 controls verified yields the same output
$\mathrm{Out}$ of \cref{sec:exports}. At a completed Espresso prefix
this gives the corresponding completed-prefix result as a special case. It does not require identical
future-control queues or source-reading checkpoints. Every accepted DATA
block keeps its content,
position and producing identity forever. Once a version is closed, no
later EXPORT adds new DATA from that version to $G$.
\end{proposition}
\begin{proof}
Induct over the common finalized Espresso sequence and, within a block,
over the steps of \cref{alg:global-rule}.
The initial checkpoints agree. Complete applicable L1 controls,
immutable authorities, Espresso headers and namespace bytes authenticate
the same next input. Advance announcement supplies controls before their
crossing; missing actual evidence pauses at the same operation. STATUS order, the first matching MSIG, target indices and historical
NextLive selection are deterministic. They give the same prelude and
immutable snapshots. Future controls not yet due change neither result.

Both observers read the same object set $\mathcal O_q$ from the same
authenticated namespace content. A private cache changes retrieval work,
not $\mathcal O_q$. The conditions of Step~3 and its hash tie-break are
functions of $\mathcal O_q$ and the current output, so both observers
append the same block. By \cref{prop:final-kernel-soundness},
\tagg{CheckFINAL} establishes intrinsic input validity from a genuine
correct signer's regional commit. The independent boundary, ancestry,
source-order and cutoff checks establish current append permission;
regional safety identifies the same next DATA. Optional admission or
native sidecars cannot veto that kernel. Executable state, native output
and DA correspondence remain subject to their separate sound interfaces. Each DATA transaction appends the
same next block and accounting, and persists the same subposition.
A terminal transaction creates the same pending boundary without waiting
for its successor authority. Later resolution selects only from the prescribed historical inputs. Missing evidence
cannot undo the already committed transaction or alter its timestamp.
These rules give agreement by induction over all committed subpositions.
Two observers that have reached different subpositions hold ordered
states on this append-only DATA sequence, including when one is paused
inside a record. Extra future-control queues and private validation
caches are not part of the compared replay result.

Inspect preservation at each step. STATUS and control ingestion remove no DATA. Force keeps all of $G$ at the start of the block and changes only
the authority for future appends. Selection or deferred resolution of an unassigned boundary
also preserves $G$. Each accepted DATA transaction only extends it. Retained DATA keeps its
original producing pair and hash. Closed versions are never reopened,
and every later DATA candidate checks the current version. Thus
neither an old regular successor nor a late receipt can append a closed
suffix. Local program phases and retirement do not mutate $G$, so they
add no exception to this induction.
\end{proof}

\begin{proposition}[Fixed entry and isolation of local programs]
\label{prop:program-isolation}
Under the regional-validity and source-verification premises, each
opened pair $(\beta,r)$ has one fixed handoff. At a correct replica,
initial Enter occurs at most once for that pair, uses exactly that
handoff's anchor and state, and never uses an obsolete task callback
as authority. Actions for one pair do not alter another pair's locks,
votes, timers or entry time. Candidate preparation creates none of
these BFT actions.
\end{proposition}
\begin{proof}
Proceed from the initial checkpoint through successive handoffs.
For a regular handoff, the predecessor's fixed boundary and quota fix
its terminal height; regional validity fixes the terminal DATA and its
signed selection reference. The least eligible position fixes the
region and designated authority. For Force, common replay fixes the
old and new versions, source context and exact retained $G$; historical
selection then fixes the target. Repeated Forces without assignment
change $\beta$ but keep the stored predecessor index. They cannot
reuse an already opened pair. Handoff verification rejects conflicting
fixed fields instead of replacing an existing record.

A preparation may have several candidate anchors over time. None is
eligible for Enter until verification installs the actual handoff.
Material callbacks keep the requested pair and anchor identity;
the gate checks the exact fixed state and separate current-activation predicate.
The serialized Ready/Enter event tests and sets the durable entry flag
before dispatching signatures. Repetition, restart and late callbacks
therefore cannot repeat Enter or reset its entry time. Each core
message is authenticated and routed by the pair it signs; local tasks
use the same pair and anchor context. Waiting does not run the core.
Closing and Terminated records prevent a late event from reopening an
old program. These checks establish the stated isolation by induction
over the local event sequence.
\end{proof}

This is local authorization relative to a node's processed public
prefix. It does not claim that all nodes learn a finalized Force at the
same real time. A lagging node can briefly run an old program, but
\cref{prop:global-prefix} still excludes its closed-version output.
\Cref{prop:retirement-safety} covers regional behavior during
retirement. Both propositions rely on Part~I and on the DATA/DRAIN and
application-validation interfaces rather than re-proving them.
\subsection{Explicit control and conditional preparation}
\label{sec:liveness}
The guarantees in this subsection do not require every region to receive
service, and they do not make the consensus layer responsible for
request persistence.

\begin{proposition}[No unannounced cross-region Force]
\label{prop:explicit-force}
Starting from an authenticated checkpoint, every later increment of
$\beta$ has a unique effective MSIG cause in the verified L1 control
history. A status update, missing DATA, a local health report or a Ready
event does not increment $\beta$.
\end{proposition}
\begin{proof}
Step~3 of \cref{alg:force} is the only transition that increments
$\beta$. It is reached only for the first valid due instruction matching
the branch version recorded at the start of the block. Control identities
and persisted subpositions prevent a second application. Other
transitions change eligibility, resolve a boundary, append DATA or update
local programs; none changes $\beta$. Induction over committed
transitions proves the claim.
\end{proof}

\begin{proposition}[Preparation under a continuing assignment]
\label{prop:ready-bound}
Let a public assignment open in block $J_o$, with finalization time
$t_0$. If the selected region and its state service satisfy
\cref{ass:ready-service}, the fixed handoff stays applicable, and no
Force supersedes it during preparation, then correct target members can
complete Ready/Enter by $t_0+T_{\rm ready}$. Entry uses the retained
anchor of that assignment and happens at most once per member and pair.
Entry does not require any DATA production.
\end{proposition}
\begin{proof}
A member that already entered satisfies the claim. For a member that has
not entered, the service assumption supplies every material condition of
\cref{def:ready} by the stated time, and continuing applicability
supplies the current-entry check. \Cref{alg:rotation,alg:bootstrap} then
perform the guarded entry in the same complete event.
\Cref{prop:program-isolation} gives anchor binding and at-most-once
entry. The gate does not look at future DATA output. A valid Force before
completion would remove the continuing-assignment premise; it would not
violate any waiting rule.
\end{proof}

\subsection{Conditions for progress}
\label{sec:progress}
Being Ready only means that a region can start. Global progress further
needs regional consensus, a FINAL receipt, publication and global
acceptance. The following conditions are separate from prefix safety.
\begin{itemize}
\item Regional service keeps Part~I's fair leader schedules, eventual
exclusion of exposed equivocators, actual entry into correct turns, and
uninterrupted commit opportunities. These are conditions inside a region,
not fairness among regions.
\item Espresso keeps finalizing retrievable history with growing
timestamps. Complete results reach the Publish Layer, and the declared
publication and finalization services hold.
\item The required authority, exact anchor state and acceptance evidence
are supplied. A selected live region satisfies \cref{ass:ready-service}.
The next required DATA is published in a form and position that pass the
current-branch acceptance checks.
\item If the current region cannot serve, recovery additionally depends
on the multisignature service actually submitting suitable controls,
keeping or enabling an eligible working region, and eventually allowing
enough uninterrupted execution for regional commitment and delivery.
\end{itemize}
The last item is a liveness condition on the control service, not a
condition for accepting an MSIG. The protocol specifies no detection or
response bound and enforces no interval between Forces, and $T_{\rm L1}$
starts only at actual submission. So an honest service can still fail to
give a useful execution opportunity unless its policy provides one.

A wholly offline region is not an operational region for these service
claims. The regional safety and lifetime key-compromise bounds still
govern every accepted certificate. Request persistence, resubmission and
eventual inclusion belong to the execution layer and the client
interface. If Espresso fails, global confirmation stops; switching the
finality service is outside this protocol.

\subsection{What is needed for global progress}
\label{sec:global-progress}
Under the Part~I conditions above, a session that keeps running can form
a complete result. The bounded preparation assumption alone does not
provide that uninterrupted regional opportunity.

If a complete result is available at time $s$, the declared services give
submission by $s+T_{\rm pub}$ and Espresso visibility by
$s+T_{\rm pub}+T_E$. This is transport of the result. To advance $G$, the
result must supply the next required DATA under its fixed boundary, while
its producing branch version is still open at the canonical acceptance
position, and its references and predecessors must pass the checks of
\cref{sec:exports}. Because the rule ignores outer hints, reads the
cumulative $\mathcal O_q$ and appends one DATA at a time, wrong hints or
a later unresolved successor do not erase an otherwise valid next DATA. They never permit a gap or an append from a
closed version.

With these regional, transport, acceptance and control-service
conditions, replay of the suitable result advances $G$. The service
bounds give no wall-clock processing deadline for every arbitrary
observer. No unconditional global-growth theorem follows from repeated
Ready events, published raw records or committee honesty alone.

\paragraph{Offline recovery requires control.}
When the current region cannot continue, progress needs the committee to
submit a suitable MSIG, to disable a failed candidate when necessary, and
to keep or enable a live eligible successor. If no such control is
submitted, nothing else forces a switch. If all regions are OFF, the
boundary can stay unassigned forever. If MSIGs keep closing every useful
interval, no regional output need ever become global. These executions
are allowed unless the governance service provides more liveness
conditions. The committee can thus stall the chain, but under the
safety assumptions it cannot use a Force to roll $G$ back.

\paragraph{Failure and safety are separate.}
An offline region does not meet the operational premises of the
preparation and regional-service claims. A violated lifetime Byzantine
weight bound is a different failure, of the certificate model. Neither
case lets the multisignature committee roll back $G$. No result here
promises fair service among regions, automatic ON, one DATA between
Forces, or eventual inclusion of a client request.

\section{Implementing the Publish Layer}
\label{sec:publish-implementation}
The preceding protocols use a delivery-service assumption. The following
leader-first design explains how existing participants can implement it.
It adds no consensus round and no exclusive publisher role.
The implementation must meet $T_{\rm pub}$; the sketch alone is not a
proof of that bound.

\subsection{Leader first, with permissionless backup}
The proposing regional leader is the preferred sender. Once it has the
complete finalized result, it builds and submits an EXPORT immediately.
Other replicas continue collecting FINAL messages, keeping payloads and
making complete results available to the service. Commitment and
Completion do not wait for a publication acknowledgement.

Another replica or an external data verifier may submit the same result
without the leader's permission. An implementation may give the leader a
short first opportunity, but the backup delay, data retrieval and retries
must fit within $T_{\rm pub}$ from complete-result availability. A private
announcement that the leader has sent something is not enough to suppress
backup indefinitely. The service checks submission and Espresso evidence.

The fallback sender uses the original DATA and receipt. It cannot replace
the DATA, rewrite its session pair or create FINAL on behalf of the region.
The service guarantees delivery through an available path; no particular
sender is guaranteed to win the submission race.

\subsection{Fee sponsorship}
A backup sender need not hold tokens or pay from its own balance. It may
use a funded relay, a prepaid submission account, or another sponsored
route that can submit to Espresso. The route carries the verified result;
it does not acquire authority to change its content.

Someone still supplies the fee resources. To justify the publication
bound, a complete result covered by the service assumption must have access to a funded,
working route even when the original leader is unavailable. A sender can
try another route if one refuses or fails. The concrete sponsorship and
admission mechanism is an implementation choice; this protocol assumes
no Espresso application smart contract for it.

\subsection{Duplicates, retries and acceptance}
Different senders may carry the same complete result in different EXPORTs.
The signed producing pair and block hash identify each DATA; changing a
receipt encoding, signer set or outer hint does not create a new output.
Accepted ancestors are reused as evidence and receive no second count.
A compatible next DATA can be extracted even when the sender's expected
base is stale or its result is listed only among supporting ancestry.

Workers track the original record and its per-DATA acceptance. A complete
next result already recorded with a bad hint is interpreted at that
record position under the same rule; a corrected EXPORT is unnecessary.
If a genuine gap or an absent required object prevents acceptance, the
objects already published stay in $\mathcal O_q$; once the missing
objects appear in a later block, \cref{alg:global-rule} accepts the
result there without republication. No bound on when the missing
objects arrive is inferred. A closed version never regains append
permission.

The two service bounds concern complete-result delivery and visibility.
The conditional progress discussion of \cref{sec:global-progress}
uses those actual recorded bytes. It does not assume a bounded retry
with a corrected expected-base hint. Terminal DATA commits before successor authority
resolution. Preparation for a selected live region has the separate
scope of \cref{ass:ready-service}; arbitrary archive catch-up has no
uniform bound. \Cref{sec:transport-scope} states what the transport
must put into Espresso for the rule to use it.

\subsection{Retention and incentives}
Before finalized Espresso inclusion, the holders of a complete result
keep its payload and evidence available to backup senders. Afterward,
Espresso's retained, authenticated history supplies returning readers and
initial hydration. Prepared regional nodes additionally keep the raw
source bytes under \cref{sec:source-retention}; finalization does not
replace that retained copy by an eventual archive pointer. Direct transfer
can still accelerate regular handoff and state recovery. Retirement and
checkpoint rules preserve required evidence and outgoing duties.

Submission rewards and reimbursement may fund this service. They must not
count duplicate DATA as new progress. A report of delayed delivery
alone does not prove that the original leader, verifier or sponsor
equivocated. Health policy and any objectively verifiable penalties
are separate from the Force acceptance rule.

\Needspace{14\baselineskip}
\begingroup\small
\begin{thebibliography}{9}
\bibitem{part1}
Sync HotStuff RR, Part I: Core Protocol. Design document, 13 October 2026.
Companion source: \href{run:sync_hotstuff_rr_part1_core.tex}{Part I core protocol}.
\bibitem{sync}
I. Abraham, D. Malkhi, K. Nayak, L. Ren, and M. Yin.
\emph{Sync HotStuff: Simple and Practical Synchronous State Machine Replication}.
IEEE Symposium on Security and Privacy, 2020, pp. 106--118.
\href{https://conferences.computer.org/sp/pdfs/sp/2020/349700a106.pdf}{Published paper}.
\bibitem{certification}
MegaETH Specification, PR~8, \emph{Block Certification}.
Commit \texttt{aedc2826663b}.
\href{https://github.com/megaeth-labs/specification/tree/aedc2826663b134ee25d26bea17e70c6d23d321c/specs/block-certification}{Pinned specification}.
Its Espresso and certifier verification interfaces are a reference;
the rotation and explicit-control rules in this document are new.
\end{thebibliography}
\endgroup
\appendix
\section{Proof of \texorpdfstring{\Cref{prop:final-kernel-soundness}}{the FINAL kernel proposition}}
\label{app:final-kernel-proof}
We restate \cref{prop:final-kernel-soundness} in the precise form the
proof uses.

\paragraph{Precise statement.}
\textit{Assume the fixed-snapshot fault bound, the accepted-lifetime
key-compromise bound, the cryptographic assumptions, sound validation
interfaces, and the regional full-validation and FINAL-generation rules.
If $C$ is the authentic incoming context of $(\beta,r)$ and
$\tagg{RawFINAL}_C(D,\Sigma)$ holds, then a correct replica committed
this exact DATA $D$ at a timer event, directly or as an ancestor. Before
that event it fully checked its own-session proposal and required input
ancestry, using established historical kernel conclusions only for
inherited anchors. So a genuine valid historical witness exists for $D$.
Consequently the exact DATA has valid historical regional admission and
native input validity under its authenticated anchor, including the
regional $\rho$ and $\xi$ validity conditions. Exact same-session DATA
ancestors on that hash-linked path inherit these commitment and
input-validity conclusions.}

\begin{proof}
By the fixed-snapshot fault bound, a set of distinct signers with weight
greater than $F_i$ includes a correct signer. The lifetime compromise
condition and signature unforgeability make its signature an actual
correct FINAL action on the stated DATA hash and context. This step uses
the raw threshold, not the definition of a valid regional receipt.

A correct replica sends FINAL only after Commit's timer event commits the
named DATA, either as the timed DATA itself or in the ancestry of the
timed DATA/DRAIN block. Vote started that timer only after all
current-session dependencies were present and the proposal and required
ancestry passed full validation. In particular the replica checked its
session's historical authority and admission, parent justifications,
quota and positions, the L1 and Espresso reference rules, and native
input validity. The sound interfaces apply to these actual regional
checks. An inherited anchor from another producer may rely on that
producer's historical kernel conclusion, justified by the induction
below; this does not replace the replica's exact native state or its
own-session validation. These facts give a genuine valid historical
witness for the committed DATA. It need not be the witness that a later
publisher attaches.

A commit triggered by a receipt, including Completion, emits no new
FINAL. So it cannot supply the first correct signature without the
earlier fully validated timer event identified above. Dependencies on
earlier accepted anchors are justified by induction over execution
events from the trusted initial checkpoint. The FINAL signatures used by
an inherited anchor exist before the new Vote event, so their own timer
origins are earlier cases of the induction. The same argument covers an
unaccepted REGULAR anchor without assuming its global acceptance. No step
assumes the validity of the kernel acceptance being proved. Hash binding
identifies the supplied DATA and exact parent-linked bodies with those in
the correct signer's validated history. The timer's ancestor-commit rule
gives the same-session coverage claim.
\end{proof}

\end{document}
